% vim: foldmethod=marker foldmarker=<<<,>>>
\ifx\danslelivre\undefined
\documentclass[10pt]{../fichiers-configuration-latex/smfart}
\usepackage{ifthen}
\newboolean{english}
\newboolean{francais}
\setboolean{english}{false}
\setboolean{francais}{true}
\input{../fichiers-configuration-latex/paquets}
\selectlanguage{francais}
\input{../fichiers-configuration-latex/lemme-proposition-theoreme}
\input{../fichiers-configuration-latex/numerotation}
\input{../fichiers-configuration-latex/macros}


\externaldocument{../00-introduction/introduction}
\externaldocument{../01-quasi-excellents/excellent}
\externaldocument{../02-topologies/topologies}
\externaldocument{../03-approximation/approximation}
\externaldocument{../04-Cohen-Gabber/Cohen-Gabber}
\externaldocument{../05-algebrisation-partielle/algebrisation}
\externaldocument{../06-log-lissite/log}
\externaldocument{../07-uniformisation-locale/uniformisation-faible}
\externaldocument{../09-uniformisation-premiere-a-l/uniformisation-locale-premiere-a-l}
\externaldocument{../11-produits-orientes/produits-orientes}
\externaldocument{../12-descente-cohomologique-orientee/descente-orientee}
\externaldocument{../13-theoreme-de-finitude/finitude}
\externaldocument{../14-fonctions-dimension/fonctions-dimension}
\externaldocument{../15-Lefschetz-affine/Lefschetz-affine}
\externaldocument{../16-classes-de-Chern/gysin}
\externaldocument{../19-exemples-contre-exemples/contre-exemples}
\externaldocument{../20-rigidite/rigidite}
\externaldocument{../21-finitude-non-abelienne/finitude-non-abelienne}

\author{Joël Riou}
\address{Université Paris-Sud 11, Bât.~425, 91405 Orsay, France}
\email{joel.riou@math.u-psud.fr}
\title{XVII. Dualité}
\date{exposés oraux des 13 avril, 16 mai et 23 mai 2007,
notes du 9 mai 2008}

\setcounter{tocdepth}{3} 

\begin{document}

\begin{center}
XVII. Dualité\\
Joël Riou
\end{center}
%\maketitle

\begin{center}
\verb$Id: dualite.tex,v 1.124 2008-05-09 13:52:48 cvs Exp $
\end{center}

\tableofcontents

\else
\setboolean{english}{false}
\setboolean{francais}{true}
\selectlanguage{francais}
\chapter{Dualité}
\addtocontents{toc}{par Joël Riou}
\begin{center}
Joël Riou
\end{center}
\fi


%Introduction %<<<1

Ce texte vise à fournir une rédaction des résultats annoncés par Ofer
Gabber dans \cite{ABELIEN@Gabber} concernant les complexes dualisants dans
le contexte étale sur les schémas noethériens excellents. 

On fixe un entier naturel $n\geq 2$, on note
$\Lambda=\mathbf{Z}/n\mathbf{Z}$ : ce sera notre anneau de coefficients.
Si $X$ est un $\ZZn$-schéma noethérien, on note $\detalebc X$ la sous-catégorie
de $\detaleb X$ formée des complexes ayant des faisceaux de cohomologie
constructibles et $\detalebctf X$ la sous-catégorie pleine de $\detalebc X$ formée
des complexes de tor-dimension finie.

\begin{definition}\label{17-definition-complexe-dualisant}
Soit $X$ un $\ZZn$-schéma noethérien. Un complexe dualisant sur
$X$ est un objet $K\in\detalebc X$ tel que le foncteur de dualité 
$\Dual_K=\R\SheafHom(-,K)$ préserve $\detalebc X$ et que pour tout
$L\in\detalebc X$, le morphisme de bidualité $L\to\Dual_K\Dual_K L$ soit un
isomorphisme.
\end{definition}

Cette définition diffère de celle de \sga{5}{i}{1.7} dans la mesure où on ne
demande pas à un complexe dualisant d'être de dimension quasi-injective
finie.

Le théorème suivant récapitule l'essentiel des résultats de Gabber que nous
allons établir dans ces notes :

\begin{theoreme}[Gabber]\label{17-theoreme-gabber}
Soit $X$ un $\ZZn$-schéma noethérien excellent muni d'une fonction de
dimension (cf.~définition~\ref{17-definition-fonction-de-dimension}).
Alors, $X$ admet un complexe dualisant $K$, unique à produit
tensoriel avec un objet inversible près (cf.
proposition~\ref{17-proposition-objets-inversibles-d-t-a}).
Ce complexe dualisant $K$ appartient à $\detalebctf X$ ; il est de dimension
quasi-injective finie (autrement dit est un complexe dualisant au sens de
\sga{5}{i}{1.7}) si et seulement si $X$ est de dimension de Krull finie.

Par ailleurs, on a les résultats suivants :
\begin{itemize}
\item si $X$ est régulier, le faisceau constant $\Lambda$ est un complexe 
dualisant ;
\item si $f\colon Y\to X$ est un morphisme plat et à fibres géométriquement
régulières avec $Y$ noethérien excellent, alors $f^!K$ est un complexe
dualisant ;
\item si $f\colon Y\to X$ est un morphisme de type fini compactifiable,
alors $f^!K$ est un complexe dualisant.
\end{itemize}
\end{theoreme}

La démonstration s'appuie sur la notion de complexe dualisant potentiel
(définition~\ref{17-definition-complexe-dualisant-potentiel}) sur un
$\ZZn$-schéma noethérien excellent $X$ muni d'une fonction de dimension
$\delta$ : il s'agit de la donnée d'un complexe $K\in\detaleplus X$ muni
d'isomorphismes (appelés épinglages) $\SheafRG_x(K)\iso
\Lambda(\delta(x))[2\delta(x)]$ pour tout point $x\in X$, qui soient
compatibles aux morphismes de transition associés aux spécialisations
immédiates de points géométriques de $X$ (cf.
section~\ref{17-section-transition-codimension-un}). Gabber montre 
dans \cite[lemma~8.1]{T@Gabber}
que si un $\ZZn$-schéma noethérien excellent admet un
complexe dualisant, alors il admet aussi une fonction de dimension globale ;
cette hypothèse du théorème~\ref{17-theoreme-gabber}
est donc bien nécessaire.

Dans la section~\ref{17-section-premieres-proprietes}, nous verrons notamment
que sur un $\ZZn$-schéma noethérien régulier, le faisceau constant
$\Lambda$ est un complexe dualisant potentiel pour la fonction de dimension
$-\codim$ (cf. proposition~\ref{17-proposition-cdpotentiel-regulier}).
Pour ce faire, nous utiliserons de manière essentielle le
théorème de pureté cohomologique absolue démontré par Gabber, ainsi que les
propriétés des morphismes de Gysin établies
dans~\refext{XVI}{16-section-gysin}.

Dans la section~\ref{17-section-classe-de-cohomologie}, nous construirons un
isomorphisme $\Lambda\iso \HH^{2d}_x(X,\Lambda(d))$ où $X$ est un 
$\ZZn$-schéma noethérien
excellent local strictement hensélien normal de dimension $d$ et de
point fermé $x$, vérifierons que cet isomorphisme est compatible aux
spécialisations immédiates et nous servirons de ce résultat pour
construire des morphismes de transition $\HH^{i}_{\overline y}(X,K)\to
\HH^{i+2c}_{\overline x}(X,K(c))$ pour une spécialisation $\overline
y\to\overline x$ de codimension $c$ arbitraire entre points géométriques d'un
$\ZZn$-schéma noethérien excellent $X$, pour tout
$K\in\detaleplus X$. On utilisera
notamment la résolution des singularités pour les schémas noethériens
excellents de dimension $2$.

Dans la section~\ref{17-section-existence-unicite-cdp}, nous montrerons
l'existence et l'unicité à isomorphisme unique près d'un complexe dualisant
potentiel sur un $\ZZn$-schéma noethérien excellent muni d'une fonction de
dimension. Dans le cas d'un schéma normal, nous nous appuierons sur les
résultats de la section~\ref{17-section-classe-de-cohomologie} et sur les
résultats généraux de Gabber sur l'existence de $t$-structures définies par
des fonctions de perversité. Nous montrerons aussi que les complexes
dualisants potentiels vérifient de bonnes propriétés de stabilité par
rapport aux morphismes plats et à fibres géométriquement régulières et aux
morphismes de type fini.

Dans la section~\ref{17-section-dualite}, nous montrerons qu'un
complexe dualisant potentiel est un complexe dualisant. Une fois les
propriétés de finitude établies, nous procéderons par récurrence sur la
dimension, en utilisant d'une part une généralisation d'un argument de
\sgaquatreetdemi{Th.~finitude}{4.3} et d'autre part le théorème
d'algébrisation partielle~\refext{V}{5-algébrisation partielle mixte raffinée géométrique}.

Dans la section~\ref{17-section-coefficients-generaux}, nous montrons qu'à
partir d'un complexe dualisant à coefficients $\Lambda$, on peut construire
des complexes dualisants pour des anneaux de coefficients plus généraux.
Ces résultats sont essentiellement indépendants des sections précédentes.
Cependant, ce n'est qu'en combinant les résultats des sections précédentes
sur les complexes dualisants potentiels avec le résultat d'unicité des
complexes dualisants de la
proposition~\ref{17-proposition-unicite-complexe-dualisant} que l'on peut
déduire le théorème~\ref{17-theoreme-gabber}. En vertu de ce
théorème~\ref{17-theoreme-gabber}, les constructions de cette section donnent
en particulier des complexes dualisants pour des anneaux de coefficients
généraux sur les schémas noethériens excellents munis de fonctions de
dimension. Il semble très probable qu'il soit possible d'étendre les
résultats d'Ofer Gabber à des énoncés de dualité avec des coefficients
$\ell$-adiques. Toutefois, le rédacteur a renoncé à les rédiger.

Enfin, quelques résultats nécessaires à ce qui précède ont été rejetés en
appendice. On y décrit notamment une construction du produit tensoriel
dérivé sur la catégorie dérivée toute entière des faisceaux de Modules sur
un topos annelé (commutatif).

%>>>1
\section{Le morphisme de transition en \texorpdfstring{codimension~$1$}{codimension 1}}%<<<1
\label{17-section-transition-codimension-un}
\subsection{Notations}%<<<2

\begin{definition2}
Soit $X$ un $\ZZn$-schéma, soit $x\in X$, soit $X\in \detaleplus X$. On pose
$\SheafRG_x(K)=i_x^! K_{|X_{(x)}}\in \detaleplus x$ où $i_x$ est
l'inclusion du point fermé du localisé $X_{(x)}$\;\footnote{On obtiendrait
une définition équivalente en remplaçant le schéma local $X_{(x)}$ par son
hensélisé.}. Si $\overline x$ est un point géométrique de $X$ au-dessus de
$x$\;\footnote{Dans la suite, pour ne pas alourdir inutilement le texte, si
on fixe un point $x$ d'un schéma $X$, $\overline x$ désignera un point
géométrique au-dessus de $x$ et inversement, si on fixe un point
géométrique $\overline x\to X$, on notera $x$ le point de l'espace
topologique sous-jacent à $X$ au-dessus duquel $\overline x$ se trouve.},
on note $\RG_{\overline
x}(K)=\SheafRG_x(K)_{\overline x}\in \dplus(\Lambda)$ ; cet objet
s'identifie à la cohomologie à supports dans le point fermé de la
restriction de $K$ à l'hensélisé strict $X_{(\overline x)}$. Les objets de
cohomologie de $\RG_{\overline x}(K)$ seront notés $\HH^i_{\overline
x}(K)$ pour tout $i\in\mathbf{Z}$.
\end{definition2}

\begin{definition2}[\sga{4}{viii}{7.2}]
Si $\overline y$ et
$\overline x$ sont deux points géométriques d'un schéma $X$, une
spécialisation $\overline y\to \overline x$ est un $X$-morphisme
$X_{(\overline y)}\to X_{(\overline x)}$ entre les localisés stricts
correspondants, ce qui revient à la donnée d'un $X$-morphisme $\overline
y\to X_{(\overline x)}$. On définit la codimension d'une spécialisation
comme étant la dimension de l'adhérence du point de $X_{(\overline x)}$
en-dessous de $\overline y$. On dit qu'une spécialisation est immédiate si
elle est de codimension $1$.
\end{definition2}

On se propose ici de définir un morphisme de transition
$\transition^X_{\overline y\to\overline x}\colon
\RG_{\overline y}(K)\to \RG_{\overline
x}(K)(1)[2]$ dans $\dplus(\Lambda)$ pour toute spécialisation immédiate
$\overline y\to \overline x$ de points géométriques sur un $\ZZn$-schéma
excellent $X$, quel que soit $K\in\detaleplus X$.

%>>>2
\subsection{Cas d'un trait strictement hensélien}%<<<2

Soit $X$ un trait strictement hensélien de point générique $\eta$ et de
point fermé $s$. 
Soit $\overline \eta$ un point géométrique au-dessus de
$\eta$. On va définir le morphisme de transition
\[\RG_{\overline \eta}(K)\to \RG_s(K)(1)[2]\]
pour tout $K\in\detaleplus X$. 

On note $p$ l'exposant caractéristique du corps résiduel de $X$, que l'on
suppose inversible dans $\Lambda$. On dispose
d'une suite exacte canonique de groupes profinis :
\[1\to S\to \Gal(\overline \eta/\eta)\to G\to 1\;\text{,}\]
où $G$ est le groupe d'inertie modérée, canoniquement isomorphe à
$\chap{\ZZ}'(1)$ et où $S$, le groupe de ramification sauvage, est un
pro-$p$-groupe (cf.~\cite[proposition~6.2.12]{ART@Gabber-Ramero}).
Remarquons que l'ordre de $G$ est multiple de
$(\#\Lambda)^\infty$ : ce fait sera utile au lecteur scrupuleux qui
voudrait vérifier en exercice les détails passés sous silence dans cette
sous-section.

\subsubsection{Algèbre du groupe \texorpdfstring{$\chap{\ZZ}'(1)$}{Z'(1)}}%<<<3

\begin{definition3}
L'algèbre de groupe $\Lambda[[G]]$ est l'anneau des endomorphismes du
foncteur d'oubli de la catégorie des $\Lambda$-modules discrets munis
d'une action continue de $G$ vers celle des $\Lambda$-modules. Cette
algèbre est naturellement topologisée : elle est munie de la topologie la
moins fine qui soit telle que pour tout $\Lambda$-module discret $M$ muni
d'une action continue de $G$ et tout élément $m\in M$, l'application
$\Lambda[[G]]\to M$ qui à $a\in\Lambda[[G]]$ associe le résultat $a.m$ de
son action sur $m$ soit continue.
\end{definition3}

On a un isomorphisme canonique d'anneaux topologiques
\[\Lambda[[G]]\iso \lim \Lambda[G/H]\;\text{,}\]
où $H$ parcourt l'ensemble ordonné des sous-groupes ouverts distingués de
$G$ et où $\Lambda[G/H]$ est l'algèbre de groupe usuelle (discrète)
du groupe fini
$G/H$. On peut identifier les $\Lambda$-modules discrets munis d'une action
continue de $G$ aux $\Lambda[[G]]$-modules discrets.

L'action naturelle de $\Lambda[[G]]$ sur $\Lambda$ muni de l'action
triviale de $G$ définit un morphisme continu d'augmentation
$\varepsilon\colon\Lambda[[G]]\to \Lambda$. On note $I_G$ le noyau de
$\varepsilon$ : c'est l'idéal d'augmentation.

\begin{proposition3}\label{17-proposition-structure-i-g}
Le $\Lambda[[G]]$-module $I_G$ est libre de rang $1$, engendré par
$1-\sigma$ si $\sigma$ est un générateur topologique de $G$.
\end{proposition3}

La démonstration de cette proposition est laissée en exercice au lecteur.

%>>>3
\subsubsection{Description de la cohomologie à supports}%<<<3

Il est clair que la catégorie des faisceaux (d'ensembles) sur $X_\et$ est
naturellement équivalente à la catégorie des flèches de $\Gal(\overline
\eta/\eta)$-ensembles (discrets) ${\mathscr M}_s\to 
{\mathscr M}_{\overline \eta}$ telles que 
l'action de $\Gal(\overline \eta/\eta)$ sur ${\mathscr M}_s$ soit triviale.

Pour tout faisceau de $\Lambda$-modules $\mathscr M$ sur $X_\et$, on note
$F(\mathscr M)$ le complexe évident de $\Lambda$-modules :
\[{\mathscr M}_s\to {\mathscr M}_{\overline \eta}^S\to
\Hom_{\Lambda[[G]]}(I_G,{\mathscr M}_{\overline \eta}^S)\]
où ${\mathscr M}_s$ est placé en degré $0$. Le foncteur qui à $\mathscr M$
associe $F(\mathscr M)$ étant additif, en utilisant la construction du
complexe simple associé à un complexe double, on peut naturellement définir
un complexe de $\Lambda$-modules $F(K)$  pour tout complexe $K$ de
faisceaux de $\Lambda$-modules sur $X_\et$. En observant en outre que le
foncteur qui à $\mathscr M$ associe $F(\mathscr M)$ est exact, on voit que 
le foncteur $F$ défini au niveau des complexes préserve les
quasi-isomorphismes. On a ainsi obtenu un foncteur $F\colon \detaleplus X\to
\dplus(\Lambda)$.

\begin{proposition3}\label{17-proposition-description-h-2-support-trait}
Pour tout $K\in\detaleplus X$, le morphisme évident
\[\RG_s(X,K)\to F(K)\]
est un isomorphisme dans $\dplus(\Lambda)$.
\end{proposition3}

On se ramène au cas où $K$ se réduit à un faisceau injectif $\mathscr M$
placé en degré $0$. Il est évident que le noyau de 
${\mathscr M}_{\overline \eta}^S\to 
\Hom_{\Lambda[[G]]}(I_G,{\mathscr M}_{\overline \eta}^S)$ s'identifie à
$\Gamma(\eta,\mathscr M)$.
On a donc bien un isomorphisme $\HH^0_s(X,\mathscr M)\iso
\HH^0(F(\mathscr M))$. Il reste à montrer que $\HH^1(F(\mathscr M))$ est nul.
Cela revient à dire qu'après application du foncteur
$\Hom_{\Lambda[[G]]}(-,{\mathscr M}_{\overline \eta}^S)$, l'inclusion de
l'idéal $I_G$ dans $\Lambda[[G]]$ induit une surjection.
Pour cela, il suffit de montrer que le $\Lambda[[G]]$-module discret
${\mathscr M}_{\overline \eta}^S$ est injectif :
par restriction à un ouvert, ${\mathscr M}_{|\eta}$
est un faisceau injectif de $\Lambda$-modules sur $\eta_\et$ et ${\mathscr
M}_{\overline \eta}^S$ est obtenu à partir de ${\mathscr M}_{|\eta}$ par
application du foncteur image directe associé au morphisme
de topos évident $\eta_\et \to \BB G$ (cf. \sga{4}{iv}{4.5.2}).

%>>>3
\subsubsection{Définition du morphisme de transition}%<<<3

Soit $\sigma$ un générateur topologique de $G$. Soit $M$ un
$\Lambda[[G]]$-module discret. On observe que l'on a un isomorphisme
canonique de groupes abéliens $M(-1)\simeq \Hom(G,M)$. On définit un
morphisme de groupes abéliens
$M(-1)\to \Hom_{\Lambda[[G]]}(I_G,M)$ \emph{via} les
isomorphismes suivants :
\[\xymatrix{M(-1)\ar[r]^-{\sim} &
\Hom(G,M)\ar[r]_-\sim^-{\ev_\sigma} & M &
\ar[l]^-\sim_-{\ev_{1-\sigma}} \Hom_{\Lambda[[G]]}(I_G,M)\;\text{.}}\]

Si $\mathscr M$ est un faisceau de $\Lambda$-modules sur $X_\et$, en
composant l'application ci-dessus dans le cas de $M={\mathscr M}_{\overline
\eta}^S$ et le projecteur canonique 
${\mathscr M}_{\overline
\eta}\to {\mathscr M}_{\overline
\eta}^S$ ($S$ étant un pro-$p$-groupe et $p$ inversible dans $\Lambda$), on
définit un morphisme ${\mathscr M}_{\overline \eta}(-1)\to
\Hom_{\Lambda[[G]]}(I_G,{\mathscr M}_{\overline \eta}^S)$ et donc un
morphisme de complexes $s_\sigma\colon {\mathscr M}_{\overline \eta}(-1)\to
F(\mathscr M)[2]$, fonctoriel en $\mathscr M$.

\begin{definition3}\label{17-definition-transition-cas-trait}
On note $\transition^X_{\overline\eta\to s}\colon \RG_{\overline \eta}(K)\to
\RG_s(K)(1)[2]$ le morphisme dans $\dplus(\Lambda)$ défini par 
$s_\sigma$ fonctoriellement pour tout objet $K\in\detaleplus X$. D'après le
lemme suivant, ce morphisme $\transition^X_{\overline\eta\to s}$ 
est indépendant du générateur $\sigma$ de $G$ : c'est le morphisme de
transition associé à la spécialisation $\overline\eta\to s$ de points
géométriques de $X$.
\end{definition3}


\begin{lemme3}
Si $\sigma$ et $\sigma'$ sont deux générateurs topologiques de $G$, il
existe une unique homotopie fonctorielle en $\mathscr M$ entre $s_\sigma$
et $s_{\sigma'}$.
\end{lemme3}

On voit aussitôt qu'on peut se limiter aux $\mathscr M$ tels que 
${\mathscr M}_s=0$ et que $S$ agisse trivialement sur ${\mathscr
M}_{\overline \eta}$. On peut identifier cette catégorie de faisceaux à
celle des $\Lambda[[G]]$-modules discrets.

Soit $M$ un $\Lambda[[G]]$-module. Notons $\mathscr M$ le faisceau sur
$X_\et$ correspondant. On note $F_\sigma(M)$ le complexe
\[\xymatrix{\dots\ar[r] & 0\ar[r] & M\ar[r]^{1-\sigma} & M\ar[r] & 0 \ar[r] & \dots}\]
concentré en les degrés $1$ et $2$. On note $\Psi_\sigma\colon F(\mathscr
M)\iso F_\sigma(M)$ l'isomorphisme de complexes défini de façon évidente
à partir de $\sigma$ :
\[
\xymatrix{
\dots \ar[r] & \ar[d]0 \ar[r] &\ar@{=}[d] {\mathscr M}_{\overline \eta}\ar[r] &
\Hom_{\Lambda[[G]]}(I_G,{\mathscr M}_{\overline \eta})
\ar[d]^{\ev_{1-\sigma}}_\sim  \ar[r] & 0 \ar[r]\ar[d] & \dots\\ 
\dots \ar[r] & 0 \ar[r] & M\ar[r]^{1-\sigma} & M \ar[r] & 0 \ar[r] &\dots}
\]

Notons $\varphi_\sigma\colon {\mathscr M}_{\overline \eta}(-1)\iso M$
l'isomorphisme défini par l'évaluation en $\sigma$ \emph{via}
l'isomorphisme canonique ${\mathscr M}_{\overline\eta}(-1)\simeq
\Hom(G,{\mathscr M}_{\overline \eta})$. Notons $t_\sigma\colon M\to
F_\sigma(M)[2]$ le morphisme de complexe représenté par les flèches
verticales ci-dessous :
\[
\xymatrix{\dots \ar[r] & \ar[d]0 \ar[r] & \ar[d]0 \ar[r] &\ar@{=}[d] M\ar[r]
&\ar[d] 0 \ar[r] & \dots\\
\dots \ar[r] & 0 \ar[r] & M\ar[r]^{1-\sigma} & M\ar[r] & 0 \ar[r] & \dots}
\]

On dispose ainsi d'un carré commutatif de complexes, fonctoriel en $M$ :
\[
\xymatrix{{\mathscr M}_{\overline\eta}(-1)\ar[r]^-{s_\sigma}
\ar[d]^{\varphi_\sigma}_\sim& F(\mathscr
M)[2] \ar[d]^{\Psi_\sigma}_\sim\\
M\ar[r]^-{t_\sigma} & F_\sigma(M)[2]}
\]
Posons $f_{\sigma,\sigma'}=\Psi_{\sigma}\circ
s_{\sigma'}\circ\varphi_{\sigma}^{-1}$.
Les flèches verticales étant des isomorphismes de complexes, montrer que
les morphismes $s_\sigma,s_{\sigma'}\colon {\mathscr
M}_{\overline\eta}(-1)\to F(\mathscr M)[2]$ sont (fonctoriellement)
homotopes revient à vérifier que les deux morphismes
$t_\sigma,f_{\sigma,\sigma'}\colon M \to F_\sigma(M)[2]$ le sont.

On peut ainsi représenter la situation de façon plus concrète :
\[
\xymatrix{
\ar[d] 0\ar[r] & M\ar@/_/[d]_{\Id}\ar@/^/[d]^{g}\\
\ar[r]^{1-\sigma} M & M
}\]
où $g$ est la transformation naturelle induite par $f_{\sigma,\sigma'}$.

Comme $\Lambda[[G]]$ désigne précisément l'anneau des transformations
naturelles $M\to M$, on peut identifier $\Id$ et
$g$ à des éléments $1$ et $g$ de $\Lambda[[G]]$
respectivement.
Montrer l'existence et l'unicité de l'homotopie fonctorielle entre
$s_\sigma$ et $s_{\sigma'}$ se ramène donc à montrer
l'existence et l'unicité de $h\in\Lambda[[G]]$ tel que $(1-\sigma)\cdot
h=1-g$. D'après la proposition~\ref{17-proposition-structure-i-g}, cela
revient à montrer que $\varepsilon(g)=1$.
Si on note $u$ l'unité de $\Lambda[[G]]$ telle que 
$(1-\sigma')=u\cdot (1-\sigma)$ et qu'on note $\alpha$ l'élément de 
${\chap{\ZZ}{'}}^\times$
tel que $\sigma'=\sigma^\alpha$, alors on a la relation
$u\cdot g=\alpha$.
On est donc ramené à montrer que
$\varepsilon(u)=\varepsilon(\alpha)$. Pour cela, on utilise la formule
suivante :
\[
\frac{1-\sigma^\beta}{1-\sigma}=\sum_{i=0}^{\beta-1}\sigma^i\;\text{.}
\]
Cette formule est évidemment juste pour $\beta\in\mathbf{N}$ ; on peut lui
donner un sens pour tout $\beta\in\chap{\ZZ}$ en prolongeant chacun des
membres par continuité. En appliquant cette formule avec
$\beta=\alpha$, on obtient le
résultat voulu :
\[\varepsilon(u)=\varepsilon\left(\sum_{i=0}^{\alpha-1}\sigma^i\right)
=\sum_{i=0}^{\alpha-1}1=\varepsilon(\alpha)\;\text{,}\]
ce qui achève la démonstration du lemme.

%>>>3
%>>>2
\subsection{Cas d'un schéma local strictement hensélien intègre excellent de dimension~1}%<<<2

Soit $X$ un $\ZZn$-schéma local strictement hensélien intègre excellent de
dimension $1$. Soit $\eta$ le point générique de $X$. Soit $\overline \eta$
un point géométrique au-dessus de $\eta$. Soit $s$ le point fermé de $X$.
Soit $\tilde X\vers f X$ la normalisation de $X$. Le schéma $\tilde X$ est
un trait strictement hensélien (de point fermé $\tilde s$)
et $f$ est un homéomorphisme
universel (en particulier, $\tilde s/s$ est une extension purement inséparable),
donc le foncteur image inverse $f^\star$ induit une équivalence
entre la catégorie des faisceaux sur $X_\et$ et sur ${\tilde X}_\et$.

\begin{definition2}\label{17-definition-transition-cas-local-dimension-un}
\emph{Via} les identifications ci-dessus, 
le morphisme de transition
$\transition^X_{\overline\eta \to s}$ est induit par $\frac 1 {[\tilde s:s]}
\transition^{\tilde X}_{\overline
\eta\to\tilde{s}}$ (cf. définition~\ref{17-definition-transition-cas-trait}).
\end{definition2}

%>>>2
\subsection{Cas général}%<<<2

\begin{definition2}\label{17-defi-transition-codimension-un}
Soit $X$ un $\ZZn$-schéma excellent. Soit $\overline y\to\overline x$
une spécialisation immédiate de points géométriques de $X$. Pour définir le
morphisme de transition $\transition^X_{\overline y\to\overline x}\colon
\RG_{\overline y}(K)\to\RG_{\overline x}(K)(1)[2]$ pour tout
$K\in\detaleplus X$, quitte à remplacer $K$ par son image inverse \emph{via} le
morphisme canonique $X_{(\overline x)}\to X$, on peut supposer que $X$ est
local strictement hensélien (excellent) de point fermé $\overline x$. On
note alors $Z$ l'adhérence du point de $X$ en-dessous de $\overline y$ et
$i\colon Z\to X$ son immersion dans $X$. Ce $\ZZn$-schéma $Z$ est local
strictement hensélien intègre excellent de dimension $1$, le morphisme de
transition $\transition^{Z}_{\overline y\to\overline z}$ a été introduit dans la
définition~\ref{17-definition-transition-cas-local-dimension-un}.
Pour tout $K\in\detaleplus X$, on
définit le morphisme de transition $\transition^X_{\overline y\to\overline x}$ de
façon à faire commuter le diagramme suivant où les flèches verticales sont
les isomorphismes évidents :
\[
\xymatrix{\ar[d]^\sim \RG_{\overline y}(K)\ar[r]^-{\transition^X_{\overline
y\to\overline x}} & \RG_{\overline
x}(K)(1)[2]\ar[d]^\sim \\
\RG_{\overline y}(i^!K)\ar[r]^-{\transition^Z_{\overline
y\to\overline x}} & \RG_{\overline x}(i^!K)(1)[2]}
\]
\end{definition2}

\begin{remarque2}\label{17-remarque-composition-transition-codimension-0-ou-1}
Pour toute spécialisation $\overline x'\to\overline x$ de codimension $0$
entre points géométriques de $X$ (essentiellement, un élément du groupoïde
de Galois absolu du corps résiduel d'un des points de $X$), on a un
isomorphisme évident $\RG_{\overline x'}(K)\iso \RG_{\overline
x}(K)$, que l'on note $\transition^X_{\overline x'\to \overline x}$. Il est
évident que si $\overline z\to\overline y$ et $\overline y\to \overline x$
sont des spécialisations composables telles que la codimension $c$ de
$\overline z\to\overline x$ soit $0$ ou $1$, on a une égalité de morphismes
\[\transition^X_{\overline z\to\overline x}=\transition^X_{\overline y\to\overline
x}\circ \transition^X_{\overline z\to\overline y}\colon \RG_{\overline
z}(K)\to \RG_{\overline x}(K)(c)[2c]\;\text{.}\]
Ainsi, les morphismes de transition associés aux
spécialisations se composent bien dans l'étendue où ces constructions ont
été faites jusqu'à présent. Nous définirons plus tard des morphismes de
transition en codimension arbitraire et ce de façon compatible à la
composition, mais seulement au niveau des groupes de cohomologie
(cf. théorème~\ref{17-theoreme-transition-toute-codimension}).
\end{remarque2}

%>>>2
%>>>1
\section{Complexes dualisants putatifs et potentiels}%<<<1
\label{17-section-premieres-proprietes}
\subsection{Définition des complexes dualisants putatifs et potentiels}%<<<2

\begin{definition2}
\label{17-definition-fonction-de-dimension}
Une fonction de dimension sur un schéma localement
noethérien $X$ est une fonction $\delta\colon X\to\mathbf{Z}$ telle que
pour toute spécialisation immédiate $\overline y\to\overline x$ de points
géométriques de $X$, on ait $\delta(y)=\delta(x)+1$.
\end{definition2}

Localement pour la topologie de Zariski, un schéma excellent admet une
fonction de dimension, cf~\refext{XIV}{14-theoreme-univ-cat-fonction-dim}.

\begin{definition2}\label{17-definition-complexe-dualisant-potentiel}
Soit $X$ un $\ZZn$-schéma noethérien
excellent muni d'une fonction de dimension $\delta$. Un
complexe dualisant putatif consiste en la donnée de $K\in\detaleplus X$ et pour
tout $x\in X$ d'un isomorphisme (appelé épinglage en $x$)
$\SheafRG_x(K)\iso \Lambda(\delta(x))[2\delta(x)]$ dans $\detaleplus
x$. Un complexe dualisant putatif est un complexe dualisant potentiel si
pour toute spécialisation immédiate $\overline y\to\overline x$, le
diagramme suivant est commutatif :
\[\xymatrix{
\ar[rd]^\sim\RG_{\overline y}(K)
\ar[r]^-{\transition^X_{\overline y\to\overline x}}& 
\RG_{\overline x}(K)(1)[2]\ar[d]^\sim \\
& \Lambda(\delta(y))[2\delta(y)] 
}\]
Autrement dit, les épinglages sont compatibles aux morphismes de transition
associés aux spécialisations immédiates.
\end{definition2}

Les notions de complexes dualisants putatifs et potentiels (à coefficients
$\Lambda$) ne sont définis que pour les $\ZZn$-schémas noethériens
excellents munis d'une fonction de dimension. Certains des énoncés à venir
contiendront donc implicitement ces hypothèses sur les schémas. La fonction
de dimension ne sera pas non plus systématiquement mentionnée dans les
énoncés.

\begin{remarque2}
Si $X$ est connexe et que $K\in\detaleplus X$ est muni de deux 
structures de complexe dualisant
potentiel, pour vérifier que les épinglages sont les mêmes en tous les
points de $X$, il suffit de le faire en un seul point.
\end{remarque2}

L'objectif de cette section est de montrer que sur un schéma régulier
excellent muni de la fonction de dimension $-\codim$, le faisceau constant
$\Lambda$ est naturellement muni d'une structure de complexe dualisant
potentiel.

%>>>2
\subsection{Fonctorialité par rapport aux morphismes étales}%<<<2

Des propriétés de stabilité importantes des complexes dualisants potentiels
par rapport à certaines classes de morphismes seront obtenues dans la
section~\ref{17-section-complements-cdp}. Pour le moment, mentionnons
simplement la compatibilité suivante pour les morphismes étales :

\begin{proposition2}\label{17-proposition-stabilite-cdp-morphisme-etale}
Soit $f\colon Y\to X$ un morphisme étale entre $\ZZn$-schémas noethériens
excellents. On suppose que $X$ est muni d'une fonction de dimension
$\delta_X$. On définit une fonction de dimension $\delta_Y$ sur $Y$ en
posant $\delta_Y(y)=\delta_X(f(y))$ pour tout $y\in Y$. Soit $K$ un
complexe dualisant putatif sur $X$. Alors, $f^\star K$ est naturellement
muni d'une structure de complexe dualisant putatif et c'est un complexe
dualisant potentiel si $K$ en est un.
\end{proposition2}

Si $y$ est un point de $Y$, que $x=f(y)$ et que l'on note $g\colon y\to x$
le morphisme induit par $f$, on a un isomorphisme canonique $g^\star
\SheafRG_x(K)\simeq \SheafRG_y(f^\star K)$. Ceci permet de
définir les épinglages sur $f^\star K$. On vérifie aussitôt que si $K$ est
un complexe dualisant potentiel, alors $f^\star K$ aussi.

\bigskip

La construction de cette proposition passe évidemment à la limite : le
résultat vaut aussi pour des localisations $X_{(x)}\to X$,
$X_{(x)}^\hensel\to X$ ou $X_{(\overline x)}\to X$. Nous utiliserons
librement ces observations simples dans la suite.

%>>>2
\subsection{Compléments sur les spécialisations}%<<<2

\begin{definition2}\label{17-definition-specialisation-au-dessus-d-une-autre}
Soit $X'\to X$ un morphisme de schémas. Soit $\overline{y'}\to \overline{x'}$
et $\overline y\to \overline x$ des spécialisations de points géométriques 
de $X'$ et $X$ respectivement. Si l'on se donne des $X$-morphismes 
$\overline {y'}\to\overline y$ et $\overline{x'}\to \overline x$ tels que
le diagramme évident ci-dessous commute, alors on dit que 
$\overline{y'}\to\overline{x'}$
est au-dessus de $\overline y\to\overline x$.
\[\xymatrix{\overline {y'}\ar[d]\ar[r] & X'_{(\overline x')} \ar[d]\\
\overline y\ar[r] & X_{(\overline x)}}
\]
\end{definition2}


\begin{proposition2}\label{17-proposition-specialisation-going-down}
Soit $X'\to X$ un morphisme de schémas. Soit
$\overline {y'}\to\overline {x'}$ une
spécialisation de points géométriques de $X'$. Alors, à des isomorphismes
uniques près, il existe une unique spécialisation $\overline y\to\overline
x$ de points géométriques de $X$ en-dessous de $\overline {y'}\to\overline
{x'}$.
\end{proposition2}

C'est évident.

\begin{proposition2}\label{17-proposition-specialisation-going-up}
Soit $X'\to X$ un morphisme fini et $\overline y\to\overline x$ une
spécialisation de points géométriques de $X$. Soit $\overline
{y'}\to X'$ un point géométrique de $X'$ au-dessus de $\overline y$
(\emph{i.e.}
on s'est donné un $X$-morphisme $\overline {y'}\to\overline y$). Alors, à
des isomorphismes uniques près, il existe une unique spécialisation
de points géométriques de $X'$ au-dessus de $\overline y\to\overline x$ de
la forme $\overline {y'}\to\overline {x'}$.
\end{proposition2}

On peut supposer que $X$ est local strictement hensélien de point fermé
$\overline x$. Le schéma $X'$ étant fini sur $X$, c'est une réunion disjointe
finie de schémas locaux strictement henséliens. Quitte à remplacer $X'$ par
la composante connexe contenant $\overline {y'}$, on peut supposer que $X'$
est lui-aussi local strictement hensélien. Il n'y a alors manifestement
plus d'alternative : $\overline {x'}$ est le point fermé de $X'$.

%>>>2
\subsection{Construction d'un complexe dualisant potentiel dans le cas régulier}%<<<2
\subsubsection{Un complexe dualisant putatif}%<<<3
\begin{proposition3}\label{17-proposition-cdputatif-regulier}
Soit $X$ un $\ZZn$-schéma noethérien régulier excellent muni de la fonction de
dimension $-\codim$. Alors $\Lambda$ est naturellement muni d'une structure
de complexe dualisant putatif.
\end{proposition3}

Soit $x\in X$. Pour définir l'épinglage en $x$, on peut supposer que $X$
est local de point fermé $x$. On note $i\colon x\to X$ l'inclusion de ce
point fermé. Le morphisme de Gysin $\Cl i\colon
\Lambda(\delta(x))[2\delta(x)]\to i^!\Lambda$ est un isomorphisme d'après
le théorème de pureté cohomologique absolue
(cf.~\refext{XVI}{16-theoreme-purete-absolue}). L'épinglage en $x$ est
l'isomorphisme inverse.

%>>>3
\subsubsection{Cas d'un trait}%<<<3

\begin{proposition3}\label{17-proposition-cdp-trait}
Soit $X$ un trait excellent, muni de la fonction de dimension
$\delta=-\codim$. On suppose que $X$ est un $\ZZn$-schéma.
\begin{enumerate}
\item[(a)] Le complexe dualisant putatif $\Lambda$ de la
proposition~\ref{17-proposition-cdputatif-regulier} est un complexe dualisant
potentiel.
\item[(b)] Si $K$ est un complexe dualisant putatif, il existe un unique
morphisme $\Lambda\to K$ compatible aux épinglages au point générique.
\item[(c)] Si $K$ est un complexe dualisant potentiel, le morphisme
$\Lambda\to K$ défini ci-dessus est un isomorphisme (compatible aux
épinglages).
\end{enumerate}
\end{proposition3}

La compatiblité (a) est facile (au moins au signe près). Établissons (b).
Soit $K$ un complexe dualisant putatif. Notons $i\colon s\to X$ l'inclusion
du point fermé $s$ de $X$ et $j\colon\eta\to X$ l'inclusion de son point
générique $\eta$. On a un triangle distingué canonique, que l'on peut
récrire en présence d'épinglages :
\[\xymatrix{\ar[d]^{\sim}i_\star i^! K\ar[r] & \ar@{=}[d] K \ar[r] &
\ar[d]^\sim \R j_\star j^\star
K\ar[r] & i_\star i^! K[1]\ar[d]^\sim \\
i_\star \Lambda(-1)[-2]\ar[r] & K \ar[r] & \R j_\star \Lambda\ar[r] & 
i_\star \Lambda(-1)[-1]}\]
En appliquant le foncteur $i$-ème faisceau de cohomologie ${\mathscr H}^i$,
on obtient l'annulation de ${\mathscr H}^i K$ pour $i<0$ et un 
isomorphisme canonique ${\mathscr H}^0K\simeq \Lambda$. En vertu des propriétés
élémentaires de la $t$-structure canonique sur $\detaleplus X$, on obtient un
unique morphisme $\Lambda\simeq {\mathscr H}^0 K\to K$ compatible à
l'épinglage au point générique.

Pour obtenir (c), supposons que $K$ soit un complexe dualisant potentiel. On
considère le diagramme commutatif induit par le morphisme canonique
$\Lambda\to K$ de (b) et la fonctorialité du morphisme de spécialisation
associé à un choix de  spécialisation $\overline \eta\to \overline s$ 
au-dessus des points $\eta$ et $s$ :
\[
\xymatrix{
\RG_{\overline\eta}(\Lambda)\ar[r]\ar[d]^{\transition^X_{\overline \eta\to\overline s}} & \RG_{\overline
\eta}(K)\ar[d]^{\transition^X_{\overline \eta\to\overline s}} \\
\RG_{\overline s}(\Lambda)(1)[2]\ar[r] & \RG_{\overline s}(K)(1)[2]}
\]
Comme $\Lambda$ et $K$ sont des complexes dualisants potentiels, les
flèches verticales sont des isomorphismes. Par construction, le morphisme
du haut est un isomorphisme. Il en résulte que le morphisme du bas aussi.
Par conséquent le morphisme $\Lambda\to K$ induit un isomorphisme après
application de $j^\star$ et de $i^!$ : c'est un isomorphisme.


%>>>3
\subsubsection{Fonctorialité par rapport aux morphismes quasi-finis}%<<<3

\begin{proposition3}\label{17-proposition-cdputatif-quasi-fini}
Soit $f\colon Y\to X$ un morphisme quasi-fini. Soit $K$ un complexe dualisant
putatif sur $X$ pour une certaine fonction de dimension $\delta_X$ sur $X$.
Alors $f^!K$ est naturellement muni d'une structure de complexe dualisant
putatif pour la fonction de dimension $\delta_Y$ sur $Y$ définie par
$\delta_Y(y)=\delta_X(f(y))$ pour tout $y\in Y$.
\end{proposition3}

Soit $y\in Y$. Notons $x=f(y)$ et $\pi\colon y\to x$ le morphisme fini
induit par $f$. On a un isomorphisme canonique dans $\detaleplus y$ :
\[\SheafRG_y(f^!K)\simeq \pi^! \SheafRG_x(K)\;\text{.}\]
L'épinglage en $x$ donne un isomorphisme $\SheafRG_x(K)\simeq
\Lambda(\delta_X(x))[2\delta_X(x)]$. Pour obtenir l'isomorphisme voulu
$\SheafRG_y(f^!K)\simeq \Lambda(\delta_Y(y))[2\delta_Y(y)]$, il
suffit de définir un isomorphisme $\Lambda\iso \pi^!\Lambda$ : on
utilise le morphisme de Gysin $\Cl \pi$.

\begin{remarque3}\label{17-remarque-transitivite-construction-f-exclam-haut-quasi-fini}
Soit $g\colon Z\to Y$ un autre morphisme quasi-fini. \emph{Via}
l'isomorphisme de transitivité $g^!f^!\simeq (f\circ g)^!$, la structure de
complexe dualisant putatif sur $g^!(f^!K)$ obtenue en appliquant cette
construction à $f$ puis à $g$ est la même que celle obtenue en appliquant
directement la construction à $f\circ g$ : cela résulte aussitôt des
propriétés de composition des morphismes de Gysin.
\end{remarque3}

\begin{proposition3}\label{17-proposition-cdpotentiel-morphisme-fini}
Soit $f\colon Y\to X$ un morphisme fini et surjectif et $K$ un complexe
dualisant putatif sur $X$. Alors $K$ est un complexe dualisant potentiel si
et seulement si le complexe dualisant putatif $f^!K$ de la 
proposition~\ref{17-proposition-cdputatif-quasi-fini} est
un complexe dualisant potentiel.
\end{proposition3}

En vertu des propositions~\ref{17-proposition-specialisation-going-down} et
\ref{17-proposition-specialisation-going-up}, on peut supposer que $X$ et $Y$
sont des schémas locaux strictement henséliens intègres de dimension $1$
et que les fonctions de dimension prennent les valeurs $0$ et $-1$.
Compte tenu de la
remarque~\ref{17-remarque-transitivite-construction-f-exclam-haut-quasi-fini}, il
vient alors qu'il suffit de traiter deux cas :
\begin{enumerate}
\item[(1)] $Y$ est le normalisé de $X$ ;
\item[(2)] $X$ et $Y$ sont des traits.
\end{enumerate}

On obtient la conclusion dans le cas (1) en utilisant le fait que le
morphisme de transition pour $X$ et $K$ est défini à partir de celui pour
$Y$ et $f^\star K$ (cf.
définition~\ref{17-definition-transition-cas-local-dimension-un}) et que si
$M/L$ est une extension finie purement inséparable de corps,
le morphisme de Gysin
$\Cl \pi\colon \Lambda\to\pi^!\Lambda$ associé au morphisme $\pi\colon
\SP(M)\to\SP(L)$ s'identifie à la multiplication par le degré de $M/L$
\emph{via} les isomorphismes tautologiques $\pi^!\Lambda\simeq
\pi^\star\Lambda\simeq \Lambda$.

La démonstration dans le cas (2) va utiliser le lemme général suivant :

\begin{lemme3}\label{17-lemme-compatibilite-cdp-f-exclamation-schemas-reguliers}
Soit $f\colon Y\to X$ un morphisme quasi-fini entre $\ZZn$-schémas
noethériens excellents réguliers de
dimension relative virtuelle $-c$. On munit $X$ de la fonction de dimension
$\delta_X=-\codim$ et $Y$ de la fonction de dimension $\delta_Y$
définie par composition avec $f$ comme dans la 
proposition~\ref{17-proposition-cdputatif-quasi-fini}.
Compte tenu de la relation
$\delta_Y(y)=-\codim_Y(y)-c$ pour tout $y\in Y$,
la proposition~\ref{17-proposition-cdputatif-regulier} munit 
$\Lambda(-c)[-2c]$ d'une structure de complexe dualisant putatif pour la
fonction de dimension $\delta_Y$ sur $Y$. La
proposition~\ref{17-proposition-cdputatif-regulier} donne une structure de
complexe dualisant putatif sur $\Lambda$ sur $X$, dont on déduit, par la
proposition~\ref{17-proposition-cdputatif-quasi-fini}, une structure de
complexe dualisant putatif sur $f^!\Lambda$ sur $Y$ pour la fonction de
dimension $\delta_Y$. Alors, l'isomorphisme de pureté $\Cl
f\colon \Lambda(-c)[-2c]\iso f^!\Lambda$ est compatible aux épinglages
de ces deux complexes dualisants putatifs.
\end{lemme3}

Étudions les épinglages en un point $y\in Y$. Notons $x=f(y)$. On peut
supposer que $X$ et $Y$ sont locaux (strictement henséliens) de points
fermés respectifs $x$ et $y$. On a un diagramme commutatif :
\[
\xymatrix{y\ar[r]^j\ar[d]_g \ar[rd]^h& Y\ar[d]^f \\ x\ar[r]^i & X}
\]
Les schémas apparaissant sur ce diagramme sont affines et réguliers et 
les morphismes entre eux sont de type fini. Ces morphismes sont donc
localement d'intersection complète (lissifiables), on peut leur appliquer
la théorie des morphismes de Gysin. Le résultat du lemme découle alors
aussitôt de leur compatibilité à la composition, puisqu'elle
donne un diagramme commutatif dans $\detaleplus y$, où l'on a
noté $c'$ la codimension de $y$ dans $Y$ :
\[
\xymatrix{\Lambda\ar[r]^{\Cl g} \ar[d]^{\Cl j} & g^!\Lambda
\ar[d]^{g^!(\Cl i)} \\
j^!\Lambda(c')[2c']\ar[r]^-{j^!(\Cl f)}&h^!\Lambda(c+c')[2c+2c']
}
\]

\bigskip

Revenons à la démonstration de la
proposition~\ref{17-proposition-cdpotentiel-morphisme-fini}. Supposons que $K$
soit un complexe dualisant potentiel. D'après la
proposition~\ref{17-proposition-cdp-trait} appliquée à $X$, $K$ est
canoniquement isomorphe à $\Lambda$ (avec les épinglages de la
proposition~\ref{17-proposition-cdputatif-regulier}). On suppose donc
$K=\Lambda$. On a un isomorphisme de pureté
$\Cl f\colon \Lambda\iso f^!\Lambda$.
D'après le lemme~\ref{17-lemme-compatibilite-cdp-f-exclamation-schemas-reguliers}, cet isomorphisme est compatible avec les épinglages de la
proposition~\ref{17-proposition-cdputatif-regulier} pour $\Lambda$ et ceux de
la proposition~\ref{17-proposition-cdputatif-quasi-fini} pour $f^!\Lambda$.
D'après la proposition~\ref{17-proposition-cdp-trait} appliquée à $Y$,
il vient que
$f^!\Lambda$ est bien un complexe dualisant potentiel.

Inversement, supposons que $f^!K$ soit un complexe dualisant potentiel. La
proposition~\ref{17-proposition-cdp-trait} pour $X$
donne un morphisme canonique
$\Lambda\to K$ compatible aux épinglages au point générique de $X$.
Le morphisme $f^!\Lambda\to f^!K$ qui s'en déduit est un morphisme entre
deux complexes dualisants potentiels compatible aux épinglages au point
générique de $Y$.
D'après la proposition~\ref{17-proposition-cdp-trait},
$f^!\Lambda\to f^!K$ est un isomorphisme. Le foncteur $f^!$ étant
conservatif, il en résulte que le morphisme canonique $\Lambda\to K$ est un
isomorphisme. On peut donc supposer que $K=\Lambda$ (de façon compatible
aux épinglages au 
point générique). 
L'épinglage de $K$ au point fermé $x$ ne peut qu'être de la forme
$\lambda\cdot \Cl {x\subset X}^{-1}\colon \SheafRG_x(\Lambda)\iso
\Lambda(-1)[-2]$ pour $\lambda\in\Lambda^\times$. Compte tenu de ce qui
précède, il vient aussitôt
que le complexe dualisant putatif $f^!K$ sur $Y$ s'identifie à $\Lambda$,
épinglé trivialement au point générique et par $\lambda\cdot \Cl {y\subset
Y}^{-1}$ au point fermé $y$ de $Y$. Compte tenu de la
proposition~\ref{17-proposition-cdp-trait}, ce complexe dualisant putatif sur
$Y$ ne peut évidemment être un complexe dualisant potentiel que si
$\lambda=1$. Ainsi, $K$ est bien un complexe dualisant potentiel sur $X$,
puisqu'il s'identifie à $\Lambda$ de façon compatible aux épinglages.

%>>>3
\subsubsection{Un complexe dualisant potentiel}%<<<3

\begin{proposition3}\label{17-proposition-cdpotentiel-regulier}
Soit $X$ un schéma régulier excellent, muni de la fonction de dimension
$-\codim$. Le complexe dualisant putatif $\Lambda$ de la
proposition~\ref{17-proposition-cdputatif-regulier} est un complexe dualisant
potentiel.
\end{proposition3}

Soit $\overline y\to\overline x$ une spécialisation immédiate de points
géométriques de $X$. On veut montrer que les épinglages sur $\Lambda$ sont
compatibles au morphisme de transition $\transition^X_{\overline y\to\overline
x}$. Pour cela, on peut supposer que $X$ est local strictement hensélien de
point fermé $\overline x$. Soit $C$ l'adhérence de l'image de $\overline
y$ dans $X$. Notons $i\colon C\to X$ l'immersion fermée de $C$ dans $X$.
Soit $n\colon \tilde C\to C$ le normalisé de $C$. Montrer la compatibilité
des épinglages avec le morphisme de transition associé à $\overline y\to
\overline x$ revient à montrer que le complexe dualisant putatif
$i^!\Lambda$ de la proposition~\ref{17-proposition-cdputatif-quasi-fini} est
un complexe dualisant potentiel, ce qui, d'après la
proposition~\ref{17-proposition-cdpotentiel-morphisme-fini},
revient encore à dire que
$n^!i^!\Lambda$ en est un. D'après le 
lemme~\ref{17-lemme-compatibilite-cdp-f-exclamation-schemas-reguliers},
le complexe dualisant
putatif $n^!i^!\Lambda$ s'identifie au complexe dualisant putatif
$\Lambda(-c)[-2c]$ obtenu par torsion et décalage à partir de $\Lambda$ sur
le trait $\tilde C$. En vertu de la
proposition~\ref{17-proposition-cdp-trait}, il s'agit bien d'un complexe
dualisant potentiel, ce qui achève la démonstration.

%>>>3
%>>>2
%>>>1
\section{Morphismes de transition généraux et classe de cohomologie en degré maximal}%<<<1
\label{17-section-classe-de-cohomologie}
\subsection{Énoncés des théorèmes principaux}%<<<2

\begin{theoreme2}\label{17-theoreme-classe-degre-maximal}
Soit $X$ un $\ZZn$-schéma local normal strictement 
hensélien excellent de dimension $d$ et de point fermé $x$.
Alors, $\HH^q_x(X,\Lambda(d))=0$ pour $q>2d$ et on a un isomorphisme 
$[x]\colon \Lambda\iso \HH^{2d}_x(X,\Lambda(d))$
compatible aux morphismes de transition associés aux spécialisations
immédiates.
\end{theoreme2}

\begin{theoreme2}\label{17-theoreme-transition-toute-codimension}
Soit $X$ un $\ZZn$-schéma noethérien
excellent. Pour toute spécialisation $\overline
y\to\overline x$ de points géométriques de $X$ de codimension $c$, on peut
définir, pour tout $K\in\detaleplus X$ et $i\in\mathbf{Z}$, un morphisme de
transition $\transition^X_{\overline y\to\overline x}\colon \HH^i_{\overline
y}(K)\to \HH^{i+2c}_{\overline x}(K(c))$, compatible à la composition des
spécialisations et induit par les définitions de la 
section~\ref{17-section-transition-codimension-un} pour $c\leq 1$. Par
ailleurs, ces morphismes de transition généralisés vérifient une
compatibilité avec les morphismes finis, énoncée dans la
proposition~\ref{17-proposition-compatibilite-transition-specialisation}.
\end{theoreme2}

\begin{remarque2}
Localement pour la topologie étale, un schéma quasi-excellent est
excellent. Il est donc évident que dans l'énoncé du
théorème~\ref{17-theoreme-transition-toute-codimension}, on peut remplacer
l'hypothèse \guil{excellent} par \guil{quasi-excellent}.
\end{remarque2}

%>>>2
\subsection{Dimension cohomologique}%<<<2

On énonce ici deux résultats de dimension cohomologique qui se déduisent du
théorème de Lefschetz affine de Gabber,
cf.~\refext{XV}{15-prop-principale-locale}.

\begin{proposition2}\label{17-proposition-dimension-cohomologique}
Soit $\ell$ un nombre premier.
Soit $X$ un $\ZZn[\ell]$-schéma local strictement hensélien excellent 
de point fermé
$x$. On note $U=X-x$. Soit $\mathscr M$ un faisceau de
$\mathbf{Z}/\ell\mathbf{Z}$-modules
sur $U$. Soit $d\in\mathbf{N}$. On suppose que si $u\in U$ est tel que
${\mathscr M}_{\overline u}\neq 0$, alors la dimension de l'adhérence de
$u$ dans $X$ est $\leq d$. Alors $\HH^q(U,\mathscr M)=0$ pour $q\geq 2d$.

En particulier, si $d=\dim X$, on a $\cd_\ell U\leq 2d-1$ et $\cd_\ell
\Gamma_x\leq 2d$.
\end{proposition2}

Par passage à la limite inductive filtrante sur les sous-faisceaux
constructibles de $\mathscr M$, on peut supposer que $\mathscr M$ est
constructible. Quitte à remplacer $X$ par 
l'adhérence dans $X$ du support de $\mathscr M$, on peut supposer que
$d=\dim X$. On est alors ramené à montrer que $\cd_\ell U\leq 2d-1$.
Comme le schéma (séparé)
$U$ peut être recouvert par $d$ ouverts affines (cf.
\cite[§B.3, Chapitre~III]{Algebre@Serre}), et qu'un
ouvert affine de $X$ est de dimension cohomologique au plus 
$d$, en utilisant convenablement les suites exactes de Mayer-Vietoris,
on obtient bien que $\cd_\ell U\leq 2d-1$.

\begin{proposition2}\label{17-proposition-dimension-cohomologique-galois}
Soit $\ell$ un nombre premier.
Soit $X$ un $\ZZn[\ell]$-schéma local strictement hensélien intègre excellent
de dimension $d$, de point générique $\eta$. Alors $\cd_\ell\eta\leq d$.
\end{proposition2}

Le point générique $\eta$ s'identifie à limite projective du système
projectif formé par les ouverts affines non vides de $X$. Comme chacun de
ces ouverts affines est de $\ell$-dimension cohomologique au plus $d$,
il en va de même pour $\eta$.

%>>>2
\subsection{Morphismes finis}%<<<2

\begin{proposition2}\label{17-proposition-compatibilite-transition-spec-immediate}
Soit $f\colon X'\to X$ un morphisme fini entre $\ZZn$-schémas noethériens
excellents. Soit $\overline{y'}\to\overline{x'}$ une spécialisation
immédiate de points géométriques de $X'$ au-dessus d'une spécialisation
immédiate $\overline y\to\overline x$ de points géométriques de $X$
(cf. définition~\ref{17-definition-specialisation-au-dessus-d-une-autre}).
Pour tout $K\in\detaleplus X$ et $i\in\mathbf{Z}$, le diagramme suivant 
commute :
\[
\xymatrix{\HH^i_{\overline{y'}}(f^!K)
     \ar[r]^-{\transition^{X'}_{\overline{y'}\to\overline{x'}}}
  & \HH^{i+2}_{\overline{x'}}(f^!K(1)) \\
\HH^i_{\overline y}(K)
     \ar[r]^-{\transition^{X}_{\overline{y}\to\overline{x}}}
     \ar[u]_{\sim}^{\Cl {y'\to y}}
  & \HH^{i+2}_{\overline{x}}(K(1))\ar[u]^{\sim}_{\Cl {x'\to x}}}
\]
\end{proposition2}

On peut supposer que $f$ est un morphisme dominant entre schémas locaux
strictement henséliens intègres de dimension $1$. Comme dans la
démonstration de la
proposition~\ref{17-proposition-cdpotentiel-morphisme-fini}, il y a deux cas à
traiter :
\begin{enumerate}
\item[(1)] $X'$ est le normalisé de $X$ ;
\item[(2)] $X$ et $X'$ sont des traits.
\end{enumerate}

Le cas (1) étant trivial, on se concentre sur le cas où $X$ et $X'$ sont
des traits. Pour vérifier la compatibilité, on peut évidemment supposer que
$i=0$. Pour tout $K\in\detaleplus X$, comme les flèches verticales sont des
isomorphismes, on peut noter 
$\delta_K\colon \HH^0_{\overline y}(K)\to \HH^2_{\overline x}(K(1))$ la
différence des flèches obtenues en suivant les deux chemins possibles. Il
s'agit de monter que pour tout $K\in\detaleplus X$, on a $\delta_K=0$.

Les résultats de la section~\ref{17-section-premieres-proprietes} montrent que
$\delta_\Lambda=0$. La suite de la démonstration va consister à se ramener
à ce cas-là.

Comme $\overline y$ est au-dessus du
point générique de $X$, $\HH^0_{\overline y}$ s'identifie à la fibre en
$\overline y$ du faisceau de cohomologie de $K$ en degré $0$. Si on note
$\tau_{\leq 0} K\to K$ le morphisme canonique déduit de la $t$-structure
canonique sur $\detaleplus X$\;\footnote{L'auteur pense qu'il eût été plus 
cohérent de noter cette troncature $\tau^{\leq 0}$ pour respecter 
la convention qui veut que les degrés cohomologiques soient indiqués
en exposant, mais la tradition ayant consacré l'usage inverse,
il s'y plie avec répugnance.}, on obtient un carré commutatif :
\[
\xymatrix{
\ar[d]^{\sim}\HH^0_{\overline y}(\tau_{\leq 0}K)\ar[r]^{\delta_{\tau_{\leq
0}K}} &\HH^2_{\overline
x}(\tau_{\leq 0}K(1)) \ar[d]\\
\HH^0_{\overline y}(K)\ar[r]^{\delta_K}& \HH^2_{\overline x}(K(1))
}
\]
Il résulte de ce diagramme que si $\delta_{\tau_{\leq 0}K}=0$, alors
$\delta_K=0$.

Notons ${\mathscr H}^0K$ le faisceau de cohomologie de $K$ en degré zéro.
Pour des raisons de dimension cohomologique, le morphisme canonique
${\mathscr H}^0K\to \tau_{\leq 0}K$ induit un isomorphisme après
application du foncteur $\HH^2_{\overline x}(-(1))$ (et aussi du foncteur
$\HH^0_{\overline y}$). Par conséquent, en considérant un carré commutatif du
type précédent, on obtient cette fois-ci que $\delta_{{\mathscr H}^0K}=0$
équivaut à $\delta_{\tau_{\leq 0}K}=0$.

Il résulte de ces remarques que pour montrer que
$\delta_K=0$  pour tout $K\in\detaleplus X$, on peut supposer que $K$ est
concentré en degré $0$.

On suppose maintenant que $K=\mathscr M$ où $\mathscr M$ est un faisceau de
$\Lambda$-modules sur $X$. Notons $j\colon y\to X$ l'inclusion du point
générique de $X$. En utilisant la description du foncteur $\HH^2_{\overline x}$
de la proposition~\ref{17-proposition-description-h-2-support-trait}, on
obtient que le morphisme canonique $j_!j^\star \mathscr M\to\mathscr M$
induit un isomorphisme après application de $\HH^0_{\overline y}$ et
$\HH^2_{\overline x}(-(1))$ : $\delta_{K}=0$ équivaut à $\delta_{j_!j^\star
K}=0$. La propriété, pour un faisceau $K$ sur $X$, d'être tel que
$\delta_K=0$ ne dépend donc que de sa restriction au point générique
$y$.

Pour tout faisceau de $\Lambda$-modules $\mathscr L$ sur $y$, on
note $\delta_{\mathscr L}=\delta_{j_!\mathscr L}$. Il s'agit de montrer que
$\delta_{\mathscr L}=0$ pour tout faisceau de $\Lambda$-modules sur $y$.
Notons $\mathscr L\to {\mathscr L}_{\constant}$ le plus grand
quotient constant de $\mathscr L$ : formellement, le foncteur $\mathscr
L\longmapsto {\mathscr
L}_{\constant}$ est le foncteur adjoint à gauche du foncteur d'inclusion
de la catégorie des faisceaux constants de $\Lambda$-modules dans la
catégorie des faisceaux de $\Lambda$-modules sur $y$.
En utilisant la
proposition~\ref{17-proposition-description-h-2-support-trait}, on montre
facilement que le morphisme canonique ${\mathscr L}\to{\mathscr
L}_{\constant}$ induit un isomorphisme après application de
$\HH^2_{\overline x}(j_!(1))$ (et une surjection après application
de $\HH^0_{\overline y}$). D'où $\delta_{\mathscr
L}=0$ si et seulement si $\delta_{{\mathscr L}_{\constant}}=0$.
Il en résulte que l'on peut supposer que $\mathscr L$ est un faisceau
constant de $\Lambda$-modules. Par ailleurs, on peut évidemment supposer
que $\mathscr L$ est constructible. Si ${\mathscr L}\to{\mathscr L}'$ est
un épimorphisme de faisceaux constants constructibles de $\Lambda$-modules,
il vient aussitôt que $\delta_{\mathscr L}=0$ implique $\delta_{{\mathscr
L}'}=0$. Il en résulte que l'on peut supposer que $\mathscr L=\Lambda^r$
pour un certain $r\in\mathbf{N}$, puis, par additivité, que $r=1$. Bref, on
s'est bien ramené au cas du faisceau constant $\Lambda$, ce qui achève la
démonstration de cette proposition.

%>>>2
\subsection{Le cas de la \texorpdfstring{dimension~$2$}{dimension 2}}%<<<2
\label{17-subsection-le-cas-de-la-dimension-2}

Compte tenu des résultats établis jusqu'à présent, les
théorèmes~\ref{17-theoreme-classe-degre-maximal} et
\ref{17-theoreme-transition-toute-codimension} peuvent être considérés comme
ayant été établis en dimension $0$ et $1$. Cette sous-section se concentre
sur le cas de la dimension~$2$ qui est l'étape cruciale pour passer au cas
général.

On fixe un $\ZZn$-schéma local strictement
hensélien excellent normal $X$ de dimension~$2$. On note $x$ le point fermé
de $X$ et on pose $U=X-x$. En s'aidant d'une résolution des singularités
$X'\to X$, nous allons définir au
paragraphe~\ref{17-subsubsection-classe-construction} une classe dans
$\HH^4_x(X,\Lambda(2))$, puis nous démontrerons au
paragraphe~\ref{17-subsubsection-classe-independance} qu'elle est indépendante
de la résolution. Enfin, nous établirons au
paragraphe~\ref{17-subsubsection-classe-transition} une compatibilité entre
cette classe, les classes définies pour les localisés stricts en $X$ en
les points de codimension~$1$ et les morphismes de transition associés aux
spécialisations immédiates correspondantes.

\subsubsection{Construction d'une classe}%<<<3
\label{17-subsubsection-classe-construction}

D'après le résultat principal de 
\cite{D2DS@Lipman}, il existe un morphisme birationnel $p\colon X'\to X$ avec
$X'$ régulier. Le morphisme $p$ induit alors automatiquement un
isomorphisme $p^{-1}(U)\iso U$. On peut
supposer de plus que $p^{-1}(x)_{\red}$ est un diviseur à croisements
normaux stricts (dont on note $D_1,\dots,D_n$ les composantes irréductibles).

En effet, on peut dans un premier temps supposer que les composantes
irréductibles $D_1,\dots,D_n$ sont de dimension~$1$ : si ce n'est pas le
cas, on peut éclater le point fermé correspondant (ceci ne se peut produire
que si $X$ est déjà régulier et que $p$ est un isomorphisme). Dans un
deuxième temps, on peut s'arranger pour que les composantes irréductibles
de $p^{-1}(x)_{\red}$ soient régulières en itérant le processus consistant
à éclater les points singuliers (cf. \cite[page~38]{TIFR@Shafarevich}). Dans un
dernier temps, on peut contraindre les croisements à devenir normaux en
éclatant les points fermés récalcitrants ; comme des invariants numériques
décroissent strictement dans cette opération (cf.
\cite[page~21]{TIFR@Shafarevich}), ce processus termine.

\begin{proposition3}\label{17-proposition-classe-sorites}
\begin{enumerate}
\item[(a)] Le morphisme de bord $\HH^{q-1}(U,\Lambda(2))\to
\HH^q_x(X,\Lambda(2))$ est un isomorphisme pour $q\geq 2$ et ces groupes sont
nuls pour $q\geq 5$ ;
\item[(b)] Le morphisme évident $\HH^q(X',\Lambda(2))\to
\HH^q(p^{-1}(x),\Lambda(2))$ est un isomorphisme pour tout $q\in\mathbf{Z}$
et ces groupes sont nuls pour $q\geq 3$ ;
\item[(c)] Le morphisme de bord $\HH^{q-1}(U,\Lambda(2))\to
\HH^q_{p^{-1}(x)}(X',\Lambda(2))$ est un isomorphisme pour $q\geq 4$ ;
\item[(d)] Le morphisme évident $\HH^q_x(X,\Lambda(2))\to
\HH^q_{p^{-1}(x)}(X',\Lambda(2))$ est un isomorphisme pour $q\geq 4$ et ces
groupes sont nuls pour $q\geq 5$.
\end{enumerate}
\end{proposition3}

(a) s'obtient en utilisant la suite exacte canonique :
\[\HH^{q-1}(X,\Lambda(2))\to \HH^{q-1}(U,\Lambda(2))\to \HH^q_x(X,\Lambda(2))\to
\HH^q(X,\Lambda(2))\;\text{.}\]
En effet, $X$ étant local strictement hensélien, on a $\HH^i(X,\Lambda(2))$
pour $i>0$. On conclut en utilisant le fait que $\cd_\ell U\leq 3$ (cf.
proposition~\ref{17-proposition-dimension-cohomologique}).

(b) résulte du théorème de changement de base pour un morphisme propre et
du fait que $p^{-1}(x)$ soit une courbe propre et donc de dimension
cohomologique $2$.

(c) se déduit aussitôt de (b).

(d) résulte de (a), (c) et de la commutativité du diagramme évident
suivant :
\[
\xymatrix{\HH^{q-1}(U,\Lambda(2))\ar[r]\ar@{=}[d] & \HH^q_x(X,\Lambda(2))\ar[d] \\
\HH^{q-1}(U,\Lambda(2))\ar[r] & \HH^q_{p^{-1}(x)}(X',\Lambda(2))}
\]

\begin{proposition3}\label{17-proposition-classe-suite-exacte-1}
On a une suite exacte 
\[
\bigoplus_{i<j} \HH^4_{D_i\cap D_j}(X',\Lambda(2))\to 
\bigoplus_i \HH^4_{D_i}(X',\Lambda(2))\to
\HH^4_{p^{-1}(x)}(X',\Lambda(2))\to 0\;\text{,}
\]
où les flèches sont induites par des morphismes d'agrandissement du
support, et leurs différences.
\end{proposition3}

Pour tout fermé $F$ de $X'$, on note $\Lambda_F$ le faisceau
$i_{F\star}\Lambda$
où $i_F$ est l'immersion de $F$ dans $X'$. On a une suite exacte courte
évidente de faisceaux de $\Lambda$-modules sur $X'$ :
\[0\to \bigoplus_{i<j} \Lambda_{D_i\cap D_j}\to 
\bigoplus_i \Lambda_{D_i}\to \Lambda_{p^{-1}(x)}\to 0\;\text{.}\]

En appliquant $\R\SheafHom(-,\Lambda(2))$ au triangle distingué de
$\detaleplus {X'}$ associé à cette suite exacte courte, on peut obtenir la
suite exacte voulue, pourvu que l'on sache montrer que $\HH^5_{D_i\cap
D_j}(X',\Lambda(2))$ est nul pour $i<j$.
D'après le théorème de pureté absolue, ce
groupe s'identifie à $\HH^1(D_i\cap D_j,\Lambda)$ qui est bien nul puisque
$D_i\cap D_j$ est une union disjointe finie de spectres de corps
séparablement clos.

\begin{proposition3}\label{17-proposition-classe-suite-exacte-2}
La suite exacte de la proposition~\ref{17-proposition-classe-suite-exacte-1}
se récrit sous la forme :
\[\bigoplus_{i<j}\bigoplus_{y\in D_i\cap D_j}\Lambda\vers {\delta}
\Lambda^n \vers\varepsilon 
\HH^4_{p^{-1}(x)}(X',\Lambda(2))\to 0\;\text{.}\]
Si on note $(\chi_y)_{i<j,y\in D_i\cap D_j}$ la base canonique du groupe de
gauche et $\chi_{D_1},\dots,\chi_{D_n}$ celle de $\Lambda^n$, la
différentielle $\delta$ vérifie la formule
\[\delta(\chi_y)=[y:x]\cdot \left(\chi_{D_i}-\chi_{D_j}\right)\;\text{.}\]
En outre, $\varepsilon(\chi_{D_i})\in \HH^4_{p^{-1}(x)}(X',\Lambda(2))$ est
obtenu par agrandissement du support à partir de la classe
$\frac{1}{[y:x]}\Cl{y\to X}\in \HH^4_y(X',\Lambda(2))$
pour tout point fermé $y$ d'un $D_i$.
\end{proposition3}

Le théorème de pureté cohomologique absolue donne des isomorphismes
\[\Cl{D_i\cap D_j\to X}\colon \HH^0(D_i\cap D_j,\Lambda)\iso \HH^4_{D_i\cap
D_j}(X',\Lambda(2))\;\text{,}\]
ce qui permet de décrire le groupe de gauche dans la suite exacte de la
proposition~\ref{17-proposition-classe-suite-exacte-1}.
Le même théorème donne aussi des isomorphismes \[\Cl{D_i\to X}\colon
\HH^2(D_i,\Lambda(1))\iso \HH^4_{D_i}(X',\Lambda(2))\;\text{.}\]
En outre, on dispose
du morphisme trace (relativement à $x$) $\tr\colon \HH^2(D_i,\Lambda(1))
\iso \Lambda$. Celui-ci est caractérisé par le fait que pour tout point fermé
$y$ de $D_i$, l'image de $\Cl{y\to D_i}\in \HH^2_y(D_i,\Lambda(1))$ par le
morphisme composé \[\HH^2_y(D_i,\Lambda(1))\to \HH^2(D_i,\Lambda(1))\vers{\tr}
\Lambda\]
soit $[y:x]$.
La description des morphismes $\delta$ et $\varepsilon$ donnée dans
l'énoncé s'obtient alors aussitôt à partir des propriétés de composition des
morphismes de Gysin.

\begin{proposition3}\label{17-proposition-classe-suite-exacte-3}
La suite suivante est exacte :
\[\bigoplus_{i<j}\bigoplus_{y\in D_i\cap D_j}\Lambda\vers {\delta}
\bigoplus_i \Lambda \vers\Sigma \Lambda\to 0\;\text{,}\]
où $\Sigma$ est défini par $\Sigma(\chi_{D_i})=1$ pour tout $i$.
\end{proposition3}

Tout d'abord, à partir de la formule donnée pour $\delta$ dans la 
proposition~\ref{17-proposition-classe-suite-exacte-2}, il est clair que
$\Sigma\circ \delta=0$. $\Sigma$ induit donc un morphisme
$\coker \delta\to \Lambda$.
Comme les degrés $[y:x]$ qui
interviennent sont inversibles dans $\Lambda$, on voit que si $i$ et $j$
sont tels que $D_i\cap D_j$ soit non vide, alors $\chi_{D_i}$ et
$\chi_{D_j}$ ont la même classe dans $\coker \delta$. La fibre $p^{-1}(x)$
étant connexe (Main Theorem), on en déduit que tous les éléments
$\chi_{D_i}$ ont la même classe dans $\coker \delta$. L'inclusion entre
l'image de $\delta$ et le noyau de $\Sigma$ est donc une égalité.

\begin{corollaire3}
Les suites exactes des propositions~\ref{17-proposition-classe-suite-exacte-2}
et \ref{17-proposition-classe-suite-exacte-3} donnent naissance à un
isomorphisme $\Lambda\iso \HH^4_{p^{-1}(x)}(X';\Lambda(2))$.
\end{corollaire3}

\begin{corollaire3}
\emph{Via} l'isomorphisme canonique $\HH^4_{x}(X,\Lambda(2))\iso
\HH^4_{p^{-1}(x)}(X',\Lambda(2))$ de la
proposition~\ref{17-proposition-classe-sorites}~(d), l'isomorphisme du
corollaire précédent donne un isomorphisme $\Lambda\iso
\HH^4_x(X,\Lambda(2))$.
\end{corollaire3}

\begin{definition3}
On note $[x]_{X'}\in \HH^4_x(X,\Lambda(2))$ le générateur défini par
l'isomorphisme du corollaire précédent.
\end{definition3}

%>>>3
\subsubsection{Indépendance en la résolution}%<<<3
\label{17-subsubsection-classe-independance}

Dans ce paragraphe, nous allons montrer que si $q\colon X''\to X$ est une
autre résolution du type envisagé dans le
paragraphe~\ref{17-subsubsection-classe-construction}, alors
$[x]_{X'}=[x]_{X''}$. Quitte à introduire une désingularisation de la
composante irréductible dominant $X$ du produit
fibré de $X''$ de $X'$ au-dessus de $X$, on peut supposer qu'une des deux
désingularisations $X'$ et $X''$ considérées coiffe l'autre. On suppose
donc par exemple qu'il existe un (unique) $X$-morphisme $\pi\colon X''\to
X'$. Le morphisme $\pi$ est projectif et birationnel entre deux schémas
réguliers de dimension $2$, il induit donc un isomorphisme au-dessus d'un
ouvert $U'$ de $X'$ tel que le fermé $X'-U'$ soit de dimension $0$. Il
existe donc certainement un point fermé $y'$ de $p^{-1}(x)$ tel que le
morphisme induit soit un isomorphisme $\pi^{-1}(y')\to y'$. On note $y''$
l'unique point fermé de $\pi^{-1}(y')$. La compatibilité des classes de
Gysin au changement de base implique que la classe $\Cl{y'\to X'}$ est
envoyée sur $\Cl{y''\to X''}$ par le morphisme de restriction
$\pi^\star\colon \HH^4_{y'}(X',\Lambda(2))\to \HH^4_{y''}(X'',\Lambda(2))$.
Cette compatibilité vaut encore après élargissement du support à
$p^{-1}(x)$ et à $q^{-1}(x)$. En considérant la composition suivante :
\[\HH^4_{x}(X,\Lambda(2))\iso \HH^4_{p^{-1}(x)}(X',\Lambda(2))\iso
\HH^4_{q^{-1}(x)}(X'',\Lambda(2))\;\text{,} \]
on obtient aussitôt que $[y':x]\cdot [x]_{X'}=[y'':x]\cdot [x]_{X''}$, ce
qui permet de conclure que $[x]_{X'}=[x]_{X''}$.

\begin{remarque3}
On peut montrer que si $p\colon X'\to X$ est un morphisme projectif
birationnel avec $X'$ régulier, alors, même sans supposer que
$p^{-1}(x)_\red$ soit un diviseur à croisements normaux strict, 
l'application $\HH^4_x(X,\Lambda(2))\to
\HH^4_{p^{-1}(x)}(X',\Lambda(2))$ est un isomorphisme et la classe $[x]$
s'envoie bien sur l'élément induit par 
$\frac 1 {[x':x]}\Cl{x'\to X'}$ pour tout
point fermé $x'$ de $p^{-1}(x)$.
\end{remarque3}

%>>>3
\subsubsection{Compatibilité aux morphismes de transition}%<<<3
\label{17-subsubsection-classe-transition}

Soit $\overline y\to X$ un point géométrique au-dessus d'un point $y$ de
codimension $1$. Autrement dit, on a une spécialisation immédiate
$\overline y\to x$ de points géométriques de $X$. Nous allons montrer que
si $\overline\eta\to\overline y$ est une spécialisation immédiate, alors on
a l'égalité
\[ [x]=\transition^X_{\overline y\to x}(\transition^X_{\overline\eta\to\overline x}(1))\]
dans $\HH^4_x(X,\Lambda(2))$, ce qui achèvera la démonstration du
théorème~\ref{17-theoreme-classe-degre-maximal} jusqu'en dimension~$2$.

On note $C$ l'adhérence de $y$ dans $X$. Il existe une désingularisation
$X'\to X$ du type envisagé dans le 
paragraphe~\ref{17-subsubsection-classe-construction} telle que l'adhérence
$C'$ de $y$ dans $X'$ soit un trait. Plus précisément, le morphisme évident
$C'\to C$ identifie $C'$ au normalisé de $C$. Le morphisme $C'\to C$ étant
un homéomorphisme universel, on peut noter $x'$
le point fermé de $C'$ (celui de
$C$ est bien entendu $x$) et $U'=X'-x'$.
On dispose d'immersions fermées évidentes
$i\colon C\to X$, $i'\colon C'\to X'$,
$k\colon y\to U$ et $k'\colon y\to U'$.

\begin{lemme3}\label{17-lemme-hexagone-commutatif}
Avec les notations ci-dessus, le diagramme évident qui suit
est commutatif (les
flèches marquées comme étant des isomorphismes devant être considérées
comme bidirectionnelles) :
\[
\xymatrix{
&\ar[ld]^{\sim} \HH^4_x(X,\Lambda(2))\ar@{=}[r] & \HH^4_x(C,i^!\Lambda(2)) & \\ 
\HH^4_{p^{-1}(x)}(X',\Lambda(2)) & & & \HH^3(y,k^!\Lambda(2))
\ar[ul]^{\sim}\ar[ld]\\
& \ar[lu]^{\sim}
\HH^4_{x'}(X',\Lambda(2))\ar@{=}[r] & \HH^4_{x'}(C',{i'}^!\Lambda(2)) & \\
}
\]
\end{lemme3}

Ceci équivaut à la commutativité du diagramme évident :
\[
\xymatrix{
\HH^3(U,\Lambda(2))\ar[r]^\sim & \HH^4_x(X,\Lambda(2))\ar@{=}[r] &
   \HH^4_x(C,i^!\Lambda(2)) & \ar[l] \HH^3(y,k^!\Lambda(2))\ar@{=}[d] \\
\ar[u]_{\sim}\HH^3(U',\Lambda(2))\ar[r]^\sim & \HH^4_{x'}(X',\Lambda(2))\ar@{=}[r] &
   \HH^4_{x'}(C',{i'}^!\Lambda(2)) & \ar[l] \HH^3(y,{k'}^!\Lambda(2)) 
}
\]

On peut identifier le carré externe de ce diagramme-ci au suivant, où les
flèches horizontales sont les flèches d'oubli du support :
\[
\xymatrix{
\HH^3(U,\Lambda(2)) & \ar[l] \HH^3_y(U,\Lambda(2)) \ar@{=}[d] \\
\ar[u]_{\sim} \HH^3(U',\Lambda(2)) & \ar[l] \HH^3_y(U',\Lambda(2))
}
\]

Ce diagramme-là est bien évidemment commutatif, ce qui achève la
démonstration du lemme~\ref{17-lemme-hexagone-commutatif}.

\bigskip

On peut prolonger sur la droite le diagramme commutatif du lemme :
\[
\xymatrix{
\HH^4_x(C,i^!\Lambda(2)) &
 \ar[l]_-{\transition^C_{\overline y\to x}} \HH^2_{\overline y}(C,i^!\Lambda(1)) & &\\
\HH^3(y,k^!\Lambda(2))\ar[u]^\sim\ar[d]^\sim  & &
({\mathscr H}^2 k^!\Lambda(1))_{\overline
y}\ar[dl]_{\sim}^1\ar[ul]_{\sim}^{[x':x]} &
\ar[l]_-{\transition^X_{\overline\eta\to\overline y}}^-{\sim}\Lambda\\
\HH^4_{x'}(C',{i'}^!\Lambda(2)) 
& \ar[l]_-{\transition^{C'}_{\overline y\to x}} \HH^2_{\overline
 y}(C',{i'}^!\Lambda(1)) & & }
\]
Du groupe $({\mathscr H}^2 k^!\Lambda(1))_{\overline y}$ partent deux
isomorphismes canoniques vers des groupes de cohomologie à supports de $C$
et de $C'$ ; ici, on les a multipliés respectivement par $[x':x]$ et $1$ de
façon à rendre le diagramme commutatif (cf.
définition~\ref{17-definition-transition-cas-local-dimension-un}).

On peut partir de l'élément $1$ dans le groupe $\Lambda$ tout à droite et
considérer son image dans le groupe $\HH^4_x(X,\Lambda(2))$ figurant sur le
diagramme du lemme. En suivant le chemin du bas, on obtient
$[x':x]\cdot [x]_{X'}$.
En suivant le chemin du haut, on obtient
$[x':x]\cdot \transition^X_{\overline y\to x}
(\transition^X_{\overline\eta\to\overline y}(1))$. On peut ainsi conclure que
l'on a bien l'egalité
\[[x]_{X'}=\transition^X_{\overline y\to x}
(\transition^X_{\overline\eta\to\overline y}(1))\]
dans $\HH^4_x(X,\Lambda(2))$.

%>>>3
%>>>2
\subsection{Morphismes de transition en codimension arbitraire}%<<<2

Nous allons maintenant démontrer le
théorème~\ref{17-theoreme-transition-toute-codimension} en nous appuyant sur
le théorème~\ref{17-theoreme-classe-degre-maximal} établi pour le moment
jusqu'en dimension $2$. Ce résultat nous permettra ensuite d'établir le
théorème~\ref{17-theoreme-classe-degre-maximal} en toute généralité.

\begin{definition2}
Soit $X$ un $\ZZn$-schéma noethérien excellent.
Soit $\overline y=\overline {x_0}\to \dots \to \overline {x_n}=\overline x$ une
suite de spécialisations de points géométriques de $X$ telle que pour tout
$0\leq i<n$, la spécialisation $\overline x_i\to \overline x_{i+1}$ soit de
codimension $0$ ou $1$. On note $c$ la codimension de la spécialisation
$\overline y\to \overline x$.
Pour tout $K\in\detaleplus X$, on
note $\transition^X_{\overline x_0\to\dots \to
\overline x_n}\colon \HH^p_{\overline y}(K)\to 
\HH^{p+2c}_{\overline x}(K(c))$ le morphisme de transition obtenu par la
composition
$\transition^X_{\overline x_{n-1}\to\overline x_n}\circ\dots\circ \transition^X_{\overline x_0\to\overline x_1}$.
\end{definition2}

\begin{definition2}\label{17-definition-propriete-c}
Soit $X$ un $\ZZn$-schéma noethérien excellent. Soit $\overline y\to
\overline x$ une spécialisation de points géométriques de $X$ de
codimension $c$. Soit $K\in\detaleplus X$. On dit que la propriété
$\propC^c_{\overline y\to \overline x,K}$ est satisfaite si le morphisme
$\transition^X_{\overline {x_0}\to\dots\to\overline {x_n}}\colon
\HH^0_{\overline y}(K)\to \HH^{2c}_{\overline x}(K(c))$
ne dépend pas du choix de la factorisation de $\overline y\to\overline x$
en $\overline {x_0}\to\dots\to\overline {x_n}$.
\end{definition2}

\begin{definition2}
On dira que la propriété $\propC^{\leq c}$ est satisfaite si toutes les
propriétés $\propC^{c'}_{\overline y\to\overline x,K}$ envisagés dans
la définition~\ref{17-definition-propriete-c} le sont pour $c'\leq c$.
On dira que la
propriété $\propC^{\leq c}_{\loc}$ est vérifiée si les propriétés
$\propC^{c'}_{\overline \eta\to x,K}$ sont vérifiées 
dans la situation, dite locale, où le schéma $X$ est supposé 
local strictement hensélien intègre de dimension $\leq c$
de point fermé $x$ et où $\overline \eta$ est au-dessus du point générique
de $X$. Enfin, on dira que la propriété $\propC^{\leq
c}_{\loc,\normal,\Lambda}$ est satisfaite si on suppose de plus que $X$ est
normal et que $K=\Lambda$.
\end{definition2}

Il s'agit donc d'établir la propriété $\propC^{\leq c}$ pour tout $c\geq
0$. D'après la
remarque~\ref{17-remarque-composition-transition-codimension-0-ou-1}, on peut
considérer que la propriété $\propC^{\leq 1}$ est connue. En outre, notons
que le théorème~\ref{17-theoreme-transition-toute-codimension} affirme en
particulier $\propC^{\leq c}_{\loc,\normal,\Lambda}$ pour tout $c\geq 0$. Les
résultats de la sous-section~\ref{17-subsection-le-cas-de-la-dimension-2}
montrent que $\propC^{\leq 2}_{\loc,\normal,\Lambda}$ est satisfaite.

\begin{proposition2}\label{17-proposition-compatibilite-transition-specialisation}
Soit $f\colon X'\to X$ un morphisme fini entre $\ZZn$-schémas noethériens
excellents. Soit $\overline{y'}\to\overline{x'}$ une spécialisation
de points géométriques de $X'$ de codimension $c$
au-dessus d'une spécialisation
$\overline y\to\overline x$ de points géométriques de $X$
(cf. définition~\ref{17-definition-specialisation-au-dessus-d-une-autre}).
Pour toute factorisation $\overline{x'_0}\to\dots\overline{x'_n}$ de
$\overline {y'}\to\overline{x'}$ en une suite de spécialisations de
codimension $0$ ou $1$, si on note $\overline{x_0}\to\dots\overline{x_n}$
la factorisation de $\overline y\to\overline x$ en-dessous de la
précédente, pour tout $K\in\detaleplus X$ et $i\in\mathbf{Z}$,
le diagramme suivant commute :
\[
\xymatrix{\HH^i_{\overline{y'}}(f^!K)
     \ar[rr]^-{\transition^{X'}_{\overline{x'_0}\to\dots\to\overline{x'_n}}}
&  & \HH^{i+2c}_{\overline{x'}}(f^!K(c)) \\
\HH^i_{\overline y}(K)
     \ar[rr]^-{\transition^{X}_{\overline{x_0}\to\dots\to\overline{x_n}}}
     \ar[u]_{\sim}^{\Cl {y'\to y}}
&  & \HH^{i+2c}_{\overline{x}}(K(c))\ar[u]^{\sim}_{\Cl {x'\to x}}}
\]
En outre, on a l'équivalence
$\propC^c_{\overline{y}\to\overline{x},K}\Longleftrightarrow
\propC^c_{\overline{y'}\to\overline{x'},f^!K}$.
\end{proposition2}

La commutativité du diagramme se déduit aussitôt de la 
proposition~\ref{17-proposition-compatibilite-transition-spec-immediate}.
L'équivalence annoncée résulte de la commutativité du diagramme et du fait
qu'il est essentiellement équivalent de 
se donner une factorisation de la spécialisation
$\overline{y'}\to\overline{x'}$ ou de s'en donner une de
$\overline y\to\overline x$ (cf.
propositions~\ref{17-proposition-specialisation-going-down} et
\ref{17-proposition-specialisation-going-up}).

\begin{lemme2}\label{17-lemme-reduction-propriete-c-situation-locale}
Pour tout entier $c\geq 0$, on a l'équivalence $\propC^{\leq
c}\Longleftrightarrow \propC^{\leq c}_{\loc}$.
\end{lemme2}

Supposons $\propC^{\leq c}_{\loc}$ et montrons $\propC^{\leq c}$.
Il suffit évidemment de montrer
$\propC^{c'}_{\overline y\to\overline x,K}$ pour tout $K\in\detaleplus X$, avec $X$
un schéma local strictement hensélien de point fermé $\overline x$ et une
spécialisation $\overline y\to\overline x$ de
codimension $c'\leq c$. Il s'agit de vérifier que
l'on peut supposer que $X$ est intègre et que
$\overline y$ est au-dessus du point générique. Pour cela, on introduit
l'immersion fermée $i\colon Z\to X$ où $Z=\overline{\left\{y\right\}}$. La
proposition~\ref{17-proposition-compatibilite-transition-specialisation}
montre que $\propC^{c'}_{\overline y\to\overline x,i^!K}$ implique
$\propC^{c'}_{\overline y\to\overline x,K}$. Pour conclure, il suffit
d'observer que $\propC^{c'}_{\overline y\to\overline x,i^!K}$ est un cas
particuler de $\propC^{\leq c}_{\loc}$.

\begin{lemme2}\label{17-lemme-propriete-c-reductions}
Soit $X$ un $\ZZn$-schéma noethérien excellent local strictement hensélien
intègre de dimension $c$,
de point fermé $x$ et de point générique $\eta$. Soit $\overline
\eta$ un point géométrique au-dessus de $\eta$. Soit $K\in\detaleplus X$.
\begin{enumerate}
\item[(a)] $\propC^c_{\overline\eta\to x,\tau_{\leq 0}K}\Longrightarrow
\propC^c_{\overline\eta\to x,K}$ ;
\item[(b)] $\propC^c_{\overline\eta\to x,\tau_{\leq 0}K}\Longleftrightarrow
\propC^c_{\overline\eta\to x,{\mathscr H}^0K}$ ;
\item[(c)] Soit $j\colon U\to X$ l'inclusion d'un ouvert dense
et $\mathscr M$ et $\mathscr N$ deux faisceaux de
$\Lambda$-modules sur $X$ tels que $j^\star \mathscr M\simeq j^\star
\mathscr N$, alors 
$\propC^c_{\overline\eta\to x,\mathscr M}\Longleftrightarrow
\propC^c_{\overline\eta\to x,\mathscr N}$.
\end{enumerate}
\end{lemme2}

L'implication (a) résulte du fait que le morphisme canonique $\tau_{\leq
0}K\to K$ induise un isomorphisme $\HH^0_{\overline \eta}(\tau_{\leq
0}K)\iso \HH^0_{\overline \eta}(K)$.

On considère ensuite le morphisme canonique
$\tau_{\leq 0}K\to {\mathscr H}^0 K$. Il induit évidemment un isomorphisme
après application de $\HH^0_{\overline \eta}$, mais aussi après celle de
$\HH^{2c}_{x}$, pour des raisons de dimension cohomologique (cf.
proposition~\ref{17-proposition-dimension-cohomologique}). L'équivalence (b)
en résulte aussitôt.

Pour montrer (c), on peut supposer que $\mathscr N=j_!j^\star \mathscr M$.
On a alors un monomorphisme canonique $j_!j^\star \mathscr M\to\mathscr M$,
dont le conoyau est $i_\star i^\star\mathscr M$, où $i\colon Z\to X$
désigne une immersion fermée complémentaire. La
proposition~\ref{17-proposition-dimension-cohomologique} montre alors que
$\HH^q(X-x,(i_\star i^\star\mathscr M)_{|X-x})=0$ pour $q\geq 2c-2$. Ainsi,
le morphisme $j_!j^\star \mathscr M\to\mathscr M$ induit un isomorphisme
non seulement après application de $\HH^0_{\overline \eta}$, mais aussi celle
de $\HH^{2c}_x$, ce qui permet de conclure.

\begin{lemme2}\label{17-lemme-reduction-propriete-c-faisceau-constant}
Pour tout entier $c\geq 0$, on a l'implication $\propC^{\leq
c}_{\loc,\normal,\Lambda}\Longrightarrow \propC^{\leq c}_{\loc}$.
\end{lemme2}

Soit $X$ un $\ZZn$-schéma noethérien excellent
local strictement hensélien intègre de dimension $c'\leq c$, de point
fermé $x$ et de point générique $\eta$. Soit $\overline\eta$ un point
géométrique au-dessus de $\eta$. D'après le 
lemme~\ref{17-lemme-propriete-c-reductions}, pour montrer 
$\propC^{c'}_{\overline\eta\to x,K}$ pour tout $K\in\detaleplus X$, il suffit
de montrer $\propC^{c'}_{\overline\eta\to x,\mathscr M}$ pour tout faisceau
de $\Lambda$-modules $\mathscr M$ sur $X$. On peut évidemment supposer que
$\mathscr M$ est constructible. Mais alors, il existe un ouvert dense de
$X$ au-dessus duquel $\mathscr M$ soit localement constant. Il existe
évidemment un morphisme fini surjectif $f\colon \tilde X\to X$, avec
$\tilde X$ normal intègre, et un ouvert dense $U$ de $X$ tel que $f$
induise un revêtement fini étale $\tilde U=f^{-1}(U)\to U$ et que
$(f^\star \mathscr M)_{|\tilde U}$ soit isomorphe à un faisceau constant de
valeur un certain $\Lambda$-module $N$. Choisissons une spécialisation
$\overline {\tilde \eta}\to \tilde x$ au-dessus de
$\overline\eta\to x$. La propriété $\propC^{\leq c}_{\loc,\normal,\Lambda}$
contient $\propC^{c'}_{\overline {\tilde \eta}\to \tilde x,\Lambda}$
comme cas particulier. Cette dernière propriété
implique à son tour la propriété 
$\propC^{c'}_{\overline {\tilde \eta}\to \tilde x,N}$. On observe que l'on
a un isomorphisme canonique $(f^!\mathscr M)_{|\tilde U}\simeq N$.
L'assertion (c) du lemme~\ref{17-lemme-propriete-c-reductions} permet
d'obtenir la propriété $\propC^{c'}_{\overline {\tilde \eta}\to \tilde
x,{\mathscr H}^0(f^!\mathscr M)}$, et les assertions (a) et (b) d'en
déduire $\propC^{c'}_{\overline {\tilde \eta}\to \tilde
x,f^!\mathscr M}$. Enfin, la 
proposition~\ref{17-proposition-compatibilite-transition-specialisation}
permet de vérifier $\propC^{c'}_{\overline\eta\to x,\mathscr M}$, ce
qui achève la démonstration du lemme.


\begin{lemme2}\label{17-lemme-propriete-c-etape-de-recurrence}
Pour tout entier $c\geq 3$, on a l'implication $\propC^{\leq
c-1}\Longrightarrow \propC^{\leq c}_{\loc,\normal,\Lambda}$.
\end{lemme2}

Soit $X=\SP(A)$
un schéma noethérien excellent local strictement hensélien normal
de dimension $c$, de point fermé $x$ et de point générique $\eta$. On
choisit un point géométrique $\overline \eta$ au-dessus de $\eta$. On note
$s\colon \overline\eta\to x$ la spécialiation canonique.
Nous
allons montrer la propriété $\propC^c_{s,\Lambda}$.

Soit $\overline z\to X$ un point géométrique tel que $z\neq x$ et $z\neq
\eta$. On a une spécialisation canonique $t\colon \overline z\to x$. On peut
choisir une spécialisation $u\colon \overline \eta\to\overline z$. On obtient
peut-être pas ainsi une factorisation de la spécialisation
$\overline\eta\to x$ fixée plus haut, mais il existe certainement
$\sigma\in\Gal(\overline\eta/\eta)$ faisant commuter le diagramme suivant
de spécialisations de points géométriques de $X$ :
\[
\xymatrix{\overline\eta\ar[r]^s\ar[d]^{\sigma}_\sim & x\\ 
\overline\eta \ar[r]^u & \overline z\ar[u]^t}
\]
Les spécialisations $\sigma$, $u$ et $t$
sont de codimension $\leq c-1$, les morphismes de transitions
associés sont donc bien définis.
On peut considérer l'image de $1\in \HH^0_{\overline\eta}(X,\Lambda)$ par le
composé de ces morphismes de transition :
\[\gamma_{\sigma,u,t}=(\transition^X_{t}\circ\transition^X_u\circ\transition^X_\sigma)(1)\in
\HH^{2c}_x(X,\Lambda(c))\;\text{.}\]
Bien entendu, on a $\transition^X_\sigma(1)=1$. On note donc
simplement $\gamma_{u,t}$ l'élément
$\gamma_{\sigma,u,t}=(\transition^X_{t}\circ\transition^X_u)(1)$. 
J'affirme que cette classe $\gamma_{u,t}$
ne dépend que du point $z$ de $X$ en-dessous duquel $\overline z$ se
trouve. En effet, soit $\overline {z'}$ un autre point géométrique
au-dessus de $z$, et $t'\colon\overline{z'}\to x$ la spécialisation
canoniquement associée. On choisit une spécialisation $u'\colon
\overline\eta\to \overline{z'}$. On peut choisir un $z$-isomorphisme 
$\sigma'\colon \overline z\iso \overline {z'}$.
Le schéma $X$ étant géométriquement unibranche, on peut montrer facilement
qu'il existe un $\eta$-isomorphisme $\sigma''\colon
\overline\eta\to\overline\eta$ induisant 
un diagramme commutatif de spécialisations de points
géométriques de $X$ :
\[
\xymatrix{
\overline\eta\ar[r]^u\ar[d]^{\sigma''}_\sim&
\overline{z}\ar[d]^{\sigma'}_\sim\ar[r]^t & x\ar@{=}[d] \\
\overline\eta\ar[r]^{u'}& \overline{z'}\ar[r]^{t'} & x
}
\]
En utilisant cette fois-ci que $\sigma''$ agit trivialement sur $1\in
\HH^0_{\overline \eta}(X,\Lambda)$, on montre que 
$\gamma_{u,t}=\gamma_{u',t'}$. On peut donc noter
simplement $\gamma_z$ cette classe.

Il s'agit de montrer que la classe $\gamma_z$ est indépendante de $z\in
X-\left\{\eta,x\right\}$. Par construction, si 
$z'\in X-\left\{\eta,x\right\}$ est tel qu'il existe une spécialisation
(Zariski) $z\to z'$, on a $\gamma_z=\gamma_{z'}$.

Comme $A$ est local normal excellent de dimension $\geq 3$, d'après
\sga{2}{xiii}{2.1}, si $f\in A-\left\{0\right\}$, le schéma $\SP\left(
A/(f)-\left\{x\right\}\right)$ est connexe. Il en résulte aussitôt que les
classes $\gamma_z$ pour $z\in \SP\left(A/(f)-\left\{x\right\}\right)$
sont égales.
Maintenant, si $z$ et $z'$ sont deux éléments de $X-\left\{\eta,x\right\}$,
il existe bien évidemment $f\in A-\left\{0\right\}$ tel que $z$ et $z'$
appartiennent à l'hypersurface définie par $f$. On obtient donc bien
$\gamma_z=\gamma_{z'}$, ce qui achève la démonstration du lemme.

\bigskip

Nous sommes maintenant en mesure de démontrer le
théorème~\ref{17-theoreme-transition-toute-codimension}. Il s'agit d'établir
la propriété $\propC^{\leq c}$ pour tout $c\geq 0$. Elle équivaut à la
propriété $\propC^{\leq c}_{\loc}$ d'après le
lemme~\ref{17-lemme-reduction-propriete-c-situation-locale}. D'après le
lemme~\ref{17-lemme-reduction-propriete-c-faisceau-constant}, elle équivaut
encore à la propriété 
$\propC^{\leq c}_{\loc,\normal,\Lambda}$. Compte tenu de ces
équivalences, le
lemme~\ref{17-lemme-propriete-c-etape-de-recurrence} montre que pour $c\geq
3$, $\propC^{\leq c-1}_{\loc,\normal,\Lambda}$ implique 
$\propC^{\leq c}_{\loc,\normal,\Lambda}$. Un raisonnement par récurrence
permet donc de conclure puisque la propriété $\propC^{\leq
2}_{\loc,\normal,\Lambda}$ a été établie plus haut.

%>>>2
\subsection{Fin de la démonstration}%<<<2

\begin{proposition2}
Soit $X$ un $\ZZn$-schéma noethérien excellent 
local hensélien normal de dimension $d$, de point fermé $x$ et
de point générique $\eta$. Soit $\overline x$ un point géométrique
au-dessus de $x$. Soit $\overline \eta$ un point géométrique au-dessus de
$\eta$. Soit $\overline\eta\to\overline x$ une spécialisation.
Alors, l'image de $1$ par le morphisme de transition
$\transition^X_{\overline\eta \to \overline x}\colon
\HH^0_{\overline \eta}(X,\Lambda)\to \HH^{2d}_{\overline x}(X,\Lambda(d))$ est
invariante par $\Gal(\overline x/x)$, et indépendante de la spécialisation
choisie $\overline \eta\to \overline x$, on la note $[x]$.
On obtient ainsi un morphisme
$\Lambda[-2d]\to \tau_{\geq 2d}\SheafRG_x(\Lambda(d))$
dans $\detale x$.
\end{proposition2}

Ceci résulte aussitôt du
théorème~\ref{17-theoreme-transition-toute-codimension} et du fait que $1\in
\HH^0_{\overline\eta}(X,\Lambda)$ est fixé par $\Gal(\overline\eta/\eta)$.

\bigskip

Par construction, les classes $[x]$ sont compatibles aux morphismes de
transition. Pour finir de démontrer le
théorème~\ref{17-theoreme-classe-degre-maximal}, il reste à montrer que dans
la situation de la proposition précédente, le morphisme $\Lambda[-2d]\to
\tau_{\geq 2d}\SheafRG_x(\Lambda(d))$ est un isomorphisme, ou
encore, dans la situation où $X$ est strictement hensélien, que le
morphisme $\Lambda\to \HH^{2d}_x(X,\Lambda(d))$ induit par $[x]$ est un
isomorphisme. En effet, des considérations de dimension cohomologique (cf.
proposition~\ref{17-proposition-dimension-cohomologique}) expliquent
l'annulation de $\HH^q_x(X,\Lambda(d))$ pour $q>2d$.

\begin{lemme2}
Soit $d\geq 3$.
On suppose le résultat du théorème~\ref{17-theoreme-classe-degre-maximal}
connu jusqu'en dimension $d-1$. Alors, pour tout $\ZZn$-schéma noethérien
excellent $X$ local strictement hensélien normal de dimension $d$, de point
fermé $x$, le morphisme $\Lambda\to \HH^{2d}_x(X,\Lambda(d))$
induit par $[x]$ est surjectif.
\end{lemme2}

On utilise la suite spectrale de coniveau calculant la cohomologie de
l'ouvert $U=X-x$ à coefficients dans $\Lambda(d)$ :
\[E_1^{pq}=\bigoplus_{y\in U^p} \HH^{p+q}_y(X_{(y)},\Lambda(d))\Longrightarrow
\HH^{p+q}(U,\Lambda(d))\;\text{,}\]
où $U^p$ désigne l'ensemble des points de codimension $p$ de $U$. La
topologie étale étant plus fine que la topologie de Nisnevich, dans
l'expression du terme $E_1$, on peut remplacer les localisés $X_{(y)}$ par
leurs hensélisés $X_{(y)}^\hensel$ pour tout $y\in U^p$ avec $0\leq p\leq
d-1$. On obtient ainsi un isomorphisme
canonique :
\[\RG_y(X_{(y)},\Lambda(d))\simeq
\RG(y,\SheafRG_y(\Lambda(p))(d-p))\;\text{.}\]
La structure de $\tau_{\geq 2p}\SheafRG_y(\Lambda(p))$ est
connue par notre connaissance limitée du
théorème~\ref{17-theoreme-classe-degre-maximal} puisque $p\leq d-1$. D'après
la proposition~\ref{17-proposition-dimension-cohomologique-galois}, on a une
majoration de la $\ell$-dimension cohomologique de $y$ pour tout nombre
premier $\ell$ divisant $n$ : $\cd_\ell y\leq d-p$.
En appliquant la suite spectrale de composition des
foncteurs dérivés au calcul de $\RG_y(X_{(y)},\Lambda(d))$, on
obtient
d'une part que $\HH^i_y(X_{(y)},\Lambda(d))=0$ pour $i>d+p$, c'est-à-dire que
$\HH^{p+q}(X_{(y)},\Lambda(d))=0$ pour $q>d$, et d'autre part que 
$\HH^{p+d}(X_{(y)},\Lambda(d))\simeq \HH^{d-p}(y,\Lambda(d-p))$.

On a donc montré que $E^{p,q}_1=0$ pour $q\geq d+1$, et par ailleurs, il est
évident que $E^{p,q}_1=0$ si $p\geq d$.
On en déduit aussitôt que l'on a
des isomorphismes canoniques $\coker(E^{d-2,d}_1\to E^{d-1,d}_1)\iso
\HH^{2d-1}(U,\Lambda(d))\iso \HH^{2d}_x(X,\Lambda(d))$. En particulier, le
morphisme canonique $E^{d-1,d}_1\to \HH^{2d}_x(X,\Lambda(d))$ est surjectif.
D'après le calcul ci-dessus, on connaît la structure de $E^{d-1,d}_1$ :
\[E^{d-1,d}_1 \simeq  \bigoplus_{y\in U^{d-1}} \HH^1(y,\Lambda(1))\;\text{.}\]
Pour tout $y\in U^{d-1}$, on a 
un isomorphisme évident $\HH^1(y,\Lambda(1))\simeq \Lambda$ (induit par une
spécialisation immédiate $\overline y\to x$ de points géométriques de $X$).
Il n'est pas difficile de vérifier que l'image du
morphisme canonique $\Lambda\simeq \HH^1(y,\Lambda(1))\subset
E_1^{d-1,d}\to \HH^{2d}_x(X,\Lambda(d))$ est 
le sous-groupe engendré par $[x]$. Par conséquent, l'image du morphisme
surjectif $E^{d-1,d}_1\to \HH^{2d}_x(X,\Lambda(d))$ est le sous-groupe
engendré par $[x]$, ce qui achève la démonstration du lemme.

\begin{lemme2}\label{17-lemme-recurrence-theoreme-classe-degre-maximal}
Soit $d\geq 3$.
On suppose le résultat du théorème~\ref{17-theoreme-classe-degre-maximal}
connu jusqu'en dimension $d-1$. Alors, pour tout $\ZZn$-schéma noethérien
excellent $X$ local strictement hensélien normal de dimension $d$, de point
fermé $x$, le morphisme $\Lambda\to \HH^{2d}_x(X,\Lambda(d))$
induit par $[x]$ est un isomorphisme, autrement dit, l'énoncé du
théorème~\ref{17-theoreme-classe-degre-maximal} vaut jusqu'en dimension $d$.
\end{lemme2}

En vertu du lemme précédent, il ne s'agit plus que de déterminer la
structure du $\Lambda$-module $\HH^{2d}_x(X,\Lambda(d))$. Compte tenu du
théorème de changement de base formel (cf.
\cite[corollary~6.6.4]{TUBULAR@Fujiwara}),
on peut supposer que $X$ est complet.
Les théorèmes de structure des anneaux locaux noethériens
complets (cf. \egazero{iv}{19.8.8})
montrent que $X$ est alors isomorphe à un
sous-schéma fermé d'un schéma régulier. Pour la fonction de dimension
$\delta$ sur $X$ telle que $\delta(\eta)=0$ où $\eta$ est le point
générique de $X$, compte tenu de ce qu'on sait déjà sur les complexes
dualisants potentiels (cf.
proposition~\ref{17-proposition-cdpotentiel-regulier}),
il existe un complexe dualisant potentiel
$K$ pour $(X,\delta)$. Nous allons utiliser le lemme général suivant :

\begin{lemme2}\label{17-lemme-structure-cohomologie-negative-cdp-normal}
Soit $X$ un $\ZZn$-schéma noethérien excellent normal intègre de point
générique $\eta$. On suppose $X$ muni de la fonction de dimension $\delta$
telle que $\delta(\eta)=0$. Soit $K$ un complexe dualisant potentiel pour
$(X,\delta)$. Alors, les faisceaux de cohomologie ${\mathscr H}^qK$  sont
nuls pour $q<0$ et l'épinglage en $\eta$ s'étend en un isomorphisme
${\mathscr H}^0K\simeq \Lambda$. Autrement dit, on a un isomorphisme
canonique $\Lambda\iso \tau_{\leq 0}K$.
\end{lemme2}

Pour obtenir le résultat pour $X$, il suffit de l'avoir pour ses hensélisés
stricts. On peut donc supposer que $X$ est strictement hensélien de
dimension $d$ et de point fermé $x$. On procède par récurrence sur $d$.
Si $d=0$, le résultat est évident. On suppose donc que $d\geq 1$ et que le
résultat est connu pour l'ouvert $U=X-x$. Notons $j\colon U\to X$ et
$i\colon x\to X$ les immersions évidentes. On a un triangle distingué
dans $\detaleplus X$ :
\[i_\star i^!K\to K\to \R j_\star j^\star K\to i_\star
i^! K[1]\;\text{.}\]
Grâce à l'épinglage en $x$, il vient que ${\mathscr H}^q i^!K=0$ pour
$q\leq 1$. L'hypothèse sur $U$ montre que $\Lambda\iso \tau_{\leq
0}j^\star K$. On en déduit que ${\mathscr H}^q K=0$ pour $q<0$ et que
${\mathscr H}^0K\simeq j_\star \Lambda$. Comme $X$ est normal, le morphisme
canonique $\Lambda\to j_\star \Lambda$ est un isomorphisme, ce qui achève
la démonstration du lemme.

\bigskip

Revenons à la démonstration du 
lemme~\ref{17-lemme-recurrence-theoreme-classe-degre-maximal}, on considère le
morphisme canonique $\Lambda\to K$ déduit du lemme ci-dessus. Soit
$\overline \eta$ un point géométrique au-dessus de $\eta$.
Choisissons une spécialisation $\overline\eta\to x$. On considère le
diagramme commutatif suivant, où les flèches verticales sont induites par
$\Lambda\to K$ et les flèches horizontales par les morphismes de transition
associés à la spécialisation $\overline\eta\to x$ :
\[
\xymatrix{
\ar[d]\HH^0_{\overline\eta}(\Lambda)\ar[r] & \HH^{2d}_x(X,\Lambda(d)) \ar[d]\\
\HH^0_{\overline\eta}(K)\ar[r] & \HH^{2d}_x(X,K(d))
}
\]
Le morphisme de gauche est évidemment un isomorphisme. Celui du bas aussi
puisque $K$ est un complexe dualisant potentiel. Le morphisme du haut est
donc injectif, mais on sait déjà qu'il est surjectif. Le $\Lambda$-module
$\HH^{2d}_x(X,\Lambda(d))$ est donc isomorphe à $\Lambda$, ce qui permet de
conclure.

%>>>2
%>>>1
\section{Compléments sur les complexes dualisants potentiels}%<<<1
\label{17-section-complements-cdp}
\subsection{Énoncés}%<<<2

Grâce aux résultats de la section~\ref{17-section-classe-de-cohomologie}, nous
allons pouvoir poursuivre l'étude de certaines propriétés des complexes
dualisants potentiels.
Le résultat principal de la
section~\ref{17-section-premieres-proprietes} était la construction d'un
complexe dualisant potentiel sur les schémas réguliers (cf.
proposition~\ref{17-proposition-cdputatif-regulier}). Les deux propositions
suivantes établissent des propriétés de stabilité des complexes dualisants
potentiels par rapport aux morphismes de type fini et aux morphismes
réguliers \ega{iv}{6.8.1}.

\begin{proposition2}\label{17-proposition-stabilite-cdp-morphisme-regulier}
Soit $f\colon Y\to X$ un morphisme régulier
entre $\ZZn$-schémas noethériens excellents. On suppose $X$ muni
d'une fonction de dimension $\delta_X$. On munit $Y$ de la fonction de
dimension $\delta_Y$ définie par l'égalité
$\delta_Y(y)=\delta_X(f(y))-\codim_{f^{-1}(f(y))}(y)$
pour tout $y\in Y$.

Si $K$ un complexe dualisant putatif sur $X$, alors $f^\star K$ est
naturellement muni d'une structure de complexe dualisant putatif sur $Y$,
et c'est un complexe dualisant potentiel si $K$ en est un.
\end{proposition2}

\begin{proposition2}\label{17-proposition-stabilite-cdp-morphisme-de-type-fini}
Soit $f\colon Y\to X$ un morphisme de type fini compactifiable
entre $\ZZn$-schémas noethériens excellents. On suppose $X$ muni
d'une fonction de dimension $\delta_X$. On munit $Y$ de la fonction de
dimension $\delta_Y$ définie par l'égalité
$\delta_Y(y)=\delta_X(f(y))+\degtr(y/f(y))$
pour tout $y\in Y$.

Si $K$ un complexe dualisant putatif sur $X$, alors $f^! K$ est
naturellement muni d'une structure de complexe dualisant putatif sur $Y$,
et c'est un complexe dualisant potentiel si $K$ en est un.
\end{proposition2}

\begin{remarque2}\label{17-remarque-cdp-compatibilite-composition}
Au cours des démonstrations, on observera que les constructions des
propositions~\ref{17-proposition-stabilite-cdp-morphisme-regulier} et
\ref{17-proposition-stabilite-cdp-morphisme-de-type-fini} sont compatibles à
la composition des morphismes : si $g\colon Z\to Y$ est un autre morphisme
du type envisagé et $K$ un complexe dualisant putatif sur $X$,
l'isomorphisme de transitivité $g^\star f^\star K\simeq (f\circ g)^\star K$
(resp. $g^! f^! K\simeq (f\circ g)^! K$) est compatible aux épinglages.
La proposition~\ref{17-proposition-stabilite-cdp-morphisme-regulier}
généralise la construction de la
proposition~\ref{17-proposition-stabilite-cdp-morphisme-etale} (cas où $f$ est
étale).
\end{remarque2}

Nous verrons plus bas que la
proposition~\ref{17-proposition-stabilite-cdp-morphisme-de-type-fini} résulte
facilement de la
proposition~\ref{17-proposition-stabilite-cdp-morphisme-regulier}. Nous allons
donc nous intéresser plus particulièrement aux morphismes réguliers. La
sous-section suivante sur le théorème de changement de base par un
morphisme régulier nous permettra de définir une structure de complexe
dualisant putatif sur $f^\star K$ dans la situation de la
proposition~\ref{17-proposition-stabilite-cdp-morphisme-regulier}. Dans le cas
d'un complexe dualisant potentiel, ce sont les résultats de la
section~\ref{17-section-classe-de-cohomologie} qui vont nous permettre de
vérifier que les épinglages sur $f^\star K$ sont compatibles aux
spécialisations.

%>>>2
\subsection{Changement de base par un morphisme régulier}%<<<2

Dans ce numéro, on étudie une généralisation du théorème de changement de
base par un morphisme lisse au cas du changement de base par un
morphisme régulier.
La démonstration utilise le théorème de
Popescu (cf.~\cite{NPD@Swan}).
On examine ensuite quelques conséquences de ce résultat.

\begin{proposition2}\label{17-proposition-changement-de-base-regulier}
Soit $f\colon Y\to X$ un morphisme régulier entre $\ZZn$-schémas
localement noethériens. Soit $g\colon X'\to X$ un
morphisme quasi-compact et quasi-séparé. On constitue le carré cartésien de
schémas :
\[
\xymatrix{
Y'\ar[r]^{f'}\ar[d]^{g'} & X'\ar[d]^{g} \\
Y\ar[r]^f & X
}
\]
Alors, pour tout $K\in \detaleplus {X'}$, le morphisme de changement de base
$f^\star \R g_\star K\to \R {g'}_\star {f'}^\star K$ est un
isomorphisme dans $\detaleplus Y$.
\end{proposition2}

Les arguments habituels montrent que l'on peut supposer que $X$, $Y$, puis
$X'$, et donc $Y'$, sont affines. D'après le théorème de Popescu, on peut
écrire $Y$ comme la limite projective d'un système projectif filtrant
$(Y_{\alpha}\vers {f_\alpha} X)_\alpha$ de $X$-schémas 
affines lisses (voir \ega{iv}{8}). Le théorème de changement de
base par un morphisme lisse \sga{4}{xvi}{1.2} montre que les morphismes de
changement de base $f_\alpha^\star \R g_\star K\to 
\R {g'_\alpha}_\star {f'}_\alpha^\star K$ sont des isomorphismes.
Ce résultat pour les morphismes
$f_\alpha$ s'étend au morphisme $f$ grâce à la technique de passage à la
limite de \sga{4}{vii}{5.8}, ce qui permet de conclure.

\begin{proposition2}\label{17-proposition-rhom-morphismes-reguliers}
Soit $f\colon Y\to X$ un morphisme régulier
entre $\ZZn$-schémas noethériens.
Soit $K\in\detalebc X$. Soit $L\in \detaleplus X$. Alors, le morphisme évident est
un isomorphisme dans $\detaleplus Y$ :
\[f^\star\R\SheafHom(K,L)\iso \R\SheafHom(f^\star K,f^\star
L)\;\text{.}\]
\end{proposition2}

Avant de la démontrer, signalons un corollaire de cette proposition :

\begin{corollaire2}\label{17-corollaire-cb-regulier-pour-cohomologie-a-supports}
Soit $f\colon Y\to X$ un morphisme régulier
entre $\ZZn$-schémas noethériens. Soit $i\colon Z\to X$ une 
immersion
fermée. Notons $i'\colon f^{-1}(Z)\to Y$ l'immersion fermée déduite par
changement de base, et $f'\colon f^{-1}(Z)\to Z$ la projection.
Alors, pour tout $K\in \detaleplus X$, le morphisme évident
${f'}^\star i^! K\to {i'}^!f^\star K$ est un isomorphisme dans $\detaleplus
{f^{-1}(Z)}$.
\end{corollaire2}

Pour démontrer la proposition~\ref{17-proposition-rhom-morphismes-reguliers},
commençons par constater que ce corollaire est un cas
particulier de la proposition~\ref{17-proposition-changement-de-base-regulier}
(l'appliquer avec pour $g$ l'immersion ouverte complémentaire de $i$).

\medskip

\begin{lemme2}\label{17-lemme-rhom-morphismes-reguliers-reduction-ferme}
Soit $f\colon Y\to X$ un morphisme régulier
entre $\ZZn$-schémas noethériens. Soit $i\colon Z\to X$ une
immersion fermée. On constitue le carré cartésien suivant :
\[\xymatrix{Z'\ar[r]^{i'}\ar[d]^{f'} & Y\ar[d]^f \\Z\ar[r]^i & X}\]
Soit $K\in \detalebc Z$. On suppose que pour tout $L\in\detaleplus Z$, le
morphisme canonique
${f'}^\star\R\SheafHom(K,L)\to\R\SheafHom({f'}^\star
K,{f'}^\star L)$ est un isomorphisme. Alors, pour tout $L\in\detaleplus X$, le
morphisme canonique $f^\star\R\SheafHom(i_\star K,L)\to
\R\SheafHom(f^\star i_\star K,f^\star L)$ est un isomorphisme.
\end{lemme2}

Ce lemme-ci résulte aussitôt du 
corollaire~\ref{17-corollaire-cb-regulier-pour-cohomologie-a-supports}.

\begin{lemme2}\label{17-lemme-rhom-morphismes-reguliers-reduction-ouvert}
Soit $f\colon Y\to X$ un morphisme régulier
entre $\ZZn$-schémas noethériens. Soit $j\colon U\to X$ une
immersion ouverte. On constitue le carré cartésien suivant :
\[\xymatrix{U'\ar[r]^{j'}\ar[d]^{f'} & Y\ar[d]^f \\U\ar[r]^j & X}\]
Soit $K\in \detalebc U$. On suppose que pour tout $L\in\detaleplus U$, le
morphisme canonique
${f'}^\star\R\SheafHom(K,L)\to\R\SheafHom({f'}^\star
K,{f'}^\star L)$ est un isomorphisme. Alors, pour tout $L\in\detaleplus X$, le
morphisme canonique $f^\star\R\SheafHom(j_! K,L)\to
\R\SheafHom(f^\star j_! K,f^\star L)$ est un isomorphisme.
\end{lemme2}

Ce lemme-là résulte de l'isomorphisme $f^\star \R j_\star \simeq
\R j'_\star f'^\star$ qui est un cas particulier de la 
proposition~\ref{17-proposition-changement-de-base-regulier}.

\bigskip

Démontrons la proposition~\ref{17-proposition-rhom-morphismes-reguliers}.
Les lemmes~\ref{17-lemme-rhom-morphismes-reguliers-reduction-ferme} et 
\ref{17-lemme-rhom-morphismes-reguliers-reduction-ouvert} permettent de
supposer que $K$ est un faisceau localement constant constructible de
$\Lambda$-modules sur $X$. En utilisant que le résultat est
trivial si $f$ est étale, on se ramène finalement au cas facile
où $K$ est un faisceau constant.

%>>>2
\subsection{Démonstration de la proposition~\ref{17-proposition-stabilite-cdp-morphisme-regulier}}%<<<2

\subsubsection{Structure de complexe dualisant putatif sur \texorpdfstring{$f^\star K$}{f*K}}%<<<3

\begin{lemme3}\label{17-lemme-construction-cdputatif-f-star}
Dans la situation de la
proposition~\ref{17-proposition-stabilite-cdp-morphisme-regulier}, si $K$ est
un complexe dualisant putatif sur $X$,
on peut munir $f^\star K$ d'une structure de
complexe dualisant putatif sur $Y$.
\end{lemme3}

Soit $K$ un complexe dualisant putatif sur $X$.
Soit $x\in X$. Introduisons le localisé $X'=X_{(x)}$ de $X$ en $x$ et
$Y'=Y\times_X X'$. Notons
$i\colon f^{-1}(x)\to Y'$ l'immersion de la fibre au-dessus de $x$,
$g\colon f^{-1}(x)\to x$ la projection et $j\colon Y'\to Y$ le morphisme
canonique.
D'après le
corollaire~\ref{17-corollaire-cb-regulier-pour-cohomologie-a-supports},
on a un isomorphisme canonique $g^\star
\SheafRG_x(K)\simeq i^!j^\star f^\star K$. Définir des épinglages
sur $f^\star K$ en les points de $f^{-1}(x)$ revient à définir une
structure de complexe dualisant putatif sur $i^!j^\star f^\star K$ pour
la fonction de dimension ${\delta_Y}_{|f^{-1}(x)}$. On vient de voir que
cet objet s'identifie à $g^\star \SheafRG_x(K)$ qui s'identifie
lui-même à $\Lambda(\delta_X(x))[2\delta_X(x)]$ en vertu de l'épinglage
donné de $K$ en $x$. La fibre $f^{-1}(x)$ étant régulière, le faisceau
constant $\Lambda$ est naturellement muni d'une structure de complexe
dualisant potentiel pour la fonction de dimension $-\codim$ (cf.
proposition~\ref{17-proposition-cdputatif-regulier}). Vu la définition de
$\delta_Y$, on en déduit aussitôt une structure de complexe dualisant
potentiel sur $\Lambda(\delta_X(x))[2\delta_X(x)]\in\detaleplus {f^{-1}(x)}$
pour la fonction de dimension ${\delta_Y}_{|f^{-1}(x)}$.
On a ainsi obtenu des
épinglages pour $f^\star K$ en tous les points de
$f^{-1}(x)$. Grâce à cette construction fibre à fibre, on a défini
une structure de complexe dualisant putatif sur $f^\star K$.

\begin{lemme3}\label{17-lemme-compatibilite-cdp-f-star-schemas-reguliers}
Soit $f\colon Y\to X$ un morphisme régulier entre $\ZZn$-schémas
noethériens excellents réguliers. On munit $X$ et $Y$ des fonctions de
dimension $\delta_X=-\codim$ et $\delta_Y=-\codim$. On munit les faisceaux
constants $\Lambda$ sur $X$ et $Y$ des structures de complexes dualisants
putatifs définies dans la proposition~\ref{17-proposition-cdputatif-regulier} ;
le lemme~\ref{17-lemme-construction-cdputatif-f-star} munit $f^\star\Lambda$
d'une structure de complexe dualisant putatif sur $Y$ pour la fonction de
dimension $\delta_Y$. Alors, l'isomorphisme canonique $\Lambda\simeq
f^\star \Lambda$ sur $Y$ est compatible aux épinglages.
\end{lemme3}

Vérifions la compatibilité de l'isomorphisme canonique $\Lambda\simeq
f^\star\Lambda$ aux épinglages en un point $y$ de $Y$. Quitte à
localiser, on peut supposer que $X$ est local de point fermé $x=f(y)$ et
que $y$ est un point fermé dans la fibre $F=f^{-1}(y)$. Notons $c=\codim_X
x$ et $d=\codim_F y$. L'épinglage de $f^\star\Lambda$ en $y$ est donné par
le produit des classes $f^\star (\Cl{x\to X})\in \HH^{2c}_F(Y,\Lambda(c))$ et
$\Cl{y\to F}\in \HH^{2d}_y(F,\Lambda(d))$, ce produit trouvant demeure dans
$\HH^{2c+2d}_y(Y,\Lambda(c+d))$. La compatibilité des classes de Gysin au
changement de base implique que $f^\star(\Cl{x\to X})$ est la classe de
Gysin $\Cl{F\to Y}\in \HH^{2c}_F(Y,\Lambda(c))$. En utilisant la
compatibilité des classes de Gysin à la composition, le produit considéré
plus haut est $\Cl{y\to Y}$, qui est précisément la classe qui définit
l'épinglage de $\Lambda$ en $y$.

\begin{remarque3}
Supposons que l'on dispose de deux morphismes réguliers composables $Z\vers
g Y\vers f X$ entre  $\ZZn$-schémas noethériens excellents. Si $K$ est un
complexe dualisant putatif sur $X$, alors appliquer la construction du
lemme~\ref{17-lemme-construction-cdputatif-f-star}
à $f$, puis à $g$, revient à l'appliquer directement à $f\circ
g$, autrement dit l'isomorphisme évident $g^\star
f^\star K\simeq (f\circ g)^\star K$ dans $\detaleplus Z$ est
compatible aux épinglages. En effet, les constructions s'effectuant fibre à
fibre, on peut supposer que $X$ est le spectre d'un corps et, quitte à
modifier la fonction de dimension, que $K=\Lambda$ ;
en particulier, $X$, $Y$ et $Z$ sont réguliers, les faisceaux
constants $\Lambda$ sur $X$, $Y$ et $Z$ sont naturellement munis de
structures de complexes dualisants putatifs (et même potentiels,
cf.~proposition~\ref{17-proposition-cdpotentiel-regulier}) pour les
fonctions de dimension considérées, la compatibilité requise est
obtenue en appliquant le
lemme~\ref{17-lemme-compatibilite-cdp-f-star-schemas-reguliers} aux
morphismes $f$, $g$ et $f\circ g$.
\end{remarque3}

%>>>3
\subsubsection{Compatibilité aux spécialisations}%<<<3

Supposons maintenant que $K$ soit un complexe dualisant potentiel.
Par définition des complexes dualisants potentiels, 
il s'agit maintenant de montrer que les épinglages sur $f^\star K$ sont
compatibles aux morphismes de transition associés aux spécialisations
immédiates. D'après le
théorème~\ref{17-theoreme-transition-toute-codimension}, on peut donner un
sens à cette compatiblité pour toute spécialisation
$\overline{y'}\to\overline y$ de points géométriques de $Y$, quelle que
soit sa codimension.

\begin{lemme3}\label{17-lemme-stabilite-cdp-morphisme-regulier}
Plaçons nous dans la situation de la
proposition~\ref{17-proposition-stabilite-cdp-morphisme-regulier}. Soit $K$ 
un complexe dualisant potentiel sur $X$. Soit ${\overline {y'}\to\overline y}$
une spécialisation de points géométriques de $Y$ au-dessus d'une
spécialisation $\overline {x'}\to\overline x$ de points géométriques de
$X$.
Si $\overline{x'}\to\overline x$ est de codimension $0$ ou $1$, alors les
épinglages sur $f^\star K$ sont compatibles à la spécialisation $\overline
{y'}\to\overline y$.
\end{lemme3}

Étant entendu que la construction de la structure de complexe dualisant
putatif sur $f^\star K$ a été réalisée fibre à fibre, on peut 
supposer que $X$ est local strictement hensélien intègre, de point
générique $x'$ et de point fermé $x$. On peut supposer que
$\delta_X(x')=0$.
Le schéma $X$ étant de dimension
$\leq 1$, si on note $n\colon \tilde X\to X$ la normalisation de $X$, le schéma
$\tilde X$ est régulier. On constitue le carré cartésien suivant :
\[
\xymatrix{\tilde Y\ar[r]^{n'}\ar[d]^{\tilde f} & Y\ar[d]^f \\
\tilde{X}\ar[r]^n & X
}
\]
Les morphismes $n$ et $n'$ sont finis surjectifs radiciels. Pour eux, on
dispose des constructions $n^!$ et ${n'}^!$ sur les complexes dualisants
putatif (cf. proposition~\ref{17-proposition-cdputatif-quasi-fini}).
Pour les morphismes
réguliers $f$ et $\tilde f$, on a les
constructions $f^\star$ et ${\tilde f}^\star$. On peut donc considérer les
complexes dualisants putatifs ${\tilde f}^\star n^!K$ et ${n'}^!f^\star K$.
D'après le lemme~\ref{17-lemme-compatibilite-f-star-homeo-universel} à suivre,
on a un isomorphisme canonique de complexes dualisant putatifs 
${\tilde f}^\star n^!K\simeq {n'}^!f^\star K$. En outre, la
proposition~\ref{17-proposition-cdpotentiel-morphisme-fini} montre que pour
montrer que $f^\star K$ est un complexe dualisant potentiel, il suffit de
montrer que ${n'}^!f^\star K$ en est un. Comme il s'identifie
${\tilde f}^\star n^!K$ et que la 
proposition~\ref{17-proposition-cdpotentiel-morphisme-fini} nous dit aussi que 
$n^!K$ est un complexe dualisant potentiel, on peut finalement remplacer
$f$ par $\tilde f$ et supposer que $X$ est
régulier de dimension $\leq 1$.
On peut alors supposer que $K=\Lambda$,
épinglé comme il convient de le faire. Le complexe $f^\star \Lambda$
s'identifie à $\Lambda$ dont on sait qu'il peut-être muni d'une structure
de complexe dualisant potentiel
(cf. proposition~\ref{17-proposition-cdpotentiel-regulier}). Il s'agit de
montrer que les épinglages sur $f^\star \Lambda$ et sur $\Lambda$ sont
compatibles : c'est le sens du
lemme~\ref{17-lemme-compatibilite-cdp-f-star-schemas-reguliers}.

\begin{lemme3}\label{17-lemme-compatibilite-f-star-homeo-universel}
Plaçons nous dans la situation de la
proposition~\ref{17-proposition-stabilite-cdp-morphisme-regulier}. Soit
$g\colon X'\to X$ un morphisme fini surjectif radiciel. On constitue le
carré cartésien :
\[
\xymatrix{
Y'\ar[r]^{g'}\ar[d]^{f'} & Y\ar[d]^f \\
X'\ar[r]^g & X
}
\]
Soit $K$ un complexe dualisant putatif sur $X$. Alors, l'isomorphisme
évident ${f'}^\star g^! K\simeq {g'}^!f^\star K$ dans $\detaleplus {Y'}$ est
compatible aux épinglages.
\end{lemme3}

Les constructions envisagées se réalisant fibre à fibre, on peut supposer
que $X$ et $X'$ sont les spectres de corps notés
respectivement $E$ et $E'$. Les morphismes $f$ et $f'$ étant plats et à
fibres géométriquement régulières, les schémas $Y$ et $Y'$ sont réguliers.
Dans cette situation, compte tenu de la 
proposition~\ref{17-proposition-cdpotentiel-morphisme-fini} et de 
la construction du lemme~\ref{17-lemme-construction-cdputatif-f-star},
il est manifeste que les complexes dualisants putatifs 
${f'}^\star g^! K$ et ${g'}^!f^\star K$ sont des complexes dualisant
potentiels. Pour montrer que l'isomorphisme évident 
${f'}^\star g^! K\simeq {g'}^!f^\star K$ est compatible aux épinglages, il
suffit donc de le faire aux points maximaux de $Y'$.
Bref, on peut supposer que
$Y$ est lui aussi le spectre d'un corps $F$. Comme $Y'$ est régulier et
homéomorphe à $Y$, le schéma $Y'$ est à son tour
le spectre d'un corps $F'$. Compte tenu
de la construction de la proposition~\ref{17-proposition-cdputatif-quasi-fini}
faisant intervenir les morphismes de Gysin associés aux morphismes
d'intersection complète $Y'\to Y$ et $X'\to X$,
la comparaison des deux structures de complexes dualisants putatifs
envisagées sur $Y'$ se ramène à l'égalité des degrés $[F':F]=[E':E]$,
qui résulte aussitôt de la définition de $Y'$ : $F'=E'\otimes_E F$.

\bigskip

Pour finir la démonstration de la
proposition~\ref{17-proposition-stabilite-cdp-morphisme-regulier}, il nous
reste à établir des compatibilités entre les épinglages sur $f^\star K$ et
les spécialisations de points géométriques de $Y$. Les lemmes suivants sur
les spécialisations nous seront utiles.

\begin{lemme3}\label{17-lemme-morphismes-reguliers-connexite}
Soit $f\colon Y\to X$ un morphisme régulier entre
schémas noethériens. On suppose que $f$ est un
morphisme local entre schémas locaux strictements henséliens. Soit
$\overline x\to X$ un point géométrique de $X$. Alors, $f^{-1}(\overline
x)=Y_{\overline x}=Y\times_X\overline x$ est un schéma intègre.
\end{lemme3}

La fibre $Y_x$ est géométriquement régulière et géométriquement connexe
d'après la proposition~\ref{17-proposition-changement-de-base-regulier}
appliquée à un faisceau constant. Le schéma $Y_{\overline x}$ est donc
intègre.

\begin{lemme3}\label{17-lemme-morphismes-reguliers-specialisations-generiques}
Soit $f\colon Y\to X$ un morphisme régulier
entre schémas noethériens excellents. Soit $\overline
{x'}\to\overline x$ une spécialisation de points géométriques de $X$. Soit
$\overline{\eta}$ un point géométrique de $X$ au-dessus de $\overline x$
tel que $\eta$ soit de codimension $0$ dans sa fibre pour $f$. Alors, à un
isomorphisme non nécessairement unique près, il existe une unique
spécialisation $\overline{\eta'}\to\overline\eta$ au-dessus de
$\overline{x'}\to\overline{x}$ telle que $\eta'$ soit de codimension $0$
dans sa fibre. La codimension de $\overline{\eta'}\to\overline\eta$
est la même que celle de $\overline {x'}\to\overline x$.
\end{lemme3}

On peut supposer que $f$ est un morphisme local entre schémas locaux
strictement henséliens $Y$ et $X$ de points fermés respectifs
$\overline\eta$ et $\overline x$. Établir le lemme dans ce cas précis
revient à montrer qu'à isomorphisme près, la fibre géométrique
$Y_{\overline x}$ n'a qu'un seul point géométrique au-dessus d'un point
maximal de $Y_{\overline x}$, ce qui résulte du 
lemme~\ref{17-lemme-morphismes-reguliers-connexite}.
Les formules liant les fonctions de dimension sur $X$ et $Y$ montrent que
la codimension de la spécialisation de 
$\overline{\eta'}\to\overline\eta$ est la même que celle de 
$\overline{x'}\to\overline{x}$.

\begin{lemme3}\label{17-lemme-morphismes-reguliers-relevement-specialisations}
Soit $f\colon Y\to X$ un morphisme régulier
entre schémas noethériens excellents. Soit
$\overline{\eta'}\to\overline\eta$ une spécialisation de points
géométriques de $Y$ au-dessus d'une spécialisation
$\overline{x'}\to\overline{x}$. On suppose que $\eta$ et $\eta'$ sont de
codimension $0$ dans leur fibre pour $f$. Soit
$\overline{x_0}\to\dots\to\overline{x_n}$ une factorisation de la
spécialisation $\overline{x'}\to\overline{x}$ en une suite de
spécialisations immédiates. Alors, on peut décomposer
$\overline{\eta'}\to\overline\eta$ en une suite de spécialisations
immédiates $\overline{\eta_0}\to\dots\to\overline{\eta_n}$ au-dessus de 
$\overline{x_0}\to\dots\to\overline{x_n}$, chacun des points $\eta_i$ étant
de codimension $0$ dans sa fibre pour $f$.
\end{lemme3}

Ceci résulte aussitôt du
lemme~\ref{17-lemme-morphismes-reguliers-specialisations-generiques},
propriété d'unicité comprise.

\begin{lemme3}\label{17-lemme-morphismes-reguliers-generisation-specialisations}
Soit $f\colon Y\to X$ un morphisme régulier
entre schémas noethériens excellents. Soit
$\overline{y'}\to\overline{y}$ une spécialisation de points géométriques de
$Y$ au-dessus de $\overline{x'}\to\overline x$.
Alors, il existe un diagramme commutatif de spécialisations de points
géométriques de $Y$ :
\[
\xymatrix{
\overline{\eta'}\ar[r]\ar[d] & \overline\eta\ar[d] \\
\overline{y'}\ar[r] & \overline y
}
\]
où les points géométriques $\overline{\eta'}$ et $\overline{\eta}$ sont
respectivement au-dessus de $\overline{x'}$ et de $\overline x$, et où
$\eta'$ et $\eta$ sont de codimension $0$ dans leur fibre pour $f$.
\end{lemme3}

On peut supposer que $X$ et $Y$ sont locaux strictement henséliens de
points fermés respectifs $\overline y$ et $\overline x$. On choisit une
spécialisation $\overline\eta\to \overline y$ de points géométriques de la
fibre $Y_{\overline x}$ avec $\eta$ de codimension $0$ dans cette fibre.
D'après le
lemme~\ref{17-lemme-morphismes-reguliers-specialisations-generiques}, il
existe une spécialisation $\overline {\eta'}\to\overline \eta$ de points
géométriques de $Y$ au-dessus de $\overline{x'}\to\overline x$, avec
$\eta'$ de codimension $0$ dans sa fibre pour $f$. Par ailleurs, la fibre
géométrique
$Y_{\overline{x'}}$ étant intègre, on peut en choisir un point géométrique
$\overline{\eta''}$ au-dessus du point générique.
Le point géométrique $\overline{y'}$ de $Y$ étant au-dessus de
$\overline{x'}$, il définit un point géométrique de $Y_{\overline{x'}}$ ;
on dispose donc d'une spécialisation $\overline{\eta''}\to\overline{y'}$ de
points géométriques au-dessus de $\overline{x'}$.
La fibre géométrique $Y_{\overline{x'}}$ étant intègre, il existe un
isomorphisme $\overline{\eta'}\simeq\overline{\eta''}$, ce qui donne le
diagramme commutatif souhaité.

%<<<3
\subsubsection{Fin de la démonstration}%<<<3

Finissons la démonstration de la
proposition~\ref{17-proposition-stabilite-cdp-morphisme-regulier}, il s'agit
de montrer que pour tout complexe dualisant potentiel $K$ sur $X$, les
épinglages sur le complexe dualisant putatif $f^\star K$ sont compatibles
aux spécialisations de points géométriques de $Y$. Soit
$\overline{y'}\to\overline y$ une telle spécialisation. Le
lemme~\ref{17-lemme-morphismes-reguliers-generisation-specialisations}
s'applique et on
obtient un diagramme commutatif de spécialisations comme ci-dessus. Pour
montrer la compatibilité pour la spécialisation $\overline{y'}\to\overline
y$, il suffit de l'obtenir pour les trois autres spécialisations qui
interviennent dans le carré commutatif. Pour les spécialisations
$\overline{\eta'}\to\overline{y'}$ et $\overline{\eta}\to\overline{y}$,
cela résulte aussitôt des faits observés dans 
la construction même du lemme~\ref{17-lemme-construction-cdputatif-f-star}, à
savoir que chaque fibre de $f$ est munie d'un complexe dualisant potentiel.
On peut appliquer le
lemme~\ref{17-lemme-morphismes-reguliers-relevement-specialisations}
pour obtenir une
factorisation de $\overline{\eta'}\to\overline\eta$ en une suite de
spécialisations immédiates $\overline{\eta_0}\to\dots\to\overline{\eta_n}$
au-dessus d'une composition de spécialisations immédiates
$\overline{x_0}\to\dots\to\overline{\eta_n}$ de points géométriques de $X$.
On peut appliquer le lemme~\ref{17-lemme-stabilite-cdp-morphisme-regulier} aux
spécialisations $\overline{\eta_i}\to\overline{\eta_{i+1}}$ pour obtenir la
compatibilité souhaitée pour $\overline{\eta'}\to\overline\eta$, ce qui
achève la démonstration de la
proposition~\ref{17-proposition-stabilite-cdp-morphisme-regulier}.

%>>>3
%>>>2
\subsection{Démonstration de la proposition~\ref{17-proposition-stabilite-cdp-morphisme-de-type-fini}}%<<<2
\label{17-subsection-morphismes-de-type-fini}

\begin{lemme2}\label{17-lemme-construction-cdputatif-f-bang}
Dans la situation de la
proposition~\ref{17-proposition-stabilite-cdp-morphisme-de-type-fini},
si $K$ est un complexe dualisant putatif sur $X$,
on peut munir $f^!K$ d'une structure de
complexe dualisant putatif sur $Y$.
\end{lemme2}

Définir la structure de complexe dualisant putatif sur $f^!K$ peut se faire
fibre à fibre. On suppose donc que $X=x$ est le spectre d'un corps.
Soit $y\in Y$. Soit $U$ un ouvert non vide de régularité de l'adhérence
(réduite) $\overline{\{y\}}$. On note $j\colon U\to Y$ l'immersion de $U$
dans $Y$ et $\pi\colon U\to x$ la projection. On a des isomorphisme
canoniques induits par le morphisme de Gysin $\Cl \pi$ et l'épinglage de
$K$ en $x$ :
\[\Lambda(\delta_Y(y))[2\delta_Y(y)]\iso
\pi^!\Lambda(\delta_X(x))[2\delta_X(x)]\giso \pi^!\SheafRG_x(K)
\simeq j^!f^!K\;\text{.}\]
En passant au point générique de $U$, on obtient l'isomorphisme voulu :
$\SheafRG_y(f^!K)\simeq \Lambda(\delta_Y(y))[2\delta_Y(y)]$. Il ne
dépend évidemment pas de l'ouvert $U$, ce qui permet de définir l'épinglage
souhaité de $f^!K$ en $y$.

\medskip

On peut vérifier que la construction de cette structure de complexe
dualisant putatif sur $f^!K$ est compatible à la composition des morphismes
de type fini (pour le sens de cette affirmation, cf.
remarque~\ref{17-remarque-cdp-compatibilite-composition}). Par ailleurs, il
est évident que dans le cas où $f$ est quasi-fini, les épinglages définis
ici sont les mêmes que ceux de la
proposition~\ref{17-proposition-cdputatif-quasi-fini}. Enfin, dans le cas où
$f$ est lisse de dimension relative $d$,
\emph{via} l'isomorphisme canonique $f^\star K\simeq
f^!K(d)[2d]$, les épinglages sont compatibles avec ceux de la
proposition~\ref{17-proposition-stabilite-cdp-morphisme-regulier}, compte tenu
du décalage entre les deux fonctions de dimension envisagées sur $Y$.

\bigskip

Montrons la
proposition~\ref{17-proposition-stabilite-cdp-morphisme-de-type-fini}. La
question étant de nature locale, on peut supposer que $f\colon X\to Y$ est
un morphisme de type fini entre schémas affines. Il existe donc une
factorisation $f=p\circ i$ avec $i$ une immersion fermée et $p$ un
morphisme lisse. Soit $K$ un complexe dualisant potentiel sur $X$.
Les remarques précédentes montrent que $p^!K$,
puis $i^!p^!K$ sont des complexes dualisants
potentiels et que ce dernier s'identifie à $f^!K$.
Par conséquent, $f^!K$ est un
complexe dualisant potentiel, ce qui achève la démonstration.

%>>>2
%>>>1
\section{Existence et unicité des complexes dualisants potentiels}%<<<1
\label{17-section-existence-unicite-cdp}
\subsection{Énoncé du théorème}%<<<2

L'objectif de cette section est d'établir le théorème suivant :

\begin{theoreme2}\label{17-theoreme-existence-unicite-cdp}
Soit $X$ un $\ZZn$-schéma excellent muni d'une fonction de dimension
$\delta$. Alors $(X,\delta)$ admet un complexe dualisant potentiel $K_X$,
unique à isomorphisme unique près, et le morphisme évident $
\Lambda\to \tau_{\leq 0}\R\SheafHom(K_X,K_X)$
est un isomorphisme dans $\detale X$. De plus,
$K_X\in\Perv^{-2\delta}(X,\Lambda)$ (cf.
sous-section~\ref{17-subsection-faisceaux-pervers}).
\end{theoreme2}

%>>>2
\subsection{Préliminaires sur les faisceaux pervers}%<<<2
\label{17-subsection-faisceaux-pervers}

Si $X$ est un schéma noethérien et $p\colon
X\to\mathbf{Z}\cup \{+\infty\}$ une fonction de perversité
(c'est-à-dire que pour toute spécialisation $\overline y\to \overline x$ de
points géométriques de $X$, on a $p(x)\geq p(y)$),
Gabber a défini dans \cite{T@Gabber} une $t$-structure
$(\detaletronque X {\leq p},\detaletronque X {\geq p})$ sur $\detaleplus X$ de sorte que
pour tout $K\in\detaleplus X$, on ait :
\[K\in\detaletronque X {\leq p}\Longleftrightarrow \forall x\in X, K_{|x}\in
\detaletronque x {\leq p(x)}\;\text{,}\]
\[K\in\detaletronque X {\geq p}\Longleftrightarrow \forall x\in X,
\SheafRG_x(K)\in
\detaletronque x {\geq p(x)}\;\text{,}\]
où l'on a muni chacune des catégories $\detaleplus x$ de sa $t$-structure
canonique.

On note $\Perv^p(X,\Lambda)\subset \detaleplus X$ le cœur de cette
$t$-structure, les foncteurs de troncatures étant notés $\tau_{\leq p}$ et
$\tau_{\geq p}$.

\begin{definition2}
Soit $f\colon Y\to X$ un morphisme de schémas. On dit que $f$ est une
pseudo-immersion ouverte si $f$ induit un homéomorphisme sur son image et
que le morphisme induit $f^{-1}\OO_X\to \OO_Y$ est un isomorphisme.
\end{definition2}

Cette classe de morphismes est stable par composition ; elle contient les
immersions ouvertes et les localisations.

\begin{proposition2}\label{17-proposition-faisceaux-pervers-immersions-ouvertes}
Soit $p\colon X\to\mathbf{Z}\cup \{+\infty\}$ une fonction de perversité.
Soit $f\colon Y\to X$ une pseudo-immersion ouverte.
\begin{enumerate}
\item[(a)] La fonction $p\circ f\colon Y\to \mathbf{Z}\cup\{+\infty\}$
est une fonction de perversité (encore notée $p$)
et pour les $t$-structures définies par $p$, $f^\star\colon\detaleplus X\to
\detaleplus Y$ est $t$-exact et $\R f_\star\colon \detaleplus Y\to
\detaleplus X$
$t$-exact à gauche ;
\item[(b)] Le foncteur $f^\star$ induit un foncteur exact $f^\star\colon
\Perv^p(X,\Lambda)\to\Perv^p(Y,\Lambda)$ qui admet un adjoint à droite
$f^p_\star\colon \Perv^p(Y,\Lambda)\to \Perv^p(X,\Lambda)$ défini par la
formule $f_\star^p K=\tau_{\leq p}\R f_\star K$ ;
\item[(c)] Le morphisme d'adjonction $f^\star f^p_\star\to
\Id_{\Perv^p(Y,\Lambda)}$ est un
isomorphisme et le foncteur $f_\star^p\colon \Perv^p(Y,\Lambda)\to
\Perv^p(X,\Lambda)$ est pleinement fidèle.
\item[(d)] Si $g\colon Z\to Y$ est une pseudo-immersion ouverte composable
avec $f$, on dispose d'un isomorphisme de transitivité $f_\star^p\circ
g_\star^p\simeq (f\circ g)_\star^p$.
\end{enumerate}
\end{proposition2}

La $t$-exactitude de $f^\star$ est triviale. En particulier, $f^\star$ est
$t$-exact à droite ; par adjonction, $\R f_\star$ est $t$-exact à
gauche. (b) résulte aussitôt de (a). (c) en résulte aussi compte
tenu du fait que le morphisme d'adjonction $f^\star \R f_\star\to
\Id_{\Perv^p(Y,\Lambda)}$
est un isomorphisme. L'isomorphisme du (d) se déduit par adjonction de
l'isomorphisme de transitivité des foncteurs images inverses.

%>>>2
\subsection{Cas d'un schéma normal}%<<<2

\begin{proposition2}\label{17-proposition-theoreme-existence-unicite-cdp-cas-normal}
L'énoncé du théorème~\ref{17-theoreme-existence-unicite-cdp} est vrai si on
suppose de plus que le schéma $X$ est normal.
Plus précisément, soit $X$ un $\ZZn$-schéma
noethérien excellent irréductible normal
de point générique $\eta$,
muni d'une fonction de dimension $\delta$ telle que $\delta(\eta)=0$. On
note $j\colon \eta\to X$ l'inclusion du point générique et on pose
$T=j^\varphi_\star \Lambda$ où $\varphi\colon X\to\mathbf{N}$ est la
fonction de perversité définie par l'égalité :
\[\varphi(x)=\max(0,2\dim \OO_{X,x}-2)\]
pour tout $x\in X$.
\begin{enumerate}
\item[(a)] Le morphisme d'adjonction $\Lambda\to j_\star^\varphi
j^\star\Lambda$ définit un morphisme évident $\Lambda\to T $ tel que pour
tout point géométrique $\overline x\to X$, l'application induite
\[\HH^{-2\delta(x)}_{\overline x}(\Lambda)\to \HH^{-2\delta(x)}_{\overline x}(T)\]
soit un isomorphisme et que $\HH^q_{\overline x}(T)=0$ si $q\neq -2\delta(x)$.
Le théorème~\ref{17-theoreme-classe-degre-maximal} donne alors un épinglage
de $T$ en $x$. Avec ces épinglages, $T$ est un complexe dualisant
potentiel pour $(X,\delta)$.
\item[(b)] Si $K$ est un complexe dualisant potentiel pour $(X,\delta)$,
alors $K$ appartient à $\Perv^{\varphi}(X,\Lambda)$
(et à $\Perv^{-2\delta}(X,\Lambda)$)
et le morphisme $K\to T$ qui s'en déduit aussitôt est un isomorphisme
compatible aux épinglages.
\item[(c)] Si $K$ est un complexe dualisant potentiel sur $X$, alors le
morphisme évident $\Lambda\to\tau_{\leq 0}\R \SheafHom(K,K)$ est un
isomorphisme.
\end{enumerate}
\end{proposition2}

Établissons (a). Le point essentiel est de montrer que pour tout $x\in X$,
le morphisme évident
$\HH^{-2\delta(x)}_{\overline x}(\Lambda)\to \HH^{-2\delta(x)}_{\overline x}(T)$
est un isomorphisme et que $\HH^q_{\overline x}(T)$ est nul si $q\neq
-2\delta(x)$. Le fait que $T$ soit un complexe dualisant
potentiel résultera alors aussitôt du
théorème~\ref{17-theoreme-classe-degre-maximal}. Pour établir le résultat
voulu, on peut supposer que $X$ est local strictement hensélien de
point fermé $x$. Par récurrence sur $d=\dim X$, on peut supposer (a) connu
sur l'ouvert $U=X-x$. Si $d=0$, le résultat est évident. On suppose donc
$d\geq 1$, de sorte que $U$ contienne le point générique de $X$.
Posons $L=T_{|U}$. D'après la
proposition~\ref{17-proposition-faisceaux-pervers-immersions-ouvertes}, si on
note $g\colon U\to X$ l'immersion ouverte évidente, on a un isomorphisme
canonique $T=g_\star^\varphi L$. On en déduit aussitôt un isomorphisme
canonique $T=\tau_{\leq \varphi(x)}\R g_\star L$.

Si $d=1$, $\varphi(\eta)=0$, donc $T=g_\star\Lambda=\Lambda$. La
proposition~\ref{17-proposition-cdp-trait} permet de conclure que $T$ est bien
un complexe dualisant potentiel avec les épinglages envisagés ici. On
suppose donc que $d\geq 2$. Dans ce cas, on a $\varphi(x)=2d-2$ et
$T=\tau_{\leq 2d-2}\R g_\star L$.

D'après le lemme~\ref{17-lemme-structure-cohomologie-negative-cdp-normal}, la
structure des objets de cohomologie ${\mathscr H}^q L$ pour $q\leq 0$ est
connue. Notons $i\colon x\to X$ l'immersion du point fermé de $X$.
Utilisons le triangle distingué canonique :
\[i_\star i^! T\to T\to \R g_\star L\to i_\star i^!T[1]\;\text{.}\]
Il en résulte une suite exacte longue :
\[\dots \to \HH^{q-1}(U,L)\to \HH^q_x(X,T)\to
({\mathscr H}^qT)_x \to \HH^q(U,L)\to \dots\]
Par construction, $({\mathscr H}^q T)_x\iso \HH^q(U,L)$ si $q\leq 2d-2$ et 
$({\mathscr H}^q T)_x=0$ si $q\geq 2d-1$. Il en résulte que $\HH^q_x(X,T)=0$
si $q\leq 2d-1$ et que l'on a des isomorphismes $\HH^{q-1}(U,L)\iso
\HH^q_x(X,T)$ pour $q\geq 2d$.
Il vient aussi que ${\mathscr H}^q T=0$ si $q<0$ et que l'on a un
isomorphisme canonique $\Lambda\iso {\mathscr H}^0T$. Ainsi, on a un
morphisme canonique
$\Lambda \to T$ induisant un isomorphisme $\Lambda\iso
\tau_{\leq 0}T$.

\begin{lemme2}\label{17-lemme-technique-dimcohom-existence-unicite-cdp}
Soit $U$ le complémentaire du point fermé dans un schéma local strictement
hensélien excellent normal $X$ de dimension $d\geq 2$.
Soit $M\in\detaletronque U {\leq \varphi}$. On suppose qu'il existe un isomorphisme 
$\Lambda\iso \tau_{\leq 0}M$. Alors, pour tout $q\geq 2d$, le morphisme
évident est un isomorphisme :
\[\HH^{q-1}(U,\Lambda)\iso \HH^{q-1}(U,M)\;\text{.}\]
\end{lemme2}

On note $M^+$ un cône du morphisme évident $\Lambda\to M$. Il suffit de
montrer que $\HH^q(U,M^+)=0$ pour $q\geq 2d-2$. Les hypothèses
impliquent que les objets de cohomologie de $M^+$ sont nuls en dehors
de l'intervalle $[1,2d-4]$. On peut aussi observer que pour tout 
$1\leq i\leq d-2$, si $y\in U$ est tel que $({\mathscr
H}^{2i-1} M^+)_{\overline y}$ ou $({\mathscr H}^{2i} M^+)_{\overline y}$
soit non nul,
alors l'adhérence (dans $X$) de $y$ est de dimension $\leq d-i-1$.
D'après la proposition~\ref{17-proposition-dimension-cohomologique} et la
suite spectrale d'hypercohomologie, on obtient bien que $\HH^q(U,M^+)=0$ si
$q\geq 2d-2$.

\bigskip

On peut appliquer le
lemme~\ref{17-lemme-technique-dimcohom-existence-unicite-cdp}
avec $M=L$. Il résulte alors de ce qui précède et du
théorème~\ref{17-theoreme-classe-degre-maximal} que $\HH^q_x(X,T)=0$ si $q\neq
2d$, que l'on a un isomorphisme canonique 
$\HH^{2d}_x(X,\Lambda(d))\iso \HH^{2d}_x(X,T(d))$, et que l'on peut ainsi 
définir des épinglages sur $T$ qui en font un complexe dualisant potentiel
sur $X$.

\bigskip

Montrons maintenant (b). Soit $K$ un complexe dualisant potentiel sur $X$.
Il est évident que $K\in\detaletronque X {\geq -2\delta}\subset
\detaletronque X {\geq
\varphi}$. Pour montrer que $K\in \detaletronque X {\leq \varphi}\subset
\detaletronque X
{\leq -2\delta}$, on peut supposer que $X$ est local strictement hensélien
de dimension $d$.
On note $i\colon x\to X$ l'inclusion du point fermé
et $g\colon U\to X$ l'inclusion de l'ouvert complémentaire $X-x$. Le cas où
$d=0$ étant trivial et celui où $d=1$ ayant été traité dans la 
proposition~\ref{17-proposition-cdp-trait}, on peut supposer que $d\geq 2$.
Par récurrence sur $d$, on peut supposer (b) connu pour l'ouvert $U$. On
note $L=K_{|U}\in\detaletronque U {\leq\varphi}$.

Comme $K$ est un complexe dualisant potentiel, 
la structure des faisceaux de cohomologie ${\mathscr H}^qK$ pour $q\leq 0$
est connue (cf.
lemme~\ref{17-lemme-structure-cohomologie-negative-cdp-normal}). On peut 
donc appliquer le
lemme~\ref{17-lemme-technique-dimcohom-existence-unicite-cdp} avec $M=L$.
Ainsi, le morphisme évident est un isomorphisme $\HH^{2d-1}(U,\Lambda)\iso
\HH^{2d-1}(U,L)$, et $\HH^q(U,L)=0$ pour $q\geq 2d$.

De même que pour établir (a), on utilise le triangle distingué
canonique :
\[i_\star i^! T\to T\to \R g_\star L\to i_\star i^!T[1]\;\text{,}\]
et la suite exacte longue qui s'en déduit :
\[\dots \to \HH^{q-1}(U,L)\to \HH^q_x(X,K)\to
({\mathscr H}^qK)_x \to \HH^q(U,L)\to \dots\]
La structure de $\HH^q_x(X,K)$ étant connue pour tout $q\in\mathbf{Z}$,
il vient aussitôt que $({\mathscr H}^qK)_x=0$ pour $q\geq 2d+1$ et que
l'on a une suite exacte
\[0\to ({\mathscr H}^{2d-1}K)_x\to \HH^{2d-1}(U,L)\to \HH^{2d}_x(X,K)
\to ({\mathscr H}^{2d}K)_x \to 0\;\text{.}\]
Pour montrer que $K\in\detaletronque X {\leq \varphi}$, il reste donc à montrer que
le morphisme canonique $\HH^{2d-1}(U,L)\to \HH^{2d}_x(X,K)$ est un
isomorphisme. Choisissons une spécialisation $\overline{\eta}\to x$. On
peut considérer le diagramme commutatif suivant, où les flèches verticales
sont induites par le morphisme canonique $\Lambda\to K$ :
\[
\xymatrix{
\ar[d]^{\sim}
\HH^0_{\overline {\eta}}(X,\Lambda)\ar[r]^-{\transition^X_{\overline{\eta}\to x}} &
   \HH^{2d}_x(X,\Lambda(d))\ar[d] & \ar[l]_{\sim}
   \HH^{2d-1}(U,\Lambda(d))\ar[d]^{\sim} \\
\HH^0_{\overline {\eta}}(X,K)\ar[r]^-{\transition^X_{\overline{\eta}\to x}} &
   \HH^{2d}_x(X,K(d)) & \ar[l] \HH^{2d-1}(U,L(d))
}\]
Les flèches de la colonne de gauche vers celle du milieu sont les 
morphismes de transition introduits au
théorème~\ref{17-theoreme-transition-toute-codimension}. Ici, ce sont des
isomorphismes : pour la flèche du haut, cela résulte du
théorème~\ref{17-theoreme-classe-degre-maximal} et pour la flèche du bas, du
fait que $K$ soit un complexe dualisant potentiel. On en déduit que le
morphisme du milieu $\HH^{2d}_x(X,\Lambda(d))\to \HH^{2d}_x(X,K(d))$ est un
isomorphisme. Compte tenu des autres isomorphismes connus, il vient que le
morphisme évident $\HH^{2d-1}(U,L(d))\to \HH^{2d}_x(X,K(d))$ est un
isomorphisme, ce qui achève de montrer que $K\in \Perv^\varphi(X,\Lambda)$.

On a alors un morphisme d'adjonction $K\to j_\star^\varphi j^\star K=T$
dans $\Perv^\varphi(X,\Lambda)$. Pour montrer que c'est un isomorphisme, on
peut se placer dans la situation locale précédente, et faire une récurrence
sur la dimension pour pouvoir supposer que le morphisme induit $K_{|U}\to
T_{|U}$ est un isomorphisme. Choisissons une spécialisation
$\overline{\eta}\to x$. On en déduit un diagramme commutatif :
\[
\xymatrix{\ar[d]\HH^0_{\overline
\eta}(X,K)\ar[r]^-{\transition^X_{\overline{\eta}\to x}} & \HH^{2d}_x(X,K(d)) \ar[d]\\
\HH^0_{\overline
\eta}(X,T)\ar[r]^-{\transition^X_{\overline{\eta}\to x}} & \HH^{2d}_x(X,T(d))
}
\]
La flèche de gauche est évidemment un isomorphisme. Les flèches
horizontales aussi puisque $K$ et $T$ sont des complexes dualisants
potentiels. Il vient donc que le morphisme induit $i^!K\to i^!T$ est un
isomorphisme. Il en découle que $K\to T$ est un isomorphisme (compatible
aux épinglages au point générique, donc à tous les épinglages),
ce qui achève la démonstration de (b).

Montrons (c). Soit $K$ un complexe dualisant potentiel sur $X$. Pour
montrer que le morphisme canonique $\Lambda\to\tau_{\leq
0}\R\SheafHom(K,K)$ est un isomorphisme, quitte à remplacer $X$ par
des schémas connexes étales sur lui,
il suffit de montrer que $\Lambda\iso
\Hom_{\detaleplus X}(K,K)$ et que $\Hom_{\detaleplus X}(K,K[q])=0$ pour $q<0$.
L'annulation de $\Hom_{\detaleplus X}(K,K[q])$ pour $q<0$ résulte aussitôt du
fait que $K$ appartienne au cœur de la $t$-structure définie par $\varphi$.
D'après (b), on a un isomorphisme canonique $K=j_\star^\varphi \Lambda$.
L'isomorphisme $\Lambda\iso\Hom_{\detaleplus X}(K,K)$ résulte alors de ce
que $j_\star^\varphi\colon \Perv^\varphi(\eta,\Lambda)\to
\Perv^\varphi(X,\Lambda)$ soit pleinement fidèle
(cf. proposition~\ref{17-proposition-faisceaux-pervers-immersions-ouvertes}).

%>>>2
\subsection{Un résultat de recollement}%<<<2
\label{17-subsection-recollement-cdp}

La proposition suivante est un résultat de recollement qui nous permettra
de passer du cas normal au cas général :

\begin{proposition2}\label{17-proposition-existence-unicite-cdp-recollement}
Soit $X$ un $\ZZn$-schéma noethérien excellent muni d'une fonction de
dimension $\delta$. On suppose donné un carré cartésien
\[
\xymatrix{
Y'\ar[d]^{p'}\ar[r]^{i'} & X'\ar[d]^p \\
Y\ar[r]^i & X }
\]
où $i$ est une immersion fermée d'ouvert complémentaire $U$
et où $p$ est un morphisme fini induisant
un isomorphisme $p^{-1}(U)\iso U$. On note $q=p\circ i'$. On suppose que
l'énoncé du théorème~\ref{17-theoreme-existence-unicite-cdp} est connu pour
$X'$, $Y$ et $Y'$ (relativement aux fonctions de dimensions déduites de
celle sur $X$ par le procédé de la
proposition~\ref{17-proposition-cdputatif-quasi-fini}). Alors, cet énoncé est
également vrai pour $X$, et si on note $K_X$, $K_{X'}$,
$K_Y$ et $K_{Y'}$ les complexes dualisants potentiels de $X$, $X'$, $Y$ et
$Y'$ respectivement, on a un triangle distingué canonique dans $\detaleplus X$
:
\[q_\star K_{Y'}\to i_\star K_Y\oplus p_\star K_{X'}\to K_X\to q_\star
K_{Y'}[1]\;\text{.}\]
\end{proposition2}

Dans un premier temps, supposons que $X$ admette un complexe dualisant
potentiel $K_X$ et montrons que l'on peut définir un triangle distingué
canonique de la forme souhaitée. Pour cela, considérons la suite exacte
courte évidente de faisceaux sur $X$ :
\[0\to \Lambda\vers {(+,+)} i_\star\Lambda\oplus p_\star\Lambda\vers{(+,-)}
q_\star\Lambda\to 0\;\text{.}\]
En appliquant $\R\SheafHom(-,K_X)$ au triangle distingué
correspondant, on obtient un triangle distingué :
\[q_\star q^!K_X\vers {(+,-)}
i_\star i^! K_X\oplus p_\star p^!K_X\vers{(+,+)}
K_X\to q_\star q^!K_X[1]\;\text{,}\]
qui, compte tenu de la
proposition~\ref{17-proposition-stabilite-cdp-morphisme-de-type-fini} et de la
vertu d'unicité des complexes dualisants potentiels sur $X'$, $Y$ et $Y'$,
se récrit sous la forme :
\[q_\star K_{Y'}\vers {(+,-)} i_\star K_Y\oplus p_\star K_{X'}\vers{(+,+)}
K_X\to q_\star K_{Y'}[1]\;\text{.}\]

Revenant aux hypothèses de la proposition, nous allons montrer
qu'inversement, si on définit $K_X$ de façon à avoir un tel triangle
distingué (mais \emph{a priori} pas de façon canonique),
on obtient bien un complexe dualisant potentiel sur $X$. On
suppose donc le théorème~\ref{17-theoreme-existence-unicite-cdp} connu
seulement pour $X'$, $Y$ et $Y'$ et on note $K_{X'}$, $K_Y$ et $K_{Y'}$ les
complexes dualisants potentiels correspondants. La propriété d'unicité pour
les complexes dualisants potentiels sur $Y'$ donne des isomorphismes
canoniques $K_{Y'}\simeq {p'}^!K_Y$ et $K_{Y'}\simeq {i'}^!K_{X'}$. Par
adjonction, on en déduit des morphismes canoniques 
$p'_\star K_{Y'}\to K_Y$ et $i'_\star K_{Y'}\to K_{X'}$, puis en appliquant
respectivement $i_\star$ et $p_\star$, on obtient des morphismes canoniques
$q_\star K_{Y'}\to i_\star K_Y$ et $q_\star K_{Y'}\to p_\star
K_{X'}$. On peut considérer leur différence et constituer
un triangle distingué dans $\detaleplus X$ :
\[q_\star K_{Y'}\vers {(+,-)} i_\star K_Y\oplus p_\star K_{X'}\to
K_X\to q_\star K_{Y'}[1]\;\text{.}\]
On obtient ainsi un objet $K_X\in\detaleplus X$ et deux morphismes privilégiés
$i_\star K_Y\to K_X$ et $p_\star K_{X'}\to K_X$.

Notons $j\colon U\to X$ et $j'\colon U\to X'$ les immersions ouvertes
évidentes. En appliquant $j^\star$ au triangle distingué ci-dessus, on peut
commencer par observer que le morphisme évident ${j'}^\star K_{X'}\simeq
j^\star p_\star K_{X'}\to j^\star K_X$ est un isomorphisme. Par conséquent,
on peut munir $K_X$ d'épinglages en les points de l'ouvert $U$ de façon
compatible avec la structure de complexe dualisant potentiel obtenue sur
${j'}^\star K_{X'}$.

Considérons le morphisme canonique $q_\star K_{Y'}\to p_\star K_{X'}$.
J'affirme qu'il induit un isomorphisme après application du foncteur $i^!$.
En effet, on a des isomorphismes évidents de foncteurs $i^!q_\star
{i'}^!\simeq i^!i_\star p'_\star {i'}^\star\simeq p'_\star {i'}^\star\simeq
i^!p_\star$ et, compte tenu de l'isomorphisme canonique $K_{Y'}\simeq
{i'}^!K_{X'}$, on obtient bien que le morphisme évident $i^!q_\star
K_{Y'}\to i^!p_\star K_{X'}$ est un isomorphisme. En appliquant $i^!$ au
triangle distingué de définition de $K_X$, il vient alors que le morphisme
canonique $i_\star K_Y\to K_X$ induit un isomorphisme $K_Y\iso i^! K_X$
après application de $i^!$. Ceci permet de définir des épinglages pour
$K_X$ en tous les points de $Y$.

Finalement, on a obtenu une structure de complexe dualisant putatif sur
$K_X$. D'après la proposition~\ref{17-proposition-cdpotentiel-morphisme-fini},
pour montrer que $K_X$ est un complexe dualisant potentiel, il suffit de
montrer que $p^!K_X$ en est un. Nous allons bien évidemment le comparer à
$K_{X'}$. Par construction de $K_X$, on a un diagramme commutatif :
\[
\xymatrix{
& i_\star K_Y \ar[rd]& \\
q_\star K_{Y'}\ar[ur]\ar[dr] & & K_X \\
& p_\star K_{X'}\ar[ur] & 
}
\]
Par adjonction, on obtient que deux définitions concurrentes d'un morphisme
$K_{Y'}\to q^!K_X$ coïncident :
\[
\xymatrix{
& {p'}^! K_Y\ar[r] & {p'}^!i^!K_X & \\
K_{Y'}\ar@{=}[ru]\ar@{=}[rd] & & & q^!K_X\ar@{=}[lu]\ar@{=}[ld] \\
& {i'}^! K_{X'}\ar[r] & {i'}^!p^!K_X & 
}
\]
On a déjà montré que le morphisme canonique $K_Y\to i^! K_X$ était un
isomorphisme. Par conséquent, sur le diagramme ci-dessus, toutes les
flèches sont des isomorphismes. Ainsi, le morphisme évident 
$ K_{X'}\to p^!K_X$ induit un isomorphisme non seulement après
application de $j^\star$, mais aussi après celle de $i^!$. Il en résulte
que ce morphisme $K_{X'}\to p^! K_X$ est un isomorphisme. En outre, sur le
diagramme ci-dessus, tous les objets sont naturellement munis d'une
structure de complexe dualisant putatif et tous les isomorphismes, sauf
peut-être celui du bas, sont compatibles aux épinglages. Cet isomorphisme 
${i'}^! K_{X'}\to {i'}^!p^!K_X$ est donc lui aussi compatible aux
épinglages. Par conséquent, l'isomorphisme $K_{X'}\iso p^!K_X$ est
compatible aux épinglages non seulement sur $U$ mais aussi sur $Y'$.
Il en résulte que
$p^!K_X$ est un complexe dualisant potentiel, et on a vu que l'on pouvait
en conclure que $K_X$ en était un aussi.

En outre, il résulte aussitôt de la construction que l'hypothèse selon
laquelle les complexes dualisants potentiels sur $X'$, $Y$ et $Y'$ sont
pervers pour la fonction de perversité $-2\delta$ que $K_X$ est aussi
pervers pour $-2\delta$.

Pour conclure, il s'agit de montrer que si
$K$ et $L$ deux complexes dualisants potentiels sur $X$,
on a un isomorphisme 
privilégié $\Lambda\iso \tau_{\leq 0}\R\SheafHom(K,L)$ qui
donne naissance à un morphisme $\psi\colon K\to L$ 
qui soit un isomorphisme de
complexes dualisants potentiels. En effet, cela montrera que si $\phi\colon
K\to L$ est un autre isomorphisme, alors $\phi=\lambda\cdot \psi$ où
$\lambda\colon X\to \Lambda^\times$ est une fonction localement constante.
Demander que $\phi$ soit compatible aux épinglages impliquant que
$\lambda=1$, on aura bien un unique isomorphisme $K\iso L$ de
complexes dualisants potentiels.

D'après la propriété d'unicité des complexes dualisants potentiels sur
$X'$, $Y$ et $Y'$, on a des isomorphismes de complexes dualisants
potentiels $p^!K\iso p^!L$ et $i^! K\iso i^!L$ induisant le
même isomorphisme $q^!K\iso q^!L_Y$. On a construit plus haut un
triangle distingué canonique :
\[q_\star q^\star K\vers{(+,-)} i_\star i^! K\oplus p_\star p^!
K\vers{(+,+)} K\to q_\star q^\star K[1]\;\text{.}\]
En lui appliquant $\R \SheafHom(-,L)$, on obtient un autre triangle
distingué :
\[\R\SheafHom(K,L)\vers{(+,+)} i_\star
\R\SheafHom(i^!K,i^!L)\oplus p_\star \R\SheafHom(p^!K,p^!L)
\vers {(+,-)} q_\star \R\SheafHom(q^!K,q^!L)\vers +
\]
L'énoncé du théorème~\ref{17-theoreme-existence-unicite-cdp} pour $X'$, $Y$ et
$Y'$ implique aussitôt que les objets de cohomologie 
${\mathscr H}^q\R\SheafHom(K,L)$ sont nuls pour $q<0$, et, compte
tenu de la suite exacte canonique de faisceaux :
\[0\to \Lambda\vers {(+,+)} i_\star\Lambda\oplus p_\star\Lambda\vers{(+,-)}
q_\star\Lambda\to 0\;\text{,}\]
que l'on a un isomorphisme privilégié $\Lambda\iso {\mathscr
H}^0\R\SheafHom(K,L)$. Le morphisme $K\to L$ correspondant induit
bien entendu les uniques isomorphismes de complexes dualisants potentiels
$i^!K\iso i^!L$ et $p^!K\iso p^!L$. Ce morphisme $K\to L$ est donc
compatible aux épinglages sur $Y$ et sur $U$ : c'est un isomorphisme de
complexes dualisants potentiels. Ceci achève la démonstration de la
proposition~\ref{17-proposition-existence-unicite-cdp-recollement}.

%>>>2
\subsection{Cas général}%<<<2

Montrons le théorème~\ref{17-theoreme-existence-unicite-cdp} dans le cas
général. Pour le montrer pour tous les $\ZZn$-schémas noethériens
excellents munis d'une fonction de dimension, par récurrence noethérienne,
on peut supposer le résultat connu pour les schémas finis sur un fermé
d'intérieur vide de $X$. On peut supposer que $X$ est réduit. Notons
$p\colon X'\to X$ la normalisation de $X$. Le morphisme $p$ est fini
surjectif et induit un isomorphisme au-dessus de l'ouvert dense de
normalité $U$ du schéma excellent $X$. Posons $Y=(X-U)_\red$ et formons
le carré cartésien suivant :
\[
\xymatrix{Y'\ar[r]^{i'} \ar[d]^{p'}& X'\ar[d]^p \\ Y\ar[r]^i & X}
\]
Comme $X'$ est normal, la 
proposition~\ref{17-proposition-theoreme-existence-unicite-cdp-cas-normal}
montre que l'énoncé du théorème~\ref{17-theoreme-existence-unicite-cdp} est
connu pour $X'$. L'hypothèse de récurrence noethérienne montre que c'est
aussi le cas pour $Y$ et $Y'$. La 
proposition~\ref{17-proposition-existence-unicite-cdp-recollement} donne la
conclusion souhaitée pour $X$.

%>>>2
%>>>1
\section{Le théorème de dualité locale}%<<<1
\label{17-section-dualite}

\subsection{Énoncé du théorème}%<<<2
\begin{theoreme2}\label{17-theoreme-dualite-locale}
Soit $X$ un $\ZZn$-schéma noethérien excellent 
muni d'une fonction de dimension $\delta$.
Soit $K$ le complexe dualisant potentiel de $(X,\delta)$. Alors
\begin{itemize}
\item $K\in\detalebctf X$ ;
\item $K$ est de dimension quasi-injective finie si et seulement si $X$ est
de dimension de Krull finie ;
\item le foncteur $\Dual_K=\R\SheafHom(-,K)$ préserve $\detalebc X$ ;
\item pour tout $M\in\detalebc X$, le morphisme de bidualité $M\to
\Dual_K\Dual_K M$ est un isomorphisme.
\end{itemize}
En particulier, si $X$ est de dimension de Krull finie, $K$ est un complexe
dualisant au sens de \sga{5}{i}{1.7}.
\end{theoreme2}

%>>>2
\subsection{Constructibilité, tor-dimension, dimension quasi-injective}%<<<2
\subsubsection{Changement de coefficients}%<<<3

\begin{proposition3}\label{17-proposition-cdp-changement-coefficients}
Soit $\Lambda=\mathbf{Z}/n\mathbf{Z}$. Soit $m$ un diviseur de $n$. Soit
$\Lambda'=\mathbf{Z}/m\mathbf{Z}$. L'anneau $\Lambda'$ est une
$\Lambda$-algèbre. Soit $K\in\detaleplus X$ un complexe dualisant potentiel
sur $X$
relativement à l'anneau de coefficients $\Lambda$. Alors
$K'=\R\SheafHom_\Lambda(\Lambda',K)\in\detalepluscoefs X {\Lambda'}$ est
naturellement muni d'une structure de complexe dualisant potentiel
relativement à l'anneau de coefficients $\Lambda'$. De plus, si
$M\in\detalepluscoefs X {\Lambda'}$, on a un isomorphisme canonique dans
$\detaleplus X$ :
\[\R\SheafHom_\Lambda(M,K)\simeq
\R\SheafHom_{\Lambda'}(M,K')\;\text{.}\]
\end{proposition3}

Ceci résulte facilement de la commutation des foncteurs de cohomologie à
supports avec le foncteur $\R\SheafHom_\Lambda(\Lambda',-)$.

%>>>3
\subsubsection{Constructibilité, tor-finitude}%<<<3

\begin{proposition3}\label{17-proposition-constructibilite-cdp}
Soit $K$ un complexe dualisant potentiel sur un $\ZZn$-schéma noethérien
excellent muni d'une fonction de dimension $\delta$. Alors $K\in\detalebc X$.
\end{proposition3}

On va utiliser le lemme suivant :

\begin{lemme3}\label{17-lemme-devissage-constructibilite}
Soit $X$ un $\ZZn$-schéma noethérien excellent. Soit $i\colon Z\to X$ une
immersion fermée. Soit $j\colon U\to X$ l'immersion ouverte complémentaire.
Soit $K\in\detaleplus X$. Les conditions suivantes sont équivalentes :
\begin{itemize}
\item $K\in\detalebc X$ ;
\item $j^\star K\in\detalebc U$ et $i^!K\in\detalebc Z$.
\end{itemize}
\end{lemme3}

Ce lemme résulte aussitôt du fait non trivial que le foncteur
$\R j_\star $
envoie $\detalebc U$ dans $\detalebc X$, cf.~\refext{XIII}{13-constructibilité}.

\bigskip

Démontrons la proposition. On peut supposer $X$ réduit. Comme $X$ est
excellent, $X$ admet un ouvert dense régulier. Notons $Z$ le sous-schéma
fermé $(X-U)_{\red}$ et $i\colon Z\to X$ son immersion fermée dans $X$.
Par récurrence noethérienne,
on peut supposer que le complexe dualisant potentiel $i^!K$ de $Z$ est dans
$\detalebc Z$. En vertu du lemme, on est ramené à montrer que $j^\star
K\in\detalebc U$. Autrement dit, on peut supposer que $X$ est régulier. On peut
supposer de plus que $X$ est connexe. Notons $\eta$ le point générique de
$X$. D'après la proposition~\ref{17-proposition-cdpotentiel-regulier} et le
théorème~\ref{17-theoreme-existence-unicite-cdp}, on a un
isomorphisme canonique $K\simeq\Lambda(\delta(\eta))[2\delta(\eta)]$.
Ainsi, $K$ appartient bien à $\detalebc X$, ce qui achève la démonstration de
la proposition.

\begin{proposition3}
Soit $X$ un schéma noethérien excellent muni d'une fonction de dimension
$\delta$. Le complexe dualisant potentiel de $(X,\delta)$ appartient à
$\detalebctf X$.
\end{proposition3}

On sait déjà que le complexe dualisant potentiel $K$ de $(X,\delta)$
appartient à $\detalebc
X$. Il s'agit donc d'obtenir un résultat de tor-finitude pour $K$. Pour
cela, on peut supposer que $\Lambda=\mathbf{Z}/\ell^\nu\mathbf{Z}$ où
$\ell$ est un nombre premier et $\nu\geq 1$. D'après la
proposition~\ref{17-proposition-cdp-changement-coefficients},
$\R\SheafHom_\Lambda(\mathbf{Z}/\ell\mathbf{Z},K)$ est un complexe
dualisant potentiel relativement à l'anneau de coefficients
$\mathbf{Z}/\ell\mathbf{Z}$ ; d'après la 
proposition~\ref{17-proposition-constructibilite-cdp},
cet objet appartient à $\detalebccoefs X {\mathbf{Z}/\ell\mathbf{Z}}$, le
critère du lemme suivant permet de conclure que $K$ appartient à
$\detalebtf X$.

\begin{lemme3}\label{17-lemme-tor-finitude-z-l-n}
Soit $X$ un schéma noethérien. Soit $\ell$ un nombre premier. Soit $\nu\geq
1$. On pose $\Lambda={\mathbf{Z}/\ell^\nu\mathbf{Z}}$.
Soit $K\in\detaleb X$. Les conditions
suivantes sont équivalentes :
\begin{itemize}
\item[(i)] $K\in\detalebtf X$ ;
\item[(ii)] $K\otimes^{\L}_{\Lambda} \mathbf{Z}/\ell\mathbf{Z}\in
\detalebcoefs X {\mathbf{Z}/\ell\mathbf{Z}}$ ;
\item[(iii)]
$\R\SheafHom_\Lambda(\mathbf{Z}/\ell\mathbf{Z},K)\in\detalebcoefs X
{\mathbf{Z}/\ell\mathbf{Z}}$.
\end{itemize}
Dans ce cas, on a un isomorphisme canonique
$\R\SheafHom_\Lambda(\mathbf{Z}/\ell\mathbf{Z},K)\simeq
K\otimes_{\Lambda}^{\L}\mathbf{Z}/\ell\mathbf{Z}$ dans $\detalebcoefs X
{\mathbf{Z}/\ell\mathbf{Z}}$.
\end{lemme3}

L'équivalence entre (i) et (ii) est simplement indiquée pour mémoire. Il
s'agit ici de montrer que les conditions (ii) et (iii) sont équivalentes.
On représente $K$ par un complexe borné. On peut considérer le complexe
complexe double dont les colonnes sont représentées ci-dessous (les degrés
horizontaux étant indiqués en exposant) :
\[\dots\vers {\ell^{\nu-1}} \overset {-1} K\vers {\ell}
\overset 0 K \vers {\ell^{\nu-1}} \overset 1 K\vers
{\ell}\overset 2 K \vers {\ell^{\nu-1}} \dots \]
On note $C$ le complexe simple de faisceaux de
$\Lambda$-modules sur $X$ associé à ce complexe
double\;\footnote{Ce complexe étant concentré sur un nombre fini de lignes,
il n'y a pas lieu de préciser si l'on définit le complexe simple en termes
de sommes ou de produits.}. De façon évidente, on a un triangle distingué
dans $\detale X$ :
\[
C\vers p K\otimes_{\Lambda}^{\L}\mathbf{Z}/\ell\mathbf{Z}\to
\R\SheafHom_\Lambda(\mathbf{Z}/\ell\mathbf{Z},K)\vers i
C[1]\;\text{.}
\]

On remarque que :
\begin{itemize}
\item ${\mathscr H}^qC\simeq {\mathscr H}^{q+2}C$ pour tout
$q\in\mathbf{Z}$ ;
\item $p$ induit un isomorphisme sur les objets de cohomologie
${\mathscr H}^q$ pour $q$ suffisamment petit ;
\item $i$ induit un isomorphisme sur les objets de cohomologie
${\mathscr H}^q$ pour $q$ suffisamment grand.
\end{itemize}

Comme
$K\otimes_{\Lambda}^{\L}\mathbf{Z}/\ell\mathbf{Z}\in\detalemoinscoefs X 
{\mathbf{Z}/\ell\mathbf{Z}}$ et 
$\R\SheafHom_\Lambda(\mathbf{Z}/\ell\mathbf{Z},K)\in \detalepluscoefs X 
{\mathbf{Z}/\ell\mathbf{Z}}$, il en résulte aussitôt que les conditions
suivantes sont équivalentes :
\begin{itemize}
\item $C\simeq 0$ ;
\item $\R\SheafHom_\Lambda(\mathbf{Z}/\ell\mathbf{Z},K)\in\detaleb X$ ;
\item $K\otimes_{\Lambda}^{\L}\mathbf{Z}/\ell\mathbf{Z}\in\detaleb X$.
\end{itemize}

Si on suppose que $K\in\detalebctf X$, 
pour obtenir l'isomorphisme 
$\R\SheafHom_\Lambda(\mathbf{Z}/\ell\mathbf{Z},K)\simeq
K\otimes_{\Lambda}^{\L}\mathbf{Z}/\ell\mathbf{Z}$ dans $\detalebcoefs X
{\mathbf{Z}/\ell\mathbf{Z}}$, on peut représenter $K$ par un complexe borné
de $\Lambda$-modules plats. On peut conclure en utilisant l'isomorphisme
évident $K/\ell K \iso {}_\ell K$ induit 
par la multiplication par $\ell^{\nu-1}$, où ${}_\ell K$ désigne le
noyau de la multiplication par $\ell$ sur $K$. En effet, on montre
facilement que pour tout complexe
borné de $\Lambda$-modules plats le morphisme évident ${}_\ell K\to
\R\SheafHom_\Lambda(\mathbf{Z}/\ell\mathbf{Z},K)$ est un isomorphisme
dans $\detalepluscoefs X {\mathbf{Z}/\ell\mathbf{Z}}$.

%>>>3
\subsubsection{Préservation de \texorpdfstring{$\detalebc X$}{Dbc(X,Lambda)}}%<<<3

\begin{proposition3}\label{17-proposition-la-dualite-preserve-dbc}
Soit $X$ un $\ZZn$-schéma noethérien excellent muni d'une fonction de dimension
$\delta$. Notons $K_X$ le complexe dualisant potentiel de $(X,\delta)$.
Alors, le foncteur $\Dual_X=\R\SheafHom(-,K_X)$ préserve $\detalebc X$.
\end{proposition3}

Grâce à la proposition~\ref{17-proposition-cdp-changement-coefficients}, on
peut supposer que l'anneau de coefficients est $\mathbf{Z}/\ell\mathbf{Z}$,
avec $\ell$ un nombre premier. Il s'agit de montrer que pour tout faisceau
constructible de $\Lambda$-modules $\mathscr M$ sur $X$,
$\Dual_X \mathscr M\in\detalebc X$. D'après la 
proposition~\ref{17-proposition-constructibilite-cdp}, si $\mathscr M$ est
constant, on a bien $\Dual_X\mathscr M\in\detalebc X$. Si $p\colon Y\to X$ est un
morphisme étale, on a un isomorphisme évident $p^\star
\Dual_X\mathscr M\simeq \Dual_Y p^\star\mathscr M$ pour tout
faisceau constructible $\mathscr M$ sur $X$. Si $\mathscr M$ est localement
constant, en choisissant pour $p$ un morphisme étale surjectif tel que
$p^\star \mathscr M$ soit constant,
on obtient que $\Dual_X\mathscr M\in\detalebc X$ si $\mathscr M$
est localement constant.

On raisonne alors par récurrence noethérienne. Pour tout faisceau
constructible $\mathscr M$, il existe un ouvert dense $U$ de $X$ sur lequel
$\mathscr M$ est localement constant. On note $j\colon U\to X$ l'immersion
ouverte correspondante et $i\colon Z\to X$ une immersion fermée
complémentaire. D'après ce qui précède, on a $j^\star \Dual_X\mathscr
M=\Dual_U{\mathscr M}_{|U}\in\detalebc U$. Par ailleurs, $i^!\Dual_X\mathscr
M\simeq \Dual_Zi^\star \mathscr M$. Par hypothèse de récurrence
noethérienne, on obtient que $i^!\Dual_X\mathscr M$ appartient à
$\detalebc Z$. Le
lemme~\ref{17-lemme-devissage-constructibilite} permet de conclure que
$\Dual_X\mathscr M$ appartient à $\detalebc X$.

\medskip

La stabilité de $\detalebc X$ par $\Dual_X$ permet d'énoncer le résultat
important suivant :

\begin{proposition3}\label{17-proposition-comportement-dualite-morphisme-regulier}
Soit $p\colon X'\to X$ un morphisme régulier
entre $\ZZn$-schémas noethériens excellents. On suppose $X$ muni
d'une fonction de dimension $\delta_X$ et on munit $X'$ de la fonction de
dimension $\delta_{X'}$ définie dans 
la proposition~\ref{17-proposition-stabilite-cdp-morphisme-regulier}. Alors,
pour tout $L\in \detalebc X$, on a un isomorphisme canonique $p^\star\Dual_X
L\iso \Dual_{X'}p^\star L$ et le morphisme de bidualité
$p^\star L\to\Dual^2_{X'} p^\star L$
s'identifie à l'image par $p^\star$ du morphisme de
bidualité $L\to \Dual^2_X L$.
\end{proposition3}

Cela résulte aussitôt des
propositions~\ref{17-proposition-stabilite-cdp-morphisme-regulier},
\ref{17-proposition-rhom-morphismes-reguliers} et
\ref{17-proposition-la-dualite-preserve-dbc}.

%>>>3
\subsubsection{Dimension quasi-injective}%<<<3

\begin{proposition3}\label{17-proposition-dimension-quasi-injective}
Soit $X$ un $\ZZn$-schéma noethérien excellent muni d'une fonction de
dimension $\delta$. Notons $K_X$ le complexe dualisant potentiel de
$(X,\delta)$. La dimension quasi-injective de $K_X$ est $-2\inf_{x\in
X}\delta(x)$. En particulier, elle est finie si et seulement si $X$ est de
dimension de Krull finie.
\end{proposition3}

Commençons par minorer la dimension quasi-injective de $K_X$. Soit $x\in
X$. On note $i\colon Z\to X$ l'inclusion du sous-schéma intègre de 
$X$ de point générique $x$. On a un isomorphisme canonique 
$i^!K_X\simeq i^\star \R\SheafHom(\Lambda_Z,K_X)$. Le complexe
$K_Z=i^!K_X$ est un complexe dualisant potentiel pour $(Z,\delta_{|Z})$. Par
conséquent $(K_Z)_{\overline x}\simeq \Lambda(\delta(x))[2\delta(x)]$. Il
en résulte que la dimension quasi-injective de $K_X$ est au moins
$-2\delta(x)$. On obtient ainsi la minoration
\[-2\inf_{x\in X}\delta(x)\leq \dimquasiinj K_X\;\text{.}\]

Montrons que cette inégalité est en fait une égalité si $X$ est de
dimension de Krull finie. On peut procéder par récurrence sur la dimension
de $X$. On peut en outre supposer que $X$ est local strictement
hensélien (réduit) de point fermé $x$.
Soit $\mathscr M$ un faisceau constructible de
$\Lambda$-modules sur $X$. Il s'agit de montrer que $\Dual_X\mathscr
M\in\detaletronque X {\leq -2\delta(x)}$.
Il existe un ouvert affine dense $U$ sur lequel
$\mathscr M$ est localement constant. Le schéma $X$ étant réduit et
excellent, quitte à rétrécir $X$, 
on peut supposer que $U$ est régulier. Notons $j\colon U\to X$ l'immersion
de $U$. Notons
$U_1,\dots,U_n$ les composantes connexes de $U$, et $\eta_1,\dots,\eta_n$
les points génériques de ces composantes. Le schéma $U$ étant régulier, on
connaît la structure du complexe dualisant potentiel
$K_U$ : pour tout $1\leq i\leq n$, on a un
isomorphisme canonique $K_{U_i}\simeq
\Lambda(\delta(\eta_i))[2\delta(\eta_i)]$. En particulier,
$K_{U^i}\in\detaletronque {U_i}{\leq -2\delta(\eta_i)}$. Le faisceau ${\mathscr
M}_{|U_i}$ étant localement constant, on obtient que $\Dual_{U_i}{\mathscr
M}_{|U_i}\in\detaletronque {U_i}{\leq -2\delta(\eta_i)}$. D'après le théorème de
Lefschetz affine appliqué aux immersions ouverts affines
$j_i\colon U_i\to X$, il
vient alors que $\R {j_i}_\star \Dual_{U_i}{\mathscr
M}_{|U_i}$ appartient à $\detaletronque X {\leq \dim
\overline{\left\{\eta_i\right\}}-2\delta(\eta_i)}$. Comme on a  
$\dim \overline{\left\{\eta_i\right\}}-2\delta(\eta_i)\leq -2\delta(x)$,
il vient que $\R j_\star \Dual_U{\mathscr M}_{|U}$ appartient à 
$\detaletronque X {\leq -2\delta(x)}$.

Notons $i\colon Z\to X$ une immersion fermée complémentaire à $j$. Grâce à
la récurrence sur la dimension, on sait que 
$\Dual_Zi^\star \mathscr  M\in\detaletronque Z {\leq -2\delta(x)}$. En utilisant le
triangle distingué canonique
\[i_\star \Dual_Zi^\star {\mathscr M}\to \Dual_X\mathscr M\to
\R j_\star \Dual_U{\mathscr M}_{|U}\;\text{,}\]
on obtient bien que $\Dual_X\mathscr M\in\detaletronque X {\leq -2\delta(x)}$, ce
qui achève la démonstration de la proposition.

%>>>3
%>>>2
\subsection{Le théorème en degré négatif ou nul}%<<<2

\begin{proposition2}\label{17-proposition-dualite-locale-degre-negatif}
Soit $X$ un $\ZZn$-schéma noethérien excellent muni d'une fonction de
dimension $\delta$. Soit $\mathscr M$ un faisceau constructible de
$\Lambda$-modules sur $X$. Alors, le morphisme canonique est un
isomorphisme dans $\detalebc X$.
\[\mathscr M\iso \tau_{\leq 0}\Dual_X\Dual_X \mathscr M\;\text{.}\]
\end{proposition2}

Au cours de cette démonstration, on dira qu'un faisceau constructible de
$\Lambda$-modules $\mathscr M$ sur $X$ est faiblement réflexif si le
morphisme canonique $\mathscr M\to \tau_{\leq 0}\Dual_X\Dual_X \mathscr
M$ de la proposition est un isomorphisme.

D'après le théorème~\ref{17-theoreme-existence-unicite-cdp}, on sait que
$\Lambda$ est faiblement réflexif. Si $g\colon Z\to X$ est une immersion
fermée et $\mathscr N$ un faisceau constructible de $\Lambda$-modules sur
$Z$, il est clair que $\mathscr N$ est faiblement réflexif sur $Z$ si et
seulement si $g_\star \mathscr N$ est faiblement réflexif sur $X$. Plus
généralement, si $f\colon Y\to X$ est un morphisme fini et $\mathscr N$ un
faisceau constructible de $\Lambda$-modules sur $Y$, alors $\mathscr N$ est
faiblement réflexif si et seulement si $f_\star \mathscr N$
l'est (voir \sga{5}{i}{1.13}). Notons aussi qu'une utilisation
appropriée du lemme des cinq montre que si on a une suite
exacte courte $0\to {\mathscr M}' \to \mathscr M\to {\mathscr M}''\to 0$ de
faisceaux constructibles de $\Lambda$-modules, et que ${\mathscr M}''$ est
faiblement réflexif, alors $\mathscr M$ est faiblement réflexif si et
seulement si ${\mathscr M}'$ est faiblement réflexif.

Grâce à la stabilité par extension énoncée plus haut et à la
proposition~\ref{17-proposition-cdp-changement-coefficients}, on peut supposer
que $\Lambda=\mathbf{Z}/\ell\mathbf{Z}$ où $\ell$ est un nombre premier.
Des remarques précédentes, il résulte que si $f\colon Y\to X$ est fini et
que $U$ est un ouvert de $Y$, alors $f_\star \Lambda_U$ est faiblement
réflexif. La classe des faisceaux constructibles de $\Lambda$-modules
faiblement réflexifs sur $X$ étant stable par facteurs directs et
extensions, on peut conclure en utilisant le dévissage des faisceaux
constructibles de \sga{4}{ix}{5.8}.

%>>>2
\subsection{L'argument de \texorpdfstring{\sgaquatreetdemi{Th.~finitude}{4.3}}{[SGA 4\textonehalf{} [Th. finitude] 4.3]}}%<<<2

\begin{definition2}
Soit $X$ un $\ZZn$-schéma 
noethérien excellent muni d'une fonction de dimension
$\delta$. Soit $\mathscr M$ un faisceau constructible de $\Lambda$-modules
sur $X$. On dit que $\mathscr M$ est réflexif si le morphisme de
bidualité $\mathscr M\to \Dual_X\Dual_X\mathscr M$ est un isomorphisme. On
dira que le morphisme de bidualité est un isomorphisme pour $X$ si tout
faisceau constructible de $\Lambda$-modules sur $X$ est réflexif.
\end{definition2}

\begin{proposition2}\label{17-proposition-argument-de-sga-4-et-demi}
Soit $d\geq 0$.
Si le morphisme de bidualité est un isomorphisme pour les $\ZZn$-schémas
excellents noethériens de dimension $\leq d$, alors il l'est aussi pour les
schémas de type fini sur de tels schémas.
\end{proposition2}

\begin{remarque2}
Dans \sgaquatreetdemi{Th.~finitude}{4.3}, un tel
isomorphisme de bidualité est construit pour les schémas de type
fini sur un $\ZZn$-schéma régulier de dimension au plus un (mais
non nécessairement excellent). Aux hypothèses d'excellence et de régularité
du schéma de base près, il s'agit essentiellement du cas $d=1$ de la
proposition. La démonstration qui suit reprend et généralise celle de 
\sgaquatreetdemi{Th.~finitude}{4.3}.
\end{remarque2}

Par récurrence sur $d$, on peut supposer que l'isomorphisme de bidualité
est un isomorphisme pour les schémas de type fini (et leurs hensélisés
stricts) sur des $\ZZn$-schémas noethériens excellents de dimension $<d$.
Pour montrer que l'isomorphisme de bidualité est un isomorphisme pour tout
schéma $Y$ de type fini sur $X$, il suffit évidemment de le faire pour
$Y=(\mathbf{P}^1)^n\times X$ pour tout $n\in\mathbf{N}$.

Montrons par récurrence sur $n$ que pour tout $\ZZn$-schéma noethérien
excellent $X$ de dimension au plus $d$ (muni d'une fonction de dimension
$\delta$), le morphisme de bidualité est un isomorphisme pour
$(\mathbf{P}^1)^n\times X$. L'hypothèse de la proposition règle le cas
$n=0$. Supposons la propriété établie jusqu'au cran $n-1$, avec $n\geq 1$.
Soit $X$ un
schéma noethérien excellent de dimension $d$ muni d'une fonction de
dimension $\delta$. Montrons que le morphisme de bidualité est un
isomorphisme pour $(\mathbf{P}^1)^n\times X$. On peut supposer que $X$ est
local strictement hensélien, de point fermé $x$. Soit $\mathscr M$ un
faisceau constructible sur $(\mathbf{P}^1)^n\times X$. Notons $C$ un cône
du morphisme de bidualité $\mathscr M\to \Dual\Dual \mathscr M$. On sait
déjà que $C\in\detalebc {(\mathbf{P}^1)^n\times X}$
(cf.~proposition~\ref{17-proposition-la-dualite-preserve-dbc}). Nous
allons dans un premier temps montrer que les faisceaux de cohomologie de
$C$ sont en gratte-ciel, c'est-à-dire supportés par des points fermés.
L'hypothèse de récurrence sur $d$ montre que le support de $C$ est contenu
dans le fermé
$(\mathbf{P}^1)^n\times x$. Posons $Y=(\mathbf{P}^1\times X)_{(y)}$ où
$y$ est le point générique de $\mathbf{P}^1\times x\subset
\mathbf{P}^1\times X$. On considère les $n$ projections canoniques
$(\mathbf{P}^1)^n\times X\to \mathbf{P}^1\times X$ et leur changement de
base $(\mathbf{P}^1)^{n-1}\times Y\to Y$ au-dessus de $Y$. Le schéma $Y$
étant de dimension $d$ et la 
proposition~\ref{17-proposition-comportement-dualite-morphisme-regulier} montrant en
particulier que \guil{la dualité commute aux localisations},
l'hypothèse de récurrence pour $n-1$ implique que si
$z\in (\mathbf{P}^1)^n\times X$ est tel que $C_{\overline z}\neq 0$, alors
les images de $z$ par les $n$ projections canoniques $(\mathbf{P}^1)^n\times
X\to \mathbf{P}^1\times X$ sont des points fermés (puisqu'elles sont
au-dessus de $x$ et différentes du point générique $y$ de
$\mathbf{P}^1\times x$).
Par conséquent, un tel point $z$ est un point fermé de
$(\mathbf{P}^1)^n\times x$. Bref, les faisceaux de cohomologie de $C$ sont
supportés par des points fermés. Il en résulte que si on note  $\pi\colon
(\mathbf{P}^1)^n\times X\to X$ le morphisme canonique, alors pour montrer
que $C\simeq 0$, il suffit de montrer que $\R\pi_\star C\simeq 0$.
D'après \sgaquatreetdemi{Th.~finitude}{4.4},
$\R\pi_\star C$ s'identifie
au cône du morphisme de bidualité $\R\pi_\star \mathscr
M\to\Dual_X\Dual_X \R\pi_\star \mathscr M$, qui est un isomorphisme
par hypothèse. Par conséquent, $C\simeq 0$, ce qui achève la démonstration
de la proposition.

%>>>2
\subsection{Fin de la démonstration}%<<<2

Démontrons le théorème~\ref{17-theoreme-dualite-locale}. Compte tenu des
résultats antérieurs, il ne reste plus qu'à montrer que le morphisme de
bidualité est un isomorphisme pour tout $\ZZn$-schéma noethérien excellent
$X$ muni d'une fonction de dimension $\delta$. Comme il suffit d'obtenir la
conclusion pour les hensélisés stricts de $X$, on peut supposer que $X$ est
de dimension de Krull finie $d$. On va raisonner par récurrence sur $d$.

\begin{definition2}
Soit $\mathscr M$ un faisceau constructible de $\Lambda$-modules sur un
$\ZZn$-schéma noethérien excellent $X$ muni d'une fonction de dimension
$\delta$. Pour tout $q\geq 1$, on dit que $\mathscr M$ vérifie la propriété
$\propD_q$ si le faisceau de cohomologie ${\mathscr H}^q(\Dual_X^2\mathscr
M)$ est nul.
\end{definition2}

D'après la proposition~\ref{17-proposition-dualite-locale-degre-negatif}, 
$\mathscr M$ est réflexif si et seulement s'il vérifie la propriété
$\propD_q$ pour tout $q\geq 1$.

La fonction de dimension $\delta$ sert à formuler la
propriété $\propD_q$. Pourtant, elle n'en dépend pas. En effet, si $\delta$
et $\delta'$ sont deux fonctions de dimension sur $X$ (connexe), il existe
un entier relatif $k$ tel que $\delta'=\delta+k$. Le complexe dualisant
potentiel $K_{X,\delta'}$ s'identifie canoniquement à
$K_{X,\delta}(k)[2k]$ : les foncteurs de bidualité $\Dual_{X,\delta}^2$ et
$\Dual_{X,\delta'}^2$ sont canoniquement isomorphes. Il n'y a donc pas lieu
de mentionner la fonction de dimension dans la notation $\Dual_X^2$, et les
propriétés $\propD_q$ définies relativement à $\delta$ et $\delta'$
sont équivalentes.

En outre, les propriétés $\propD_q$ sont clairement locales pour la
topologie étale. Comme les schémas noethériens excellents admettent
localement pour la topologie étale des fonctions de dimension, on peut leur
donner un sens même en l'absence d'une fonction de dimension globale. Par
recollement, on peut même donner un sens aux faisceaux de cohomologie
${\mathscr H}^q(\Dual_X^2\mathscr M)$.

\medskip

Soit $d\geq 0$. On suppose que l'isomorphisme de bidualité est un
isomorphisme pour tout $\ZZn$-schéma noethérien excellent de dimension au
plus $d-1$ muni d'une fonction de dimension. 

Nous allons montrer par récurrence sur $q\geq 1$ que tout faisceau
constructible de $\Lambda$-modules $\mathscr M$ sur un schéma noethérien
excellent $X$ de dimension $\leq d$ vérifie la propriété $\propD_q$.

Soit $q\geq 1$. On suppose que pour tout $1\leq q'<q$, tout faisceau de
$\Lambda$-modules sur un schéma noethérien excellent $X$ de dimension $\leq
d$ vérifie la propriété $\propD_{q'}$.

\begin{lemme2}\label{17-lemme-propriete-d-extensions-et-sous-objets}
Les entiers $d$ et $q$ ayant été fixés comme ci-dessus, la propriété
$\propD_q$ pour les faisceaux constructibles de $\Lambda$-modules sur les
schémas noethériens excellents de dimension au plus $d$ est stable par
extensions et sous-objets.
\end{lemme2}

En effet, si on une suite
exacte courte $0\to{\mathscr M}'\to\mathscr M\to {\mathscr M}''\to 0$
de faisceaux constructibles de $\Lambda$-modules sur un
tel schéma $X$, l'hypothèse de récurrence si $q\geq 2$ ou la
proposition~\ref{17-proposition-dualite-locale-degre-negatif} si $q=1$
implique que l'on a une suite exacte de faisceaux :
\[0\to \mathscr \HH^q(\Dual_X^2{\mathscr M}')\to 
\mathscr \HH^q(\Dual_X^2{\mathscr M})\to
\mathscr \HH^q(\Dual_X^2{\mathscr M}'')\]
La propriété $\propD_q$ est donc stable par extensions et sous-objets.

\begin{lemme2}\label{17-lemme-propriete-d-et-morphismes-finis}
Soit $p\colon Y\to X$ un morphisme fini entre schémas noethériens
excellents. Soit $\mathscr M$ un faisceau constructible de
$\Lambda$-modules sur $Y$. Alors, $\mathscr M$ vérifie la propriété
$\propD_q$ si et seulement si $p_\star\mathscr M$ la vérifie.
\end{lemme2}

Ceci résulte aussitôt de l'isomorphisme canonique
$p_\star {\mathscr H}^q(\Dual_Y^2\mathscr M)\simeq {\mathscr
H}^q(\Dual_X^2p_\star \mathscr M)$ (voir \sga{5}{i}{1.12~(a)})
et de la conservativité du foncteur $p_\star$.

\begin{lemme2}\label{17-lemme-devissage-standard-normal}
Soit $X$ un schéma noethérien excellent. 
Soit $\mathscr C$ une sous-catégorie strictement pleine de la catégorie
$\Cons(X,\Lambda)$ des faisceaux constructibles de $\Lambda$-modules sur
$X$, stable par facteurs directs et extensions. On suppose que pour tout
morphisme fini $p\colon Y\to X$, toute immersion ouverte $j\colon U\to Y$
avec $Y$ normal intègre
et tout nombre premier $\ell$ divisant $n$, on a $p_\star j_!
\mathbf{Z}/\ell\mathbf{Z}\in\mathscr C$. Alors, $\mathscr
C=\Cons(X,\Lambda)$.
\end{lemme2}

Il s'agit d'une variante facile de \sga{4}{ix}{5.8}.

\begin{lemme2}
Les entiers $d$ et $q$ ayant été fixés comme ci-dessus, 
si la propriété $\propD_q$ est satisfaite par le faisceau constant $\Lambda$
sur les schémas noethériens excellents normaux (strictement henséliens)
de dimension au plus $d$, alors la propriété $\propD_q$ est satisfaite par
tout faisceau constructible de $\Lambda$-modules sur un schéma noethérien
excellent de dimension au plus $d$.
\end{lemme2}

D'après le lemme~\ref{17-lemme-devissage-standard-normal}, il suffit d'établir
la propriété $\propD_q$ pour un faisceau de la forme $p_\star
j_!\mathbf{Z}/\ell\mathbf{Z}$ avec $\ell$ un diviseur premier de $n$, $p$
un morphisme fini et $j$ une immersion ouverte entre schémas normaux
intègres. D'après le 
lemme~\ref{17-lemme-propriete-d-et-morphismes-finis}, il suffit d'établir la
propriété $\propD_q$ pour $j_!\mathbf{Z}/\ell\mathbf{Z}$ avec $j\colon U\to
Y$ une immersion ouverte, avec $Y$ normal intègre. D'après la 
stabilité par sous-objet et extensions de la propriété $\propD_q$
(cf. lemme~\ref{17-lemme-propriete-d-extensions-et-sous-objets}), il suffit de
traiter le cas du faisceau $j_!\Lambda$, qui est lui-même un sous-faisceau
du faisceau constant $\Lambda$ sur le schéma normal $Y$, ce qui achève la
démonstration du lemme.

\bigskip

On est ainsi ramené à montrer la propriété $\propD_q$ pour le faisceau
constant $\Lambda$ sur les schémas noethériens excellents normaux $X$ de
dimension $d$. On peut supposer
$X$ local strictement hensélien de point fermé $x$ et de point générique
$\eta$. Si $d\leq 1$, $X$ est régulier, et alors, si on choisit la
fonction de dimension $\delta$ sur $X$ telle que $\delta(\eta)=0$,
le complexe dualisant potentiel associé $K_X$ sur $X$ est le
faisceau constant $\Lambda$, et alors il est évident que $\Lambda$ vérifie
la propriété $\propD_q$ puisque l'on a tautologiquement
$\Dual_X\Dual_X\Lambda\simeq \Lambda$. On
peut donc supposer que $d\geq 2$. En appliquant le lemme suivant à la
complétion $\widehat{X}\to X$, on voit qu'on peut supposer que $X$ est
complet :

\begin{lemme2}\label{17-lemme-propriete-d-et-morphismes-reguliers}
Soit $q\geq 1$.
Soit $p\colon X'\to X$ un morphisme régulier
entre $\ZZn$-schémas noethériens excellents. Soit $\mathscr M$
un faisceau constructible de $\Lambda$-modules sur $X$.
On suppose que $p$ est surjectif.
Alors, le faisceau $\mathscr M$ vérifie la propriété $\propD_q$ si et
seulement si $p^\star\mathscr M$ la vérifie.
\end{lemme2}

Ceci résulte aussitôt de la 
proposition~\ref{17-proposition-comportement-dualite-morphisme-regulier}.

\bigskip

Il nous reste à montrer que si $X$ est un schéma local strictement
hensélien noethérien normal complet de dimension $d\geq 2$, alors
le faisceau constant $\Lambda$ sur $X$ vérifie la propriété $\propD_q$.
D'après le théorème d'algébrisation
partielle \refext{V}{5-algébrisation partielle mixte raffinée géométrique},
il existe un morphisme $p\colon X'\to X$ fini surjectif tel que
\begin{itemize}
\item le schéma $X'$ soit normal ;
\item il existe un schéma local noethérien complet $Y$ de
dimension $<d$, un morphisme de type fini $Z\to Y$, un point géométrique
$\overline z\to Z$ et un isomorphisme $X'\simeq \widehat{Z_{(\overline z)}}$.
\end{itemize}

Le schéma $Y$ est noethérien excellent et de dimension $<d$. L'hypothèse
de récurrence sur $d$ et la
proposition~\ref{17-proposition-argument-de-sga-4-et-demi} impliquent que le
morphisme de bidualité est un isomorphisme pour $Z_{(\overline z)}$.  En
particulier, le faisceau constant $\Lambda$ sur $Z_{(\overline z)}$ est
réflexif. Appliqué au morphisme de complétion $X'\to Z_{(\overline z)}$,
la proposition~\ref{17-proposition-comportement-dualite-morphisme-regulier}
montre que le faisceau constant $\Lambda$ sur $X'$ est réflexif.
En particulier, le faisceau constant $\Lambda$ sur $X'$ satisfait
la propriété $\propD_q$.
D'après le lemme~\ref{17-lemme-propriete-d-et-morphismes-finis}, on peut en
déduire que le faisceau
$p_\star\Lambda$ sur $X$ vérifie la propriété $\propD_q$. Le morphisme fini
$p$ étant surjectif, le morphisme canonique $\Lambda\to p_\star\Lambda$ est
un monomorphisme de faisceaux. La propriété $\propD_q$ étant stable par
sous-objets (cf. lemme~\ref{17-lemme-propriete-d-extensions-et-sous-objets}),
le faisceau constant $\Lambda$ sur $X$ vérifie bien la propriété
$\propD_q$, ce qui achève la démonstration du théorème de dualité locale.

%>>>2
%>>>1
\section{Anneaux de coefficients généraux}%<<<1
\label{17-section-coefficients-generaux}
\subsection{Énoncés}%<<<2

\begin{definition2}
Soit $X$ un schéma noethérien. Soit $A$ un anneau noethérien. On appelle
complexe dualisant sur $\detalebccoefs X A$ (resp. $\detalebctfcoefs X A$)
un objet $K\in\detalebccoefs X A$ tel
que le foncteur $\Dual_K=\R\SheafHom(K,-)$ préserve $\detalebccoefs X A$
(resp. $\detalebctfcoefs X A$)
et que pour tout $M\in\detalebccoefs X A$ (resp. $M\in\detalebctfcoefs X A$),
le morphisme de bidualité $M\to \Dual^2_K M$ soit un isomorphisme.
\end{definition2}

Cette section vise à établir les deux théorèmes suivants :

\begin{theoreme2}\label{17-theoreme-dualite-d-b-c-coefficients}
Soit $A$ une $\Lambda$-algèbre noethérienne. Soit $X$ un schéma
noethérien. S'il en existe, les complexes dualisants sur $\detalebccoefs X A$
sont uniques au produit tensoriel près avec des objets inversibles.
Soit $R\in\dtotale(A)$ un complexe ponctuellement
dualisant fort au sens de \cite[page~120]{GDBC@Conrad}\;\footnote{On rappelle
que cela signifie ici que $R$ appartient à
$\dbc(A)$ et que pour tout $x\in \SP(A)$,
$R_{(x)}\in\dtotale(A_{(x)})$ est un complexe dualisant pour $A_{(x)}$ au sens
de \cite[page~258]{RD@Hartshorne}, ce qui signifie que $R_{(x)}$ est de
dimension injective finie et que le foncteur
$\R\Hom(-,R_{(x)})$ induit une involution de
$\dbc(A_{(x)})$. D'après,
\cite[lemma~3.1.5]{GDBC@Conrad}, il revient au même de demander que le foncteur
$\R\Hom_A(-,R)$ induise une involution de
$\dbc(A)$.}.
Soit $K$ un complexe dualisant sur $\detalebccoefs X \Lambda$. Alors,
$R\overset{\L}\otimes_{\Lambda} K$ est un complexe dualisant sur
$\detalebccoefs
X A$.
\end{theoreme2}

\begin{theoreme2}\label{17-theoreme-dualite-d-b-ctf-coefficients}
Soit $A$ une $\Lambda$-algèbre noethérienne. Soit $X$ un schéma
noethérien. S'il en existe, les
complexes dualisants sur $\detalebctfcoefs X A$ sont uniques au produit
tensoriel près avec des objets inversibles.
Soit $K$ un complexe dualisant sur $\detalebccoefs X \Lambda$.
Alors, $A\overset{\L}\otimes_{\Lambda} K$ est un complexe dualisant sur
$\detalebctfcoefs X A$.
\end{theoreme2}

%>>>2
\subsection{Systèmes locaux}%<<<2
\label{17-subsection-systemes-locaux}

\begin{definition2}
Soit $X$ un schéma noethérien. On appelle système local (d'ensembles)
sur $X$
un faisceau d'ensembles sur $X_\et$ isomorphe à une limite inductive
filtrante de faisceaux représentés par des revêtements étales finis de $X$.
Un système local fini est un système local représenté par un revêtement
étale fini.
\end{definition2}

\begin{proposition2}\label{17-proposition-systemes-locaux-ind-objets-dans-rev}
Soit $X$ un schéma noethérien. La catégorie des systèmes locaux sur $X$ est
équivalente à la catégorie $\Ind(\Rev(X))$ des ind-objets dans la catégorie
$\Rev(X)$ des revêtements étales finis de $X$.
\end{proposition2}

Le foncteur qui à un revêtement étale fini $Y\to X$ associe le faisceau
d'ensembles sur $X_\et$ représenté par $Y$ est pleinement fidèle. En
utilisant \sga{1}{v}{8.7.5 a)}, on en déduit, par passage à la limite
inductive, un foncteur pleinement fidèle de la catégorie $\Ind(\Rev(X))$
vers celle des faisceaux d'ensembles sur $X_\et$. Par définition, l'image
essentielle de ce foncteur est la catégorie des systèmes locaux.

\begin{proposition2}
Soit $X$ un schéma noethérien connexe. Soit $\overline x$ un point
géométrique de $X$. Le foncteur qui à un système local $\mathscr F$ associe
la fibre ${\mathscr F}_{\overline x}$ est naturellement muni d'une action
de $\pi_1(X,\overline x)$ et définit une équivalence entre la
catégorie des systèmes locaux sur $X$ et
la catégorie $\pi_1(X,\overline x)-\Ens$ des ensembles sur lesquels le
groupe profini $\pi_1(X,\overline x)$ agit continûment.
Autrement dit, la catégorie des
systèmes locaux d'ensembles sur $X$
s'identifie à la catégorie des faisceaux
d'ensembles sur le topos classifiant du
groupe profini $\pi_1(X,\overline x)$.
\end{proposition2}

Notons $\Ensf$ la catégorie des ensembles finis.
D'après \sga{1}{v}{7}, le foncteur $\Rev(X)\to \Ensf$ qui à $Y$ associe
l'ensemble sous-jacent au schéma $Y_{\overline x}$ s'enrichit d'une action
du groupe $\pi_1(X,\overline x)$ pour définir une équivalence de catégories 
$\Rev(X)\iso \pi_1(X,\overline x)-\Ensf$ où $\pi_1(X,\overline x)-\Ensf$
est la catégorie des ensembles finis (discrets) munis d'une action continue
du groupe profini $\pi_1(X,\overline x)$. En passant cette équivalence aux
ind-objets, on obtient une équivalence $\Ind(\Rev(X))\iso
\pi_1(X,\overline x)-\Ens$, ce qui permet de conclure d'après la
proposition~\ref{17-proposition-systemes-locaux-ind-objets-dans-rev}.

\bigskip

À partir de la définition des systèmes locaux d'ensembles, on peut définir
les systèmes locaux de groupes abéliens, de groupes abéliens de torsion, de
modules, etc.

\begin{proposition2}\label{17-proposition-stabilites-systemes-locaux}
Pour tout schéma noethérien $X$, la catégorie des systèmes locaux
d'ensembles (resp. de groupes abéliens) sur $X$
admet des limites inductives et des limites projectives finies et le
foncteur d'inclusion de la catégorie des systèmes locaux d'ensembles (resp.
de groupes abéliens) sur $X$ dans la
catégorie des faisceaux d'ensembles (resp. de groupes abéliens)
sur $X_\et$ y commute.

Soit $p\colon Y\to X$ un morphisme entre schémas noethériens.
Si $\mathscr F$ est un système local sur $X$, alors $p^\star \mathscr F$
est un système local sur $Y$ ; la réciproque est vraie si $p$ est un
revêtement étale fini surjectif.

Si $\mathscr G$ est un système local sur $Y$ et $p$ un revêtement étale
fini, alors $p_\star \mathscr G$ est un système local sur $X$.
\end{proposition2}

Le seul fait non trivial à montrer ici est que si $p$ est un revêtement
étale fini surjectif et $\mathscr F$ un faisceau d'ensembles sur $X_\et$
tel que $p^\star \mathscr F$ soit un système local, alors $\mathscr F$ est
un système local. On a un isomorphisme évident entre $\mathscr F$ et
l'égalisateur des deux morphismes évidents $p_\star p^\star \mathscr F\to
p_\star p^\star p_\star p^\star \mathscr F$ déduits du couple de foncteurs
adjoints $(p^\star,p_\star)$. D'après les autres propriétés triviales de
stabilité des systèmes locaux, $p_\star p^\star \mathscr F$ et $p_\star
p^\star p_\star p^\star \mathscr F$ sont des systèmes locaux ;
l'égalisateur de deux morphismes entre systèmes locaux est encore un
système local.

\begin{proposition2}\label{17-proposition-devissage-systemes-locaux}
Soit $X$ un schéma noethérien. La catégorie abélienne des systèmes locaux
de groupes abéliens de torsion sur $X$ est stable par extension dans la
catégorie des faisceaux de groupes abéliens sur $X_\et$. Plus précisément,
si $0\to \mathscr F'\to\mathscr F\to\mathscr F''\to 0$ est une suite exacte
de faisceaux de groupes abéliens sur $X$ telle que $\mathscr F'$ et
$\mathscr F''$ soient des systèmes locaux et que $\mathscr F'$
soit de torsion, alors $\mathscr F$ est un système local de groupes
abéliens.
\end{proposition2}

Cette proposition résulte aussitôt du lemme suivant :

\begin{lemme2}
Soit $X$ un schéma noethérien. Soit $\mathscr G$ un faisceau de groupes
abéliens de torsion agissant librement sur un faisceau d'ensembles
$\mathscr T$. On note $p\colon \mathscr T\to\mathscr T/\mathscr G$ le
morphisme quotient. Si $\mathscr G$ et $\mathscr T/\mathscr G$ sont des
systèmes locaux, alors $\mathscr T$ aussi.
\end{lemme2}

On peut supposer que $X$ est connexe.
Notons $\mathscr Y=\mathscr T/\mathscr G$. On peut écrire $\mathscr Y$
comme réunion de sous-systèmes locaux finis $\mathscr Y'$. Pour chacun de
ces $\mathscr Y'$, on peut considérer $p^{-1}(\mathscr Y')$ : il s'agit
d'un faisceau d'ensembles sur $X_\et$ sur lequel $\mathscr G$ agit
librement avec $\mathscr Y'$ pour quotient. Le faisceau $\mathscr T$ étant
réunion des sous-faisceaux $p^{-1}(\mathscr Y')$ associés, pour montrer que
$\mathscr T$ est un système local, il suffit de montrer que pour tout
sous-système local fini $\mathscr Y'$ de $\mathscr Y$, $p^{-1}(\mathscr
Y')$ est un système local. Bref, on peut supposer que $\mathscr Y$ est un
système local fini.

On suppose ainsi que $\mathscr Y$ est représenté par un revêtement étale
fini $q\colon Y\to X$. On peut identifier le faisceau $\mathscr T$ sur $X$
muni du morphisme $p\colon \mathscr T\to \mathscr Y$ à un faisceau étale
$\tilde{\mathscr T}$ sur $Y$ \emph{via}
un isomorphisme canonique $\mathscr T\simeq q_!\tilde{\mathscr T}$.
L'action libre de $\mathscr G$ sur $\mathscr T$ correspond à une action
libre du système local $q^\star \mathscr G$ sur $\tilde{\mathscr T}$.
D'après la proposition~\ref{17-proposition-stabilites-systemes-locaux}, pour
montrer que $\mathscr T$ est un système local, il suffit de montrer que
$\tilde{\mathscr T}$ en est un. Bref, on peut supposer que $Y=X$.

On s'est ramené à la situation où $\mathscr T/\mathscr G$ est l'objet final
de la catégorie des faisceaux sur $X$, c'est-à-dire que $\mathscr T$ est un
torseur sous $\mathscr G$. Si $\mathscr G$ est un système local fini, alors
$\mathscr T$ est représentable par un revêtement étale fini et est donc un
système local ; nous allons nous ramener à ce cas-là.

La classe d'isomorphisme du $\mathscr G$-torseur $\mathscr T$ est définie
par un élément dans l'ensemble $\HH^1(X_\et,\mathscr G)$. Comme
$\HH^1(X_\et, -)$ commute aux limites inductives filtrantes, il existe un
sous-système local de groupes abéliens finis
$\mathscr G'$ de $\mathscr G$ (supposé de torsion),
un $\mathscr G'$-torseur $\mathscr T'$ et un $\mathscr G$-isomorphisme
$\mathscr T\simeq \mathscr G\otimes_{\mathscr G'}\mathscr T'$ où l'on a
noté $\otimes$ le foncteur d'extension du groupe structural (cf.
\cite[proposition~1.3.6, Chapitre~III]{CNA@Giraud}). L'extension du groupe
structural commutant aux limites inductives filtrantes, $\mathscr T$
s'identifie à la limite inductive des $\mathscr G''\otimes_{\mathscr
G'}\mathscr T'$ pour $\mathscr G''$ parcourant l'ensemble ordonné des
sous-systèmes locaux de groupes abéliens finis de $\mathscr G$ contenant
$\mathscr G'$. D'après ce qui précède, $\mathscr G''\otimes_{\mathscr
G'}\mathscr T'$ est un système local d'ensembles sur $X$.  Par passage à la
limite inductive, $\mathscr T$ est bien un système local.

\bigskip

Le résultat de l'exercice suivant montre que l'hypothèse \guil{de torsion}
est bien nécessaire dans la
proposition~\ref{17-proposition-devissage-systemes-locaux}, et que par
ailleurs, un faisceau qui est un système local localement pour la topologie
étale n'est pas forcément un système local.

\begin{exercice2}
Soit $A$ le sous-anneau de $\mathbf{C}[t]$ formé des
polynômes $f$ tels que $f(0)=f(1)$ : le schéma $C=\SP(A)$ correspondant
est obtenu en identifiant $0$ et $1$ dans la droite affine complexe
$\mathbf{A}^1_{\mathbf{C}}$.

\begin{itemize}
\item Montrer que $C$ est isomorphe à la cubique plane d'équation $x^3-y^2+xy=0$
dans le plan affine complexe $\SP(\mathbf{C}[x,y])$
(envoyer $x$ et $y$ respectivement sur $t(t-1)$ et $t^2(t-1))$.

\item Montrer que $C$ admet un unique point singulier $O$.

\item Montrer que le morphisme évident $p\colon
\mathbf{A}^1_{\mathbf{C}}\to C$ est le normalisé de $C$ et que le
sous-schéma fermé (réduit) $p^{-1}(O)$ de $\mathbf{A}^1_{\mathbf{C}}$
est $\{0,1\}$.

\item Construire un isomorphisme $\HH^1_\et(C,\mathbf{Z})\simeq \mathbf{Z}$.

\item Montrer qu'il existe un faisceau de groupes abéliens sur
$C_\et$ tel que :
\begin{enumerate}
\item[(i)] $\mathscr F$ soit extension de deux systèmes locaux, et s'insère
plus précisément dans une suite exacte courte 
$0\to\mathbf{Z}\to\mathscr F\to\mathbf{Z}\to 0$ ;
\item[(i')] Localement pour la topologie étale, $\mathscr F$ soit un système
local ;
\item[(ii)] $\mathscr F$ ne soit pas un système local.
\end{enumerate}
\end{itemize}
\end{exercice2}

%>>>2
\subsection{Partitions galoisiennes}%<<<2

\begin{definition2}
Une partition galoisienne d'un schéma noethérien $X$ consiste en la donnée
d'une partition finie de $X$ par des sous-schémas (localement fermés)
réduits connexes $(S_i)_{i\in I}$ et d'un revêtement étale galoisien
$S'_i\to S_i$ pour tout $i\in I$.
\end{definition2}

\begin{definition2}
Soit $p\colon Y\to X$ un revêtement fini étale galoisien entre schémas
noethériens. Soit $A$ une
$\Lambda$-algèbre.
On dit qu'un système local de $A$-modules $\mathscr F$ sur $X$ est 
rendu ind-unipotent par $Y$, si pour un point géométrique $\overline y$ de
$Y$ (et donc pour tous) au-dessus d'un point géométrique $\overline x$ de
$X$, le $\pi_1(X,\overline x)$-module discret ${\mathscr F}_{\overline x}$
est ind-unipotent pour le sous-groupe distingué $\pi_1(Y,\overline y)$
(cf.~sous-section~\ref{17-subsection-modules-ind-unipotents-sous-groupe-distingue}),
autrement dit que $\mathscr F$ est limite inductive filtrante de faisceaux
localement constants extensions successives de faisceaux dont l'image
inverse par $p$ soit un faisceau constant.
\end{definition2}

\begin{definition2}
Soit $X$ un schéma noethérien muni d'une partition galoisienne $\mathscr
P=(S'_i\to S_i)_{i\in I}$. Soit $A$ une $\Lambda$-algèbre.
On dit d'un faisceau de
$A$-modules sur $X$ qu'il est faiblement constructible par rapport à
$\mathscr P$ si pour tout $i\in I$, sa restriction à $S_i$ est un système
local rendu ind-unipotent par $S'_i$. On note $\FCons^{\mathscr P}(X,A)$ la
sous-catégorie pleine de la catégorie des faisceaux de $A$-modules sur $X$
formée des faisceaux faiblement constructibles pour $\mathscr P$.
Si $\mathscr P'$ est une deuxième partition galoisienne, on dit que
$\mathscr P'$ raffine $\mathscr P$ si on a l'inclusion $\FCons^{\mathscr
P}(X,\Lambda)\subset \FCons^{\mathscr P'}(X,\Lambda)$ (et donc aussi
$\FCons^{\mathscr P}(X,A)\subset \FCons^{\mathscr
P'}(X,A)$ pour toute $\Lambda$-algèbre $A$).
\end{definition2}

\begin{proposition2}
Soit $X$ un schéma noethérien muni d'une partition galoisienne. Soit $A$
une $\Lambda$-algèbre.
La catégorie $\FCons^{\mathscr P}(X,A)$ est abélienne ; son foncteur
d'inclusion dans la catégorie des faisceaux de $A$-modules sur $X$ est
exact et commute aux limites inductives. $\FCons^{\mathscr P}(X,A)$ est 
stable par extensions dans la catégorie des faisceaux de $A$-modules
sur $X$.
\end{proposition2}

Ceci résulte aussitôt des propriétés des modules ind-unipotents pour un
sous-groupe
(cf.~proposition~\ref{17-proposition-proprietes-unipotents-sous-groupe-distingue})
et des propriétés générales des systèmes locaux
(cf.~sous-section~\ref{17-subsection-systemes-locaux}).

\begin{definition2}
On dit d'une partition galoisienne $\mathscr P=(S'_i\to S_i)_{i\in I}$ sur
un schéma noethérien $X$
qu'elle est dirigée si on a muni $I$ d'un ordre total tel que, soit $I$ est
vide, soit, si on note $i_0$ le plus petit élément de $I$, $S_{i_0}$ est
ouvert et, récursivement, $(S'_i\to S_i)_{i\in I-\{i_0\}}$ est une
partition galoisienne dirigée du fermé réduit $X-S_{i_0}$. On dit qu'une
partition galoisienne est dirigeable s'il existe un ordre total sur
l'ensemble d'indices qui en fasse une partition galoisienne dirigée.
\end{definition2}

\begin{proposition2}\label{17-proposition-raffinements-dirigeables}
Toute partition galoisienne d'un schéma noethérien est raffinée par une
partition galoisienne dirigeable.
\end{proposition2}

\begin{lemme2}
Soit $X'\to X$ un revêtement étale galoisien. Soit $(S_i\to X)_{i\in I}$
une partition de $X$ par un nombre fini de sous-schémas réduits connexes.
On note $S'_i$ une composante connexe du produit fibré $S_i\times_X X'$.
Alors, $(S'_i\to S_i)_{i\in I}$ est une partition galoisienne de $X$ qui
raffine la partition galoisienne $(X'\to X)$.
\end{lemme2}

Cela résulte aussitôt de la théorie de Galois.

\bigskip

Montrons la proposition par récurrence noethérienne sur $X$. Soit $\mathscr
P=(S'_i\to
S_i)_{i\in I}$ une partition galoisienne d'un schéma noethérien non vide
$X$. Tout d'abord, montrons que, quitte à 
raffiner $\mathscr P$, on peut supposer qu'il
existe un indice $i_0\in I$ tel que $S_{i_0}$ soit un ouvert. En effet, si
on choisit un $i_0\in I$ tel que $S_{i_0}$ contienne un point
maximal de $X$, $S_{i_0}$ contient un ouvert non vide $U$ de
$S_{i_0}$. Notons $V_1,\dots,V_n$ les composantes connexes du
fermé réduit $S_{i_0}-U$ de $S_{i_0}$.
D'après le lemme, il existe une partition
galoisienne $\mathscr Q=
(U'\to U,V'_1\to V_1,\dots,V'_n\to V_n)$ de $S_{i_0}$ qui
raffine la partition galoisienne $(S'_{i_0}\to S_{i_0})$ de $S_{i_0}$.
Quitte à remplacer la partition galoisienne initiale de $X$ par son
raffinement $\mathscr Q\cup \mathscr P'$ avec
$\mathscr P'=(S'_i\to S_i)_{i\in I-\{i_0\}}$, on peut effectivement
supposer que $S_{i_0}$ est ouvert.

On peut appliquer l'hypothèse de
récurrence noethérienne à la partition galoisienne $\mathscr P'$ 
de $X-S_{i_0}$ pour en obtenir un raffinement $\mathscr P''$ indexé par un
certain ensemble totalement ordonné $J$ qui fasse de $\mathscr P''$ une
partition galoisienne dirigée de $X-S_{i_0}$. La partition galoisienne
$(S'_{i_0}\to S_{i_0})\cup \mathscr P''$ raffine $\mathscr P$, et si on
étend l'ordre sur $J$ en un ordre sur la réunion disjointe $\{i_0\}\sqcup
J$ de façon à faire de $i_0$ le plus petit élément, on a obtenu une
partition dirigée.

\begin{exercice2}\label{17-exercice-partitions-ensemble-filtrant}
Montrer que si $\mathscr P$ et $\mathscr P'$ sont deux partitions
galoisiennes d'un schéma noethérien $X$, il existe une partition
galoisienne (dirigeable) raffinant à la fois $\mathscr P$ et $\mathscr P'$.
\end{exercice2}

%>>>2
\subsection{Dévissages}%<<<2

Le but de cette sous-section est d'établir le résultat suivant :

\begin{proposition2}\label{17-proposition-devissage-coefficients}
Soit $A$ une $\Lambda$-algèbre noethérienne.
Soit $X$ un schéma noethérien.
Soit $\mathscr T$ une sous-catégorie triangulée strictement pleine de
$\detalebccoefs X A$ stable par facteurs directs. On suppose que pour tout
nombre premier $\ell$ divisant $n$ et tout $A/\ell A$-module $N$, le
foncteur $N\overset{\L}{\otimes}_{\mathbf{Z}/\ell\mathbf{Z}}-\colon
\detalebccoefs X
{\mathbf{Z}/\ell\mathbf{Z}}\to \detalebccoefs X A$ prend ses valeurs dans
$\mathscr T$. Alors, $\mathscr T=\detalebccoefs X A$.
\end{proposition2}

La démonstration de cette proposition est repoussée à la fin de cette
sous-section.

\begin{proposition2}\label{17-proposition-devissage-partitions-galoisiennes-trivial}
Soit $A$ une $\Lambda$-algèbre.
Soit $\mathscr P=(S'_i\to S_i)_{i\in I}$ avec
$I=\{1,\dots,N\}$ une
partition galoisienne dirigée d'un schéma noethérien $X$. 
Notons $k_i\colon S_i\to X$ l'immersion canonique pour tout $i\in I$. Pour
tout $i\in I$, le foncteur ${k_i}_!\colon \FCons^{(S'_i\to S_i)}(S_i,A)\to
\FCons^{\mathscr P}(X,A)$ est pleinement fidèle. Tout objet $\mathscr F$ de
$\FCons^{\mathscr P}(X,A)$ admet une filtration croissante (fonctorielle)
$(\Fil_n \mathscr F)_{n\in\mathbf{Z}}$ telle que le $\Fil_{-1}\mathscr F=0$,
$\Fil_{N}\mathscr F=\mathscr F$ et que pour tout $1\leq i\leq n$, le
quotient $\Fil_i\mathscr F/\Fil_{i-1}\mathscr F$ soit dans
${k_i}_!\FCons^{S'_i\to S_i}(S_i,A)$.
\end{proposition2}

C'est trivial.

\begin{proposition2}\label{17-proposition-objets-simples-faiblement-constructibles}
Soit $A$ une $\Lambda$-algèbre.
Soit $\mathscr P=(S'_i\to S_i)_{i\in I}$ une partition
galoisienne dirigeable d'un schéma noethérien $X$. Notons $k_i\colon S_i\to
X$ l'immersion canonique pour tout $i\in I$.
Si $A$ est noethérien, alors $\FCons^{\mathscr P}(X,A)$ est une catégorie
abélienne localement noethérienne (cf.~\cite[pages~325--326]{CA@Gabriel}).
Si $A$ est artinien, $\FCons^{\mathscr
P}(X,A)$ est localement finie et admet un nombre fini d'objets simples ;
plus précisément, si on note $k_i\colon S_i\to X$ les inclusions
canoniques, $W_i$ un ensemble (fini) représentatif des objets simples de la
catégorie des $\Gal(S'_i/S_i)$-modules (\emph{via} le choix d'un point
géométrique de $S'_i$, on identifie ces objets à des systèmes locaux sur
$S_i$ trivialisés par $S'_i$), alors les objets ${k_i}_!\mathscr F$ pour
$i\in I$ et $\mathscr F\in W_i$ forment un ensemble représentatif des
objets simples de $\FCons^{\mathscr P}(X,A)$.
\end{proposition2}

Supposons $A$ noethérien.
Montrons que $\FCons^{\mathscr P}(X,A)$ est localement noethérienne. On sait
déjà que cette catégorie abélienne admet des limites inductives filtrantes
et que celles-ci sont exactes. Il s'agit de montrer que tout objet de
$\FCons^{\mathscr P}(X,A)$ est limite inductive d'objets noethériens (ou
plus précisément, mais cela revient au même,
\guil{réunion} de ses sous-objets noethériens). Si
l'ensemble d'indice $I$ de $\mathscr P$ est vide, c'est trivial. Sinon, on
peut choisir un ordre total sur $I$ qui fasse de $\mathscr P$ une partition
galoisienne dirigée, et noter $i_0$ le plus petit élément de $I$. Notons
$j\colon S_{i_0}\to X$ l'immersion (ouverte) correspondante et $k\colon
X-S_{i_0}\to X$ l'immersion du fermé réduit complémentaire. Pour tout objet
$\mathscr F\in \FCons^{\mathscr P}(X,A)$, on a une suite exacte courte :
\[0\to j_!j^\star\mathscr F\to \mathscr F\vers\pi k_\star k^\star \mathscr F\to
0\;\text{.}\]
En raisonnant par récurrence sur
le cardinal de $I$, on peut supposer que $k^\star\mathscr F$ est
\guil{réunion} de ses sous-objets noethériens dans $\FCons^{\mathscr
P'}(X-S_{i_0},A)$ avec $\mathscr P'=\mathscr P-(S'_{i_0}\to S_{i_0})$. Bien
entendu, un objet $\mathscr R\in \FCons^{\mathscr P'}(X-S_{i_0},A)$ est
noethérien si et seulement si $k_\star\mathscr R$ l'est dans
$\FCons^{\mathscr P}(X,A)$. On dispose
donc d'un système inductif $({\mathscr H}_a)_{a\in A}$ indexé par un ensemble
ordonné filtrant $A$ de sous-objets noethériens de $k_\star k^\star
\mathscr F$ tel que $\mathscr F$ soit la réunion des sous-objets
$\pi^{-1}(\mathscr H_a)$ de $\mathscr F$ pour $a\in A$.
Si chacun des $\pi^{-1}(\mathscr H_a)$ est réunion de ses sous-objets
noethériens, alors $\mathscr F$ aussi. Ceci permet de supposer que
$k^\star \mathscr F$ est noethérien. Concernant $j^\star\mathscr F$, en
utilisant que 
la catégorie des $A[G]$-modules discrets (avec $G$ groupe profini) est
localement noethérienne, on obtient que l'objet $j^\star\mathscr F$ de
$\FCons^{S'_{i_0}\to S_{i_0}}(S_{i_0},A)$ est réunion de ses sous-objets
noethériens ; on en déduit aussitôt que $j_!j^\star\mathscr F$ est aussi
réunion de ses sous-objets noethériens dans $\FCons^{\mathscr P}(X,A)$. Pour
conclure que $\mathscr F$ est réunion de ses sous-objets noethériens, on
utilise le lemme suivant :

\begin{lemme2}
Soit $A$ une $\Lambda$-algèbre noethérienne.
Soit $\mathscr P$ une partition galoisienne dirigeable
d'un schéma noethérien $X$.
Soit $0\to \mathscr H\to\mathscr F\to \mathscr G\to 0$ une suite exacte
courte dans $\FCons^{\mathscr P}(X,A)$. On suppose que $\mathscr G$ est un
objet noethérien de $\FCons^{\mathscr P}(X,A)$ et que
$\mathscr H$ est réunion de ses sous-objets noethériens.
Alors, $\mathscr F$ est aussi réunion de ses sous-objets noethériens.
\end{lemme2}

Il est évident que $\mathscr G$ est un faisceau de $A$-modules
constructible sur $X$. D'après \sga{4}{ix}{2.7.3}, le foncteur
$\Ext^1(\mathscr G,-)$ de la catégorie des faisceaux de $A$-modules sur $X$
vers celle des $A$-modules commute aux limites inductives filtrantes. La
suite exacte donnée définissant un élément dans $\Ext^1(\mathscr G,\mathscr
H)$ et $\mathscr H$ s'écrivant comme une limite inductive filtrante de ses
sous-objets noethériens, il existe un sous-objet noethérien ${\mathscr
H}'$ de $\mathscr H$, une suite exacte courte $0\to \mathscr H'\to \mathscr
F'\to\mathscr G\to 0$ et un diagramme commutatif de la forme suivante, où
le carré de gauche est cocartésien :
\[
\xymatrix{0\ar[r] & \ar[d]\mathscr H'\ar[r] & \mathscr F'\ar[r]\ar[d]
 & \mathscr G\ar@{=}[d] \ar[r]& 0 \\
0\ar[r] & \mathscr H\ar[r] & \mathscr F\ar[r] &
\mathscr G \ar[r]& 0\;\text{,}}
\]
À vrai dire, on peut remplacer $\mathscr H'$ par tout sous-objet (noethérien)
$\mathscr H''$ de $\mathscr H$ contenant $\mathscr H'$, et $\mathscr F$
s'identifie à la réunion des sous-objets $\mathscr F''$ ainsi définis. Pour
conclure, il suffit de montrer qu'un tel $\mathscr F''$ est un objet 
noethérien de $\FCons^{\mathscr P}(X,A)$, ce qui est évident puisqu'il est
extension de deux objets noethériens $\mathscr H''$ et $\mathscr G$.

\bigskip

Supposons $A$ artinien. Il s'agit de trouver un ensemble fini d'objets
simples à partir desquels tous les objets noethériens s'obtiennent par
extensions successives. Compte tenu du dévissage de la
proposition~\ref{17-proposition-devissage-partitions-galoisiennes-trivial}, on
peut supposer que $\mathscr P$ est constitué d'un unique revêtement
galoisien $X'\to X$. Choisissons un point géométrique $\overline x'$ de
$X'$ au-dessus d'un point géométrique $\overline x$ de $X$. Notons
$G=\pi_1(X,x)$ et $H=\pi_1(X',x)$. Le groupe profini $H$ s'identifie à un
sous-groupe distingué fermé de $G$. Notons $K=G/H$ le groupe fini quotient.
L'anneau $A[K]$ est évidemment artinien à gauche, on peut en noter $W$ un
ensemble représentatif fini d'objets simples ; considérés comme des
$A[G]$-modules discrets, les éléments de $W$ sont encore simples.

La catégorie $\FCons^{\mathscr P}(X,A)$ s'identifie à la catégorie des
$G$-modules discrets ind-unipotents pour $H$. Tout objet de cette catégorie
s'écrit comme une réunion de sous-objets de type fini unipotents pour $H$,
et de tels sous-objets se dévissent eux-mêmes en extensions successives
d'objets sur lesquels $H$ agit trivialement, ces derniers s'identifiant à
des $A[K]$-modules de type fini, ils se dévissent en extensions successives
d'éléments de $W$.

\begin{definition2}
Soit $A$ une $\Lambda$-algèbre noethérienne.
Soit $X$ un schéma noethérien. Soit
$\mathscr P$ une partition galoisienne. On note $\Cons^{\mathscr P}(X,A)$
la sous-catégorie pleine de $\FCons^{\mathscr P}(X,A)$ formée des faisceaux
de $A$-modules constructibles (ce qui revient ici à dire que
les fibres sont des $A$-modules de type fini). Il s'agit bien entendu de la
sous-catégorie abélienne des objets noethériens dans $\FCons^{\mathscr
P}(X,A)$.
\end{definition2}

\begin{proposition2}\label{17-proposition-devissage-gabber}
Soit $A$ une $\Lambda$-algèbre noethérienne. Soit $X$ un schéma noethérien.
Soit $\mathscr P$ une partition galoisienne dirigeable de $X$. Soit $\mathscr
F$ un objet de $\Cons^{\mathscr P}(X,A)$. 
Il existe une filtration finie de $\mathscr F$ dans $\Cons^{\mathscr P}(X,A)$
dont les quotients successifs soient des facteurs directs d'objets de la
forme $M\otimes_{\mathbf{Z}/\ell\mathbf{Z}} {\mathscr F}_0$ où $\ell$ est
un nombre premier divisant $n$, $M$ un $A/\ell A$-module de type fini et
${\mathscr F}_0$ un objet de $\Cons^{\mathscr
P}(X,\mathbf{Z}/\ell\mathbf{Z})$.
\end{proposition2}

D'après la
proposition~\ref{17-proposition-devissage-partitions-galoisiennes-trivial}, on
peut supposer que $\mathscr P$ est constituée d'un unique revêtement
galoisien $X'\to X$. En reprenant les notations utilisées dans la
démonstration du cas artinien de la
proposition~\ref{17-proposition-objets-simples-faiblement-constructibles}, on
peut identifier $\mathscr F$ à un $A[G]$-module discret unipotent pour le
sous-groupe distingué fermé $H$. Cette propriété d'unipotence permet de
supposer que $H$ agit trivialement, de sorte qu'on se retrouve avec une
action du groupe fini $K=G/H=\Gal(X'/X)$ (bref, on peut supposer que 
$\mathscr F$ est un système local de $A$-modules trivialisé par $X'$). Le
lemme suivant appliqué à l'algèbre de groupe $B=A[K]$ permet de conclure.

\begin{lemme2}
Soit $A$ une $\Lambda$-algèbre noethérienne. Soit $B$ une $\Lambda$-algèbre
finie non nécessairement commutative. Tout $A\otimes_\Lambda B$-module (à
gauche) de type fini admet une filtration finie dont les quotients successifs
soient des facteurs directs de $A\otimes_\Lambda B$-modules de la forme
$N\otimes_{\mathbf{F}_\ell} L$ où $\ell$ est un nombre premier divisant
$n$, $N$ un $A/\ell A$-module de type fini et $L$ un $B/\ell B$-module simple.
\end{lemme2}

En premier lieu, comme $A\otimes_\Lambda B$ est évidemment noethérien à
gauche, on peut procéder à une récurrence noethérienne ; il suffit donc de
montrer que tout $A\otimes_\Lambda B$-module non nul admet un sous-module
non nul facteur direct d'un module de la forme $N\otimes_{\mathbf{F}_\ell}
L$ avec $N$ un $A/\ell A$-module de type fini, $L$ un $B/\ell B$-module
simple et $\ell$ un nombre premier divisant $n$. Ceci permet de supposer
que $B$ est un anneau semi-simple. En effet, pour tout $B$-module non nul
$M$, l'annulateur du radical de Jacobson $\mathscr N$ de $B$ dans $M$
est un sous-$B$-module non nul de $M$ (cf.~\cite[§4]{GTM131@Lam}) ;
si $M$ est un $A\otimes_{\Lambda} B$-module non nul, l'annulateur de
$\mathscr N$ dans $M$ s'identifie donc à un $A\otimes_{\Lambda}(B/\mathscr
N)$-module non nul et l'anneau $B/\mathscr N$ est bien semi-simple.

En deuxième lieu, l'énoncé du lemme est vrai pour un produit
$B=B_1\times\dots\times B_k$ d'anneaux si et seulement s'il est vrai pour
chacun des $B_i$ et l'énoncé est aussi invariant par équivalence de Morita
(cf.~\cite[§18]{GTM189@Lam}) puisqu'on peut le formuler intrinsèquement
en termes de la catégories des $B$-modules. Compte tenu du théorème
d'Artin-Wedderburn de structure des anneaux semi-simples
(cf.~\cite[3.5]{GTM131@Lam}), on peut donc supposer que $B$ est un
corps fini, \emph{a priori} non commutatif, mais effectivement commutatif
en vertu du théorème de Wedderburn (cf.~\cite[13.1]{GTM131@Lam}).

En troisième lieu, l'énoncé est vrai dans le cas particulier auquel on
s'est ramené ci-dessus. Soit $B=L$ une extension finie de
$\mathbf{F}_\ell$, pour un certain nombre premier $\ell$ divisant $n$. On
peut supposer que $\ell$ annule $A$.
Soit $M$ un
$A\otimes_{\mathbf{F}_\ell} L$-module. L'extension $L/\mathbf{F}_\ell$ est
galoisienne, notons $G$ son groupe de Galois.
On considère le $A\otimes_{\mathbf{F}_\ell} L$-module
$M'=\oplus_{\sigma\in G} \sigma^\star M$ où $\sigma^\star M$ est le
$A\otimes_{\mathbf{F}_\ell} L$-module obtenu par image inverse par
l'automorphisme de $\SP(A\otimes_{\mathbf{F}_\ell} L)$ induit par
$\sigma^{-1}\colon L\to L$ (on peut identifier $\sigma^\star M$ au
$A$-module sur $M$ sur lequel on ferait agir $\lambda\in L$ par la
multiplication d'origine par $\sigma(\lambda)$).
Le $A\otimes_{\mathbf{F}_\ell} L$-module $M$ est évidemment un facteur
direct de $M'$. On va montrer que $M'$ est de la forme
$N\otimes_{\mathbf{F}_\ell} L$ où $N$ est un $A$-module. On observe que
pour tout $\sigma\in G$, on dispose d'un isomorphisme évident de
$A\otimes_{\mathbf{F}_\ell}L$-modules $M'\simeq \sigma^\star M'$. Ces
isomorphismes vérifient tautologiquement les relations de cocycles qui
permettent de munir le $A\otimes_{\mathbf{F}_\ell} L$-module $M'$ d'une
donnée de descente relative au morphisme de schémas
$\SP(A\otimes_{\mathbf{F}_\ell} L)\to\SP(A)$ (qui est un revêtement étale
surjectif). Par descente fidèlement plate (cf.~\sga{1}{viii}{1.1}),
il existe un $A$-module $N$ tel
que $N\otimes_{\mathbf{F}_\ell} L$ s'identifie à $M'$ de façon compatible
avec les isomorphismes $M'\simeq \sigma^\star M'$. En particulier, on a un
isomorphisme de la forme voulue entre les $A\otimes_{\mathbf{F}_\ell}
L$-modules $M'$ et $N\otimes_{\mathbf{F}_\ell} L$, ce qui permet de
conclure.

\bigskip

Nous sommes maintenant en mesure de démontrer la
proposition~\ref{17-proposition-devissage-coefficients}. Soit $\mathscr T$ une
telle sous-catégorie triangulée de $\detalebccoefs X A$. Soit $\mathscr F$ un
faisceau constructible de $A$-modules sur $X$. Il s'agit de montrer que
$\mathscr F$ appartient à $\mathscr T$. Il existe évidemment une partition
galoisienne $\mathscr P$ telle que $\mathscr F$ appartienne à
$\Cons^{\mathscr P}(X,A)$. D'après la
proposition~\ref{17-proposition-raffinements-dirigeables}, on peut supposer
que $\mathscr P$ est dirigeable. On peut alors appliquer la
proposition~\ref{17-proposition-devissage-gabber} pour conclure que $\mathscr
F$ appartient à $\mathscr T$.

%>>>2
\subsection{Complexes dualisants sur \texorpdfstring{$\detalebccoefs X A$}{Dbc(X,A)}}%<<<2
\subsubsection{Unicité}%<<<3

\begin{proposition3}\label{17-proposition-unicite-complexe-dualisant}
Soit $X$ un schéma noethérien. Soit $A$ un anneau noethérien. Si
$K$ et $K'$ sont deux complexes dualisants sur $\detalebccoefs X A$ (resp.
$\detalebctfcoefs X A$), alors il
existe un objet inversible $L\in\detalebctfcoefs X A$ 
(cf.~proposition~\ref{17-proposition-objets-inversibles-d-t-a} pour plus de
précisions) tel que $K'$ soit isomorphe à $L\otimes^{\L}_{A}K$.
\end{proposition3}

\begin{lemme3}\label{17-lemme-dual-dans-d-b-implique-dans-d-b}
Soit $X$ un schéma noethérien. Soit $A$ un anneau noethérien. On suppose
que $K\in\detalebccoefs X A$ est dualisant. Pour tout
$F\in\detalemoinsccoefs X A$, si $\Dual_K F\in\detalebccoefs X A$, alors
$F\in\detalebccoefs X A$.
\end{lemme3}

Commençons par montrer que l'on peut supposer que $\Dual_K F=0$. Pour tout
$F\in\detalecoefs X A$, notons $\varepsilon_F\colon F\to \Dual_K\Dual_K F$ le
morphisme de bidualité. Le morphisme composé
\[\xymatrix{\Dual_K F\ar[r]^-{\varepsilon_{\Dual_K F}} &\Dual_K\Dual_K\Dual_K
F\ar[r]^-{\Dual_K(\varepsilon_F)}& \Dual_K F}\]
est l'identité de $\Dual_K F$. Comme $\Dual_K F$ est dans $\detalebccoefs X A$,
le fait que $K$ soit dualisant implique que $\varepsilon_{\Dual_K F}$ est
un isomorphisme. Par conséquent, $\Dual_K(\varepsilon_F)$ est un
isomorphisme. Quitte à remplacer $F$ par un cône de $\varepsilon_F$, on
peut supposer que $\Dual_K F=0$.

Par l'absurde, supposons que $F$ soit non nul. Il existe un morphisme
non nul $p\colon F\to F'$ avec $F'\in\detalebccoefs X A$ (par exemple, le
morphisme canonique $F\to \tau_{\geq n}F$ pour un entier $n$ bien choisi).
Considérons le carré commutatif :
\[
\xymatrix{F \ar[r]^p \ar[d]^{\varepsilon_F}& F'\ar[d]^\sim_{\varepsilon_{F'}}\\
\Dual^2_K F\ar[r]^{\Dual^2_K(p)} & \Dual^2_KF'}
\]
D'un côté, $\Dual^2_K F$ est nul, donc $\varepsilon_{F'}\circ p=0$, mais de
l'autre, $\varepsilon_{F'}$ est un isomorphisme, d'où $p=0$, ce qui conduit
à une contradiction.

\bigskip

Établissons la proposition~\ref{17-proposition-unicite-complexe-dualisant}
dans le cas des complexes dualisants sur $\detalebccoefs X A$, la démonstration
qui suit vaudra aussi pour $\detalebctfcoefs X A$ (à ceci près qu'il
ne sera plus alors nécessaire de recourir au lemme ci-dessus).
On pose $Y=\Dual_KK'\in\detalebccoefs X A$. Comme
$K$ est dualisant, on a aussi un isomorphisme privilégié $K'=\Dual_K Y$.
Pour tout $Z\in\detalebccoefs X A$, on a un isomorphisme fonctoriel
$\Dual_{K'}Z\simeq \Dual_K(Z\overset{\L}\otimes_A{Y})$ dans $\detalecoefs
X A$. Grâce au lemme, on en déduit que pour tout $Z\in\detalebccoefs X A$, on a
$Z\overset{\L}\otimes_A Y\in\detalebccoefs X A$. On a ainsi un triangle
commutatif de catégories et de foncteurs (à isomorphismes près de
foncteurs) :
\[
\xymatrix{\ar[rd]_{\Dual_{K'}}\detalebccoefs X A
\ar[r]^-{-\overset{\L}\otimes_A Y} & \ar[d]^{\Dual_K} \detalebccoefs X A\\
& (\detalebccoefs X A)^\opp }
\]
Comme $\Dual_K$ et $\Dual_{K'}$ sont des équivalences, le foncteur
$-\overset{\L}\otimes_A Y$ aussi. En particulier, $Y$ est un complexe
inversible. Notons $Y'$ l'inverse de $Y$. On a un isomorphisme de
foncteurs $\R\SheafHom(Y,-)\simeq Y'\overset{\L}\otimes_{A}-$
(si un foncteur est une équivalence, son adjoint à droite est un
quasi-inverse). Comme $K'=\Dual_K Y$, on peut en déduire que $K'\simeq
Y'\overset{\L}\otimes_A K$.

%>>>3
\subsubsection{Réduction au cas \texorpdfstring{$\Lambda=\mathbf{Z}/\ell\mathbf{Z}$}{Lambda=Z/lZ}}%<<<3

\begin{proposition3}\label{17-proposition-reduction-complexe-dualisant-coh-ferme}
Soit $A$ un anneau noethérien. Soit $R\in \dtotale(A)$ un complexe
ponctuellement dualisant fort. Soit $J$ un idéal de $A$. On pose $A'=A/J$
et $R'=\R\Hom_A(A',R)\in\dtotale(A')$. Alors, $R'$ est un complexe
ponctuellement dualisant fort. Si on note $\Dual$ (resp. $\Dual'$) le
foncteur $\R\Hom_A(-,R)$ (resp. $\R\Hom_{A'}(-,R')$) sur
$\dtotale(A)$ (resp. $\dtotale(A')$) et $\oub\colon
\dtotale(A')\to\dtotale(A)$ le foncteur de \guil{restriction des scalaires},
on a un isomorphisme canonique :
\[\oub\circ \Dual'\simeq \Dual\circ \oub\;\text{.}\]
\end{proposition3}

C'est essentiellement trivial.

\begin{proposition3}\label{17-proposition-complexes-dualisants-quotient-par-ideal}
Soit $X$ un schéma noethérien. Soit $K\in\detalebctfcoefs X \Lambda$.
Soit $A$ une
$\Lambda$-algèbre noethérienne. Soit $J$ un idéal de $A$. Soit
$R\in\dtotale(A)$ un complexe ponctuellement dualisant fort. On pose
$A'=A/J$ et $R'=\R\Hom_A(A',R)\in\dtotale(A')$. On note
$K_R=K\overset{\L}\otimes_\Lambda R\in\detalebccoefs X A$ et $K_{R'}=
K\overset{\L}\otimes_\Lambda {R'}\in\detalebccoefs X {A'}$. On note
$\oub\colon \detalebccoefs X {A'}\to\detalebccoefs X A$ le foncteur conservatif
évident. Alors, pour tout $M\in\detalecoefs X {A'}$, on a un isomorphisme
canonique dans $\detalecoefs X A$ :
\[\oub(\R\SheafHom_{A'}(M,K_{R'}))\simeq \R\SheafHom_A(\oub
M,K_R)\;\text{.}\]
De plus, si $K_R$ est un complexe dualisant sur $\detalebccoefs X A$, alors
$K_{R'}$ en est un sur $\detalebccoefs X {A'}$ et la réciproque est vraie si
$J$ est nilpotent.
\end{proposition3}

Les autres assertions en étant des conséquences faciles, il s'agit de
montrer que l'on a un isomorphisme canonique
$\R\SheafHom_A(A',K_R)\simeq \oub(K_{R'})$ dans $\detalebccoefs X A$, ce
qui résulte de la proposition~\ref{17-proposition-coefficients-universels-2}.

\begin{corollaire3}\label{17-corollaire-reduction-theoreme-dualite-d-b-c-coefs-z-l}
Pour démontrer le théorème~\ref{17-theoreme-dualite-d-b-c-coefficients}, on
peut supposer que $A=\mathbf{Z}/\ell\mathbf{Z}$ où $\ell$ est un nombre
premier.
\end{corollaire3}

Étant entendu que la propriété d'unicité des complexes dualisants a déjà
été obtenu, il est évident que pour démontrer le
théorème~\ref{17-theoreme-dualite-d-b-c-coefficients}, on peut supposer que
$\Lambda=\mathbf{Z}/\ell^\nu\mathbf{Z}$ où $\ell$ est un nombre premier et
$\nu\geq 1$. Posons $A'=A/\ell A$. Notons $R\in\dtotale(A)$ un complexe
ponctuellement dualisant fort. D'après la
proposition~\ref{17-proposition-reduction-complexe-dualisant-coh-ferme}, le
complexe $R'=\R\SheafHom_A(A',R)\in\dtotale(A')$ en est un pour $A'$.
Appliquant dans un premier temps la
proposition~\ref{17-proposition-complexes-dualisants-quotient-par-ideal} au
cas où $\Lambda\to A$ est $\Lambda\to \mathbf{Z}/\ell\mathbf{Z}$, nous
obtenons que $K''=\R\SheafHom_\Lambda(\mathbf{Z}/\ell\mathbf{Z},K)$
est un complexe dualisant sur $\detalebccoefs X {\mathbf{Z}/\ell\mathbf{Z}}$.
D'après le lemme~\ref{17-lemme-tor-finitude-z-l-n}, on a aussi un isomorphisme
$K''\simeq K\overset{\L}\otimes_\Lambda\mathbf{Z}/\ell\mathbf{Z}$
dans $\detalebccoefs X {\mathbf{Z}/\ell\mathbf{Z}}$. Appliquons le
théorème~\ref{17-theoreme-dualite-d-b-c-coefficients} dans le cas de la
$\mathbf{Z}/\ell\mathbf{Z}$-algèbre $A'$ : on obtient que
$K''\overset{\L}\otimes_{\mathbf{Z}/\ell\mathbf{Z}} R'$ est un complexe
dualisant sur $\detalebccoefs X {A'}$. Cet objet
$K''\overset{\L}\otimes_{\mathbf{Z}/\ell\mathbf{Z}} R'$ s'identifie
aussi à $K\overset{\L}\otimes_\Lambda R'=K_{R'}$. D'après la
proposition~\ref{17-proposition-complexes-dualisants-quotient-par-ideal}, il
vient que $K\overset{\L}\otimes_\Lambda R=K_R$ est un complexe
dualisant sur $\detalebccoefs X A$, ce qui achève la démonstration de ce
corollaire.

%>>>3
\subsubsection{Démonstration du théorème~\ref{17-theoreme-dualite-d-b-c-coefficients}}%<<<3

L'énoncé d'unicité des complexes dualisants sur $\detalebccoefs X A$
a déjà été obtenu,
cf.~proposition~\ref{17-proposition-unicite-complexe-dualisant}.
Pour l'énoncé d'existence, d'après le
corollaire~\ref{17-corollaire-reduction-theoreme-dualite-d-b-c-coefs-z-l}, on
peut supposer que $\Lambda=\mathbf{Z}/\ell\mathbf{Z}$ avec $\ell$ un nombre
premier. On se donne $K$ un complexe dualisant sur $\detalebccoefs X
{\mathbf{Z}/\ell\mathbf{Z}}$. Soit $A$ une $\mathbf{Z}/\ell\mathbf{Z}$-algèbre
noethérienne et $R\in\dtotale(A)$ un complexe ponctuellement dualisant
fort. Notons $\Dual_X=\R\SheafHom_{\Lambda}(-,K)$ le foncteur de
dualité sur $\detalebccoefs X {\Lambda}$ induit par $K$ et $\Dual_A$ celui
induit par $R$ sur $\dbc(A)$. Notons $\Dual_{X,A}$ le
foncteur $\R\SheafHom_A(-,K_R)$ sur $\detalecoefs X A$ où
$K_R=K\overset{\L}\otimes_A R$. La
proposition~\ref{17-proposition-coefficients-universels-3} ($\Lambda$ est un
corps) montre que l'on a
un isomorphisme canonique, pour tout $N\in\dbc(A)$ et
$\mathscr F\in\detalebccoefs X {\mathbf{Z}/\ell\mathbf{Z}}$ :
\[\Dual_{X,A}(N
\overset{\L}\otimes_{\mathbf{Z}/\ell\mathbf{Z}} \mathscr F)\simeq
(\Dual_A
N)\overset{\L}\otimes_{\mathbf{Z}/\ell\mathbf{Z}}(\Dual_X\mathscr F)\;
\text{.}\]
Par hypothèse, $\Dual_A N\in\dbc(A)$ et
$\Dual_X\mathscr F\in\detalebccoefs X {\mathbf{Z}/\ell\mathbf{Z}}$, ce qui
permet de déduire que $\Dual_{X,A}(N
\overset{\L}\otimes_{\mathbf{Z}/\ell\mathbf{Z}} \mathscr F)$
appartient à $\detalebccoefs X A$, puis que $\Dual_{X,A}$ préserve
$\detalebccoefs
X A$ grâce au dévissage de la
proposition~\ref{17-proposition-devissage-coefficients}.
Avec les mêmes notations, les morphismes de bidualité $N\to\Dual_A^2N$
et $\mathscr F\to\Dual_X^2\mathscr F$ sont des isomorphismes, le morphisme
de bidualité $N
\overset{\L}\otimes_{\mathbf{Z}/\ell\mathbf{Z}} \mathscr
F\to\Dual_{X,A}^2(
N
\overset{\L}\otimes_{\mathbf{Z}/\ell\mathbf{Z}} \mathscr F)$ en est
donc un aussi ; le même dévissage permet de conclure que $\Dual_{X,A}$
définit
une involution de $\detalebccoefs X A$, c'est-à-dire que $K_R$ est un complexe
dualisant sur $\detalebccoefs X A$, ce qui achève la démonstration du
théorème~\ref{17-theoreme-dualite-d-b-c-coefficients}.

%>>>3
%>>>2
\subsection{Complexes dualisants sur \texorpdfstring{$\detalebctfcoefs X A$}{Dbctf(X,A)}}%<<<2

L'assertion d'unicité des complexes dualisants sur $\detalebctfcoefs X A$ a
déjà été énoncée dans la
proposition~\ref{17-proposition-unicite-complexe-dualisant}. L'essentiel de
cette sous-section vise à établir le théorème suivant, dont on va déduire
dans quelques lignes le
théorème~\ref{17-theoreme-dualite-d-b-ctf-coefficients} :

\begin{theoreme2}\label{17-theoreme-stabilite-d-b-ctf-par-rhom}
Soit $X$ un schéma noethérien. On suppose qu'il existe un complexe
dualisant sur $\detalebc X$. Soit $A$ une $\Lambda$-algèbre noethérienne.
Pour tous $K$ et $L$ objets de $\detalebctfcoefs X A$, l'objet
$\R\SheafHom_A(K,L)$ appartient à $\detalebctfcoefs X A$ et pour tout
$M\in\dplus(A)$, le morphisme canonique est un isomorphisme :
\[M\overset{\L}\otimes_A \R\SheafHom_A(K,L)\iso
\R\SheafHom_A(K,M\overset{\L}\otimes_A L)\;\text{.}\]
Si $A'$ est une $A$-algèbre,
que $K$ et $L$ sont des objets de
$\detalebctfcoefs X A$, alors on a un isomorphisme canonique dans
$\detalebcoefs X
{A'}$ :
\[A'\overset{\L}\otimes_A\R\SheafHom_A(K,L)\iso
\R\SheafHom_{A'}(A'\overset{\L}\otimes_A K,
A'\overset{\L}\otimes_A L)\;\text{.}\]
\end{theoreme2}

Montrons que l'on peut déduire le
théorème~\ref{17-theoreme-dualite-d-b-ctf-coefficients} de 
ce théorème~\ref{17-theoreme-stabilite-d-b-ctf-par-rhom} et du
théorème~\ref{17-theoreme-dualite-d-b-c-coefficients} qui a déjà été établi.
Commençons par un lemme :

\begin{lemme2}
Soit $A$ un anneau noethérien. Soit $K\in\dmc(A)$.
Alors $K$ est nul si et seulement si pour 
tout idéal maximal $\mathfrak m$ de $A$, l'objet
$A/\mathfrak m\overset{\L}\otimes_A K$ est nul. 
\end{lemme2}

On suppose que $K$ n'est pas nul. On veut montrer qu'il existe un idéal
maximal $\mathfrak m$ de $A$ tel que $A/\mathfrak
m\overset{\L}\otimes_A K$ ne soit pas nul.
Soit $q$ le plus grand entier tel que
$\HH^q(K)$ soit non nul. On peut supposer que $K$ est un complexe formé de
$A$-modules projectifs de type fini et nuls en degrés strictement plus
grands que $q$. Par construction du produit tensoriel dérivé, on a une
surjection $\HH^q(A/\mathfrak m\overset{\L}\otimes_A K)\to
\HH^q(K)/\mathfrak m \HH^q(K)$ pour tout idéal maximal $\mathfrak m$ de $A$. 
Pour conclure, il suffit donc de montrer qu'il existe un idéal maximal
$\mathfrak m$ tel que $\HH^q(K)\neq \mathfrak m \HH^q(K)$ ou encore, d'après le
lemme de Nakayama, que $\HH^q(K)\otimes_A A_{\mathfrak m}\neq 0$. Le support
de $\HH^q(K)$ est un fermé non vide de $\SP(A)$ (cf.
\egazero{i}{1.7}), il contient un point fermé que l'on
identifie à un idéal maximal $\mathfrak m$ de $A$,
et cet idéal maximal vérifie la condition voulue.

\bigskip

Soit $K$ un
complexe dualisant sur $\detalebc X$ et $A$ une $\Lambda$-algèbre noethérienne.
Grâce au lemme~\ref{17-lemme-tor-finitude-z-l-n}, il vient que $K$ appartient
à $\detalebctf X$ et donc que
$K_A=A\overset{\L}\otimes_\Lambda K$ appartient
à $\detalebctfcoefs X A$. D'après le
théorème~\ref{17-theoreme-stabilite-d-b-ctf-par-rhom}, le foncteur
$\Dual_A=\R\SheafHom_A(-,K_A)$ préserve $\detalebctfcoefs X A$. Il reste à
montrer que le morphisme de bidualité $L\to \Dual_A^2 L$ est un
isomorphisme pour tout $L\in\detalebctfcoefs X A$.
D'après le lemme, il suffit de
montrer qu'après produit tensoriel dérivé avec $A/\mathfrak m$, le
morphisme $L\to\Dual_A^2 L$ induit un isomorphisme. D'après le
théorème~\ref{17-theoreme-stabilite-d-b-ctf-par-rhom}, le foncteur de dualité
considéré commute au changement d'anneau, ainsi, après produit tensoriel
avec $A/\mathfrak m$, on obtient le morphisme de bidualité pour
$A/\mathfrak m\overset{\L}\otimes_A L$ dans $\detalebctfcoefs X
{A/\mathfrak m}$. Bref, on peut supposer que l'anneau $A$ est un corps.
Dans ce cas, on peut conclure en utilisant le
théorème~\ref{17-theoreme-dualite-d-b-c-coefficients}.

\begin{proposition2}\label{17-proposition-r-j-star-preserve-d-b-c-si-cplx-dualisant}
Soit $X$ un schéma noethérien. On suppose qu'il existe un complexe
dualisant sur $\detalebccoefs X \Lambda$. Alors, pour toute immersion $j$ d'un
ouvert $U$ de $X$, le foncteur $\R j_\star$ envoie
$\detalebc U $ dans $\detalebc X$.
\end{proposition2}

Soit $K$ un complexe dualisant sur $\detalebc X$. Pour des raisons
évidentes, $j^\star K$ est un complexe dualisant sur $\detalebc U$.
On note $\Dual_X$ (resp. $\Dual_U$) les dualités induites par $K$
et $j^\star K$ sur les catégories triangulées $\detalebc X $ (resp.
$\detalebc U $). On a un isomorphisme canonique $\Dual_X\circ
j_!\simeq \R j_\star \Dual_U$. On en déduit que pour tout
$M\in\detalebc U $, $\R j_\star M\simeq \Dual_Xj_!\Dual_U M$,
ce qui permet de conclure que $\R j_\star M$ appartient à $\detalebc X$.

\begin{definition2}
Si $X$ est un schéma noethérien, $A$ une $\Lambda$-algèbre noethérienne
et $\mathscr P$
une partition galoisienne dirigeable de $X$, on note
${\detalecoefs X A}^{\mathscr P}$ la sous-catégorie triangulée de
$\detalecoefs X A$
dont les objets de cohomologie sont dans $\FCons^{\mathscr P}(X,A)$. On
définit de même les variantes ${\detalebcoefs X A}^{\mathscr P}$,
${\detalebccoefs X A}^{\mathscr P}$, etc.
\end{definition2}

\begin{proposition2}\label{17-proposition-r-j-star-preserve-d-b-partitionne}
Soit $j\colon U\to X$ une immersion ouverte entre schémas noethériens. On
suppose que $\R j_\star$ applique $\detalebc U$ dans $\detalebc X$. Pour
toute partition galoisienne dirigeable $\mathscr P$ de $U$, il existe une
partition galoisienne $\mathscr P'$ de $X$ et un entier $c$ tel que
$\R j_\star$ envoie ${\detaleb U}^{\mathscr P}$ dans
${\detaleb X}^{\mathscr P'}$ et que pour tout $q>c$ et
$\mathscr F\in\FCons^{\mathscr P}(U,\Lambda)$, on ait $\R^q j_\star
\mathscr F=0$.
\end{proposition2}

On sait que la catégorie $\FCons^{\mathscr P}(U,\Lambda)$ est localement
finie et admet même un nombre fini d'objets simples (et ceux-ci sont des
faisceaux constructibles). Comme $\R j_\star$ applique $\detalebc U$ dans
$\detalebc X$, on peut choisir un entier $c$ et une partition galoisienne
$\mathscr P'$ de $X$ tels que pour tout objet simple $\mathscr F$ de
$\FCons^{\mathscr P}(U,\Lambda)$, $\R j_\star \mathscr F$
appartienne à ${\detaleb X}^{\mathscr P'}$ et ait des objets de cohomologie
nuls en les degrés strictement supérieurs à $c$. Ce résultat s'étend par
dévissage aux objets noethériens de $\FCons^{\mathscr P}(U,\Lambda)$ puis à
cette catégorie toute entière du fait de la commutation des foncteurs
$\R^q j_\star$ aux limites inductives filtrantes.

\begin{corollaire2}\label{17-corollaire-i-bang-preserve-d-b-partitionne}
Soit $i\colon Z\to X$ une immersion ouverte entre schémas noethériens. On
suppose que $i^!$ applique $\detalebc X$ dans $\detalebc Z$. Pour toute partition
galoisienne dirigeable $\mathscr P$ sur $X$, il existe une partition
galoisienne $\mathscr P'$ sur $Z$ et un entier $c$ tel que
$i^!$ envoie ${\detaleb X}^{\mathscr P}$ dans ${\detaleb Z}^{\mathscr P'}$ et que
pour tout $q>c$ et $\mathscr F\in\FCons^{\mathscr P}(X,\Lambda)$, on ait
$\mathscr H^q(i^!\mathscr F)=0$.
\end{corollaire2}

Si on note $j\colon U \to X$ l'immersion ouverte complémentaire,
l'hypothèse sur $i^!$ énoncée ici équivaut à celle exigée sur
$\R j_\star$ dans la
proposition~\ref{17-proposition-r-j-star-preserve-d-b-partitionne}. Quitte à
raffiner $\mathscr P$, on peut supposer que les constituants de $\mathscr
P$ sont soit au-dessus de $U$, soit au-dessus de $Z$. On peut ainsi écrire
$\mathscr P=\mathscr P_U\cup \mathscr P_Z$ où $\mathscr P_U$ et $\mathscr
P_Z$ sont des partitions galoisiennes de $U$ et $Z$ respectivement. On
applique la proposition~\ref{17-proposition-r-j-star-preserve-d-b-partitionne}
à $\mathscr P_U$. On obtient une partition galoisienne $\mathscr P''$ de
$X$ telle que $\R j_\star $ applique ${\detaleb U}^{\mathscr
P_U}$ dans ${\detaleb X}^{\mathscr P''}$ et un entier naturel  $c'$
tel que pour tout
$\mathscr F\in\FCons^{\mathscr P_U}(U,\Lambda)$, on ait $\R^q j_\star
\mathscr F=0$  pour $q>c'$. Quitte à raffiner $\mathscr P''$, on peut
supposer que $\mathscr P''=\mathscr P''_Z\cup \mathscr P''_U$ comme
ci-dessus. Quitte à raffiner $\mathscr P''_Z$, on peut supposer que cette
partition galoisienne de $Z$ raffine $\mathscr P_Z$. En utilisant le
triangle distingué
\[i^!K\to i^\star K\to i^\star\R j_\star j^\star K\vers +\]
pour tout $K\in\detaleplus X$, on obtient aussitôt que $\mathscr P'=\mathscr
P''_Z$ et $c=c'+1$ conviennent.

\begin{proposition2}\label{17-proposition-r-j-star-preserve-d-b-ctf-coefficients}
Soit $j\colon U\to X$ une immersion ouverte entre schémas noethériens. On
suppose que $\R j_\star$ envoie $\detalebc U$ dans $\detalebc X$. Alors, pour
toute $\Lambda$-algèbre noethérienne $A$, $\R j_\star$ envoie
$\detalebccoefs U A$ (resp. $\detalebctfcoefs U A$)
dans $\detalebccoefs X A$ (resp. $\detalebctfcoefs X A$).
En outre, pour tout $Y\in\detalebctfcoefs U A$ et $M\in\dplus(A)$, le morphisme
canonique $M\overset{\L}\otimes_A \R j_\star Y\to \R
j_\star (M\overset{\L}\otimes_A Y)$ est un isomorphisme.
Par ailleurs, si $i\colon Z \to X$ est une immersion fermée complémentaire
à $j$, alors $i^!$ envoie $\detalebccoefs X A$ (resp. $\detalebctfcoefs X A$) dans
$\detalebccoefs Z A$ (resp. $\detalebctfcoefs Z A$), et pour tout
$M\in\dplus(A)$
et $Y\in\detalebccoefs X A$, le morphisme canonique
$M\overset{\L}\otimes_A i^!Y\to i^!(M\overset{\L}\otimes_A Y)$
est un isomorphisme.
\end{proposition2}

L'énoncé sur $i^!$ se déduit aussitôt de celui sur $\R j_\star$, on
se concentre donc sur celui-là. Pour montrer que $\R j_\star$ envoie
$\detalebccoefs U A$ dans $\detalebccoefs X A$, on peut supposer par un dévissage
évident que $\Lambda=\mathbf{Z}/\ell\mathbf{Z}$, où $\ell$ est un nombre
premier. Par conséquent, pour tout $\Lambda$-module $N$ et tout objet
$Y\in\detalebccoefs U \Lambda$, on a un isomorphisme $\R j_\star
(N\overset{\L}\otimes_\Lambda Y)\simeq
N\overset{\L}\otimes_\Lambda \R j_\star Y$ (cf.
proposition~\ref{17-proposition-coefficients-universels-r-j-star}). Si $N$ a
une structure de $A$-module de type fini, comme on sait que $\R
j_\star Y$ appartient à $\detalebc X$, on peut en déduire que $\R j_\star
(N\overset{\L}\otimes_\Lambda Y)$ appartient à $\detalebccoefs X A$.
D'après le dévissage de la
proposition~\ref{17-proposition-devissage-coefficients}, il vient que
$\R j_\star$ envoie $\detalebccoefs U A$ dans $\detalebccoefs X A$.

Montrons maintenant que $\R j_\star$ envoie $\detalebctfcoefs U A$ dans
$\detalebctfcoefs X A$. Compte tenu du résultat précédent, il suffit de montrer
que si $\mathscr F$ est un faisceau de $A$-modules plat et constructible
sur $U$, alors $\R j_\star \mathscr F\in\detalebtfcoefs X A$,
c'est-à-dire que $M\overset{\L}\otimes_A\R j_\star \mathscr F$
est borné indépendamment du $A$-module $M$. D'après la
proposition~\ref{17-proposition-coefficients-universels-r-j-star}, il suffit
de montrer que $\R j_\star(\mathscr F\otimes_A M)$ est borné
indépendamment du $A$-module $M$, et si tel est le cas, la formule des
coefficients universels énoncée ici sera satisfaite. Il existe une
partition galoisienne dirigeable $\mathscr P$ de $U$ tel que $\mathscr F$
appartienne à $\FCons^{\mathscr P}(U,A)$. Pour tout $A$-module $M$,
$\mathscr F\otimes_A M$ est un objet de $\FCons^{\mathscr P}(U,A)$, ainsi,
il suffit de montrer qu'il existe un entier $c$ tel que pour tout objet
$\mathscr G$ de $\FCons^{\mathscr P}(U,A)$, on ait $\R^q j_\star
\mathscr G=0$ pour $q>c$, ce qui résulte de la
proposition~\ref{17-proposition-r-j-star-preserve-d-b-partitionne}.

\begin{definition2}
Soit $X$ un schéma noethérien. Soit $A$ une $\Lambda$-algèbre noethérienne.
Soit $K\in\detalebctfcoefs X A$. On dira que $K$ vérifie la condition (B) si
pour toute partition galoisienne $\mathscr P$ de $X$, il existe un entier
$c$ et une partition galoisienne $\mathscr P'$ de $X$ tels que pour
tout $L\in{\detalebcoefs X A}^{\mathscr P}$, $\R\SheafHom_A(K,L)$
appartienne à ${\detalebcoefs X A}^{\mathscr P'}$, que si on suppose que
$\mathscr H^qL=0$ pour $q>0$, alors 
$\mathscr H^q\R\SheafHom_A(K,L)=0$ pour $q>c$ et enfin, que si $L$
appartient à ${\detalebccoefs X A}^{\mathscr P}$, alors
$\R\SheafHom_A(K,L)$ appartient à ${\detalebccoefs X A}^{\mathscr P'}$.
\end{definition2}

\begin{proposition2}\label{17-proposition-tout-objet-de-d-b-ctf-verifie-b}
Soit $X$ un schéma noethérien tel qu'il existe un complexe dualisant sur
$\detalebc X$. Alors, tout objet $K\in\detalebctfcoefs X A$ vérifie la condition
(B).
\end{proposition2}

\begin{lemme2}
Soit $X$ un schéma noethérien tel qu'il existe un complexe dualisant sur
$\detalebc X$. Soit $i\colon Z\to X$ une immersion fermée. Soit $j\colon U\to
X$ l'immersion ouverte complémentaire. Soit $K\in\detalebctfcoefs X A$. On
suppose que $i^\star K$ et $j^\star K$ satisfont la condition (B). Alors, $K$
satisfait la condition (B).
\end{lemme2}

Pour tout $L\in\detalecoefs X A$, on a un
triangle distingué dans $\detalecoefs X A$ :
\[i_\star \R\SheafHom_A(i^\star K,i^!
L)\to\R\SheafHom_A(K,L)\to\R j_\star
\R\SheafHom_A(j^\star K,j^\star L)\vers +\]
Grâce au résultat de l'exercice~\ref{17-exercice-partitions-ensemble-filtrant},
on peut combiner d'une part le résultat sur $i^!$ du
corollaire~\ref{17-corollaire-i-bang-preserve-d-b-partitionne} et la condition
(B) pour $i^\star K$ et d'autre part la 
proposition~\ref{17-proposition-r-j-star-preserve-d-b-partitionne} concernant
$\R j_\star$ et la condition (B) pour $j^\star K$ pour obtenir que
$K$ vérifie la condition (B).

\bigskip

Démontrons la 
proposition~\ref{17-proposition-tout-objet-de-d-b-ctf-verifie-b}. La condition (B) définit une sous-catégorie
triangulée de $\detalebctfcoefs X A$. Pour montrer la proposition, il suffit de
montrer que si $K$ est un faisceau de $A$-modules plat et constructibles,
alors $K$ satisfait la condition (B).
L'existence d'un complexe dualisant étant une
condition préservée par passage à un sous-schéma, le lemme précédent
fournit un moyen de dévisser la situation pour se ramener au cas où $K$ est
localement constant. On est ramené au lemme suivant :

\begin{lemme2}
Soit $X$ un schéma noethérien. Soit $\mathscr F$ un faisceau de $A$-modules
constructible, plat et localement constant. Alors, $\mathscr F$ satisfait
la propriété (B).
\end{lemme2}

Tout d'abord, pour tout $L\in\detalebcoefs X A$, si
$\mathscr H^qL=0$ pour $q>0$, alors pour tout $q>0$, $\mathscr
H^q(\R\SheafHom_A(\mathscr F,L))=0$, et si $L\in\detalebccoefs X A$, alors
$\R\SheafHom_A(K,L)\in\detalebccoefs X A$.
Il reste donc à montrer que si $\mathscr P$ est une partition galoisienne
de $X$, il existe une partition galoisienne $\mathscr P'$ de $X$ telle que
pour tout $L\in{\detalebcoefs X A}^{\mathscr P}$, alors
$\R\SheafHom_A(\mathscr F,L)$ appartient à
$L\in{\detalebcoefs X A}^{\mathscr P'}$. On peut supposer que $\mathscr P$ est
constitué d'un unique revêtement étale galoisien $X'\to X$. On choisit
un revêtement étale galoisien $X''\to X$ tel que l'image inverse de
$\mathscr F$ sur $X''$ soit un faisceau constant, puis un
revêtement galoisien $X'''\to X$ coiffant $X'$ et $X''$. On voit aussitôt
que la partition galoisienne $\mathscr P'=(X'''\to X)$ de $X$ convient.

\bigskip

Démontrons le théorème~\ref{17-theoreme-stabilite-d-b-ctf-par-rhom}. Soit
$K\in\detalebctfcoefs X A$. D'après la
proposition~\ref{17-proposition-tout-objet-de-d-b-ctf-verifie-b}, $K$ vérifie
la condition (B). Soit $L\in\detalebctfcoefs X A$. Il existe une partition
galoisienne $\mathscr P$ de $X$ telle que $L$ appartienne à 
${\detalebcoefs X A}^{\mathscr P}$. Pour tout $A$-module $M$, l'objet
$M\overset{\L}\otimes_A L$ appartient encore à cette catégorie (et
est borné indépendamment de $M$). Il résulte de la condition (B) de $K$ que
$\R\SheafHom_A(K,L)$ appartient à $\detalebccoefs X A$ et qu'il
existe une partition galoisienne $\mathscr P'$ telle que
$\R\SheafHom_A(K,M\overset{\L}\otimes_A L)$ soit un objet de
${\detalebcoefs X A}^{\mathscr P'}$ borné indépendamment du $A$-module $M$.
La proposition~\ref{17-proposition-coefficients-universels-1} permet de
déduire que $\R\SheafHom_A(K,L)$ appartient à $\detalebtfcoefs X A$ et
que pour tout $M\in\dplus(A)$, le morphisme canonique
$M\overset{\L}{\otimes}_A \R\SheafHom_A(K,L)\to
\R\SheafHom_A(K,M\overset{\L}{\otimes}_A L)$ est un
isomorphisme. On déduit aussitôt de cette formule la compatibilité au
changement d'anneau $A\to A'$ pour toute $A$-algèbre $A'$.

%>>>2
\subsection{Élimination d'hypothèses noethériennes sur \texorpdfstring{$A$}{A}}%<<<2

\begin{definition2}\label{17-definition-c-perfection}
Soit $A$ un anneau commutatif.
Soit $X$ un schéma noethérien. On dit d'un complexe
$K\in \detalecoefs X A$ qu'il est $\derivec$-parfait s'il existe une partition
finie $(U_i)_{i\in I}$ de $X$ par des sous-schémas (réduits) telle que pour
tout $i\in I$, $K_{|U_i}\in \detalecoefs {U_i} A$ soit un complexe parfait (cf.
\sga{6}{i}{4.8}). On note $\detalebcparfcoefs X A$ la sous-catégorie triangulée de
$\detalecoefs X A$ formée des complexes $\derivec$-parfaits.
\end{definition2}

Bien entendu, pour tout morphisme d'anneaux $A\to A'$, le foncteur
$A'\overset{\L}\otimes_A -\colon \detalecoefs X A\to \detalecoefs X {A'}$
induit un foncteur $\detalebcparfcoefs X A\to \detalebcparfcoefs X {A'}$.
En outre, si $A$ est un anneau noethérien,
$\detalebcparfcoefs X A=\detalebctfcoefs X A$.

\begin{theoreme2}\label{17-theoreme-dualite-d-b-c-parf-coefficients}
Soit $A$ une $\Lambda$-algèbre commutative. Soit $X$ un schéma
noethérien. S'il en existe, les
complexes dualisants sur $\detalebcparfcoefs X A$ sont uniques au produit
tensoriel près avec des objets inversibles.
Soit $K$ un complexe dualisant sur $\detalebccoefs X \Lambda$.
Alors, $A\overset{\L}\otimes_{\Lambda} K$ est un complexe dualisant
sur $\detalebcparfcoefs X A$.
En outre, le bifoncteur $\R\SheafHom_A$ préserve $\detalebcparfcoefs X A$
et commute à tout changement d'anneau $A\to A'$.
\end{theoreme2}

Le théorème~\ref{17-theoreme-stabilite-d-b-ctf-par-rhom} énonce
que si $B$ est une $\Lambda$-algèbre noethérienne, que $K$ et $L$
sont deux objets de $\detalebcparfcoefs X B$, alors pour toute $B$-algèbre
$A$, on a un isomorphisme canonique
\[\R\SheafHom_A(A\overset{\L}\otimes_B K,
A\overset{\L}\otimes_B
L)\simeq A\overset{\L}\otimes_B 
\R\SheafHom_B(K,L)\;\text{.}\]
Comme l'objet de droite appartient à $\detalebcparfcoefs X A$, il s'agit d'un
isomorphisme dans $\detalebcparfcoefs X A$.
Bref, compte tenu du
théorème~\ref{17-theoreme-stabilite-d-b-ctf-par-rhom} et
du théorème~\ref{17-theoreme-dualite-d-b-ctf-coefficients} (dont l'énoncé
d'unicité des complexes dualisants vaut aussi pour $\detalebcparfcoefs X A$ avec
la même démonstration), le théorème ci-dessus est ramené au lemme suivant :

\begin{lemme2}\label{17-lemme-approximation-noetherienne}
Soit $A$ un anneau commutatif.
Soit $X$ un schéma noethérien. Pour tout objet $K$ de
$\detalebcparfcoefs
X A$, il existe un sous-anneau noethérien (et même de type fini sur
$\mathbf{Z}$) $B$ de $A$ et $K'\in \detalebcparfcoefs X B$ tel que les objets
$K$ et $A\overset{\L}\otimes_B K'$ de $\detalebcparfcoefs X A$ soient
isomorphes.
\end{lemme2}

Comme il nécessite un examen plus attentif de la notion de
$\derivec$-perfection, on repousse la démonstration de ce lemme à la fin de
cette sous-section.

\begin{lemme2}
Soit $A$ un anneau commutatif. Soit $\mathscr F$ un faisceau de $A$-modules
sur $X_\et$. Les conditions suivantes sont équivalentes :
\begin{itemize}
\item[(i)] Il existe une partition $(U_i)_{i\in I}$ de $X$ par des
sous-schémas réduits tels que pour tout $i\in I$,
$\mathscr F_{|U_i}$ soit localement constant et que pour tout point
géométrique $\overline x$ de $X$, le $A$-module $\mathscr F_{\overline x}$
soit projectif de type fini ;
\item[(ii)] Le faisceau de $A$-modules $\mathscr F$ est
constructible\;\footnote{On utilise la définition donnée dans
\sga{4}{ix}{2.3} même si $A$ n'est pas noethérien, et non pas la définition
suggérée en note à cet endroit ; un faisceau constructible pour cette autre
définition est ce que pourrions appeler un faisceau
$\derivec$-pseudo-cohérent.} et pour tout point géométrique $\overline x$ de
$X$, le $A$-module $\mathscr F_{\overline x}$ est projectif  de type fini ;
\item[(iii)] Le faisceau de $A$-modules $\mathscr F$ est plat et
constructible.
\end{itemize}
\end{lemme2}

Par définition des faisceaux constructibles, on a évidemment l'équivalence
$\text{(i)}\Longleftrightarrow\text{(ii)}$. Si $\mathscr F$ est
constructible, les fibres $F_{\overline x}$ sont des $A$-modules de
présentation finie, il est alors équivalent d'exiger que ces modules soient
plats ou projectifs de type fini, ce qui montre l'équivalence
$\text{(ii)}\Longleftrightarrow\text{(iii)}$. 

On note $\mathscr C$ la catégorie fibrée au-dessus du site $X_\et$ qui à
$U\in X_\et$ fait correspondre la catégorie des faisceaux de $A$-modules
sur $U_\et$ et $\mathscr C_\derivec$ la sous-$X_\et$-catégorie de $\mathscr C$
formée des faisceaux de $A$-modules plats et constructibles. Nous allons
utiliser la terminologie de \sga{6}{i}{1.2}.
Il est évident qu'un objet de $\mathscr C$ qui est localement dans
$\mathscr C_\derivec$ est dans $\mathscr C_\derivec$ et que $\mathscr C_\derivec$ est stable par
noyau d'épimorphisme.
D'après \sga{4}{ix}{2.7}, un faisceau de $A$-modules sur $X_\et$ est
constructible si et seulement s'il est isomorphe au conoyau d'un morphisme
$A_V\to A_U$ pour $U$ et $V$ deux $X$-schémas étales et de présentation
finie. Les faisceaux $A_U$ pour $U$ étale et de présentation finis
sur $X$ sont donc évidemment plats et constructibles.
Il résulte de ces résultats qu'un objet de $\mathscr C_X$ est de
$\mathscr C_\derivec$-type fini si et seulement s'il est engendré par un nombre
fini de sections (ce qui revient à demander qu'il soit de $\mathscr
C_{0X}$-type fini) et qu'un objet de $\mathscr C_X$ est de $\mathscr
C_\derivec$-présentation finie si et seulement s'il est constructible (ce qui
revient encore à demander qu'il soit de $\mathscr C_{0X}$-présentation
finie). Il est par ailleurs évident que $\mathscr C_\derivec$ est quasi-relevable
dans $\mathscr C$ et même que $\mathscr C_{0X}$ est quasi-relevable dans
$\mathscr C_X$. Les catégories $\mathscr C_\derivec$ et $\mathscr C$ vérifient
donc les hypothèses de \sga{6}{i}{2.0} et de \sga{6}{ii}{1.1} (mais en général
pas de \sga{6}{i}{4.0}). On dispose donc d'une notion de complexe
pseudo-cohérent (relativement à $\mathscr C_\derivec)$ et celle-ci peut-être
définie de façon globale :

\begin{definition2}
Soit $A$ un anneau. Soit $X$ un schéma noethérien.
Soit $K\in C(X_\et,A)$. On dit que $K$ est strictement
$\derivec$-pseudo-cohérent
(resp. strictement $\derivec$-parfait) si
$K$ est un complexe borné supérieurement (resp. borné)
formé de faisceaux de $A$-modules plats et constructibles.
\end{definition2}

\begin{definition2}
Soit $A$ un anneau. Soit $X$ un schéma noethérien.
Soit $K\in \detalecoefs X A$. On dit que $K$ est
$\derivec$-pseudo-cohérent s'il est
isomorphe à l'image dans $\detalecoefs X A$ d'un complexe strictement
$\derivec$-pseudo-cohérent\;\footnote{En raison de l'existence de
\guil{résolutions globales}, la définition globale donnée ici
équivaut à la définition locale de \sga{6}{i}{2.3}.}.
Les objets $\derivec$-pseudo-cohérents forment une
sous-catégorie triangulée anonyme de $\detalecoefs X A$.
\end{definition2}

\begin{remarque2}
La $X_\et$-catégorie $\mathscr C$ contient aussi la $X_\et$-catégorie
$\mathscr C_0$ des faisceaux de $A$-modules facteurs directs de faisceaux
libres de $A$-modules. Bien entendu, $\mathscr C_0$ est contenue dans
$\mathscr C_{\derivec}$. Ainsi, les notions de stricte pseudo-cohérence (resp.
stricte perfection) définies relativement à $\mathscr C_0$ impliquent
les notions correspondantes relativement à $\mathscr C_\derivec$.
\end{remarque2}

\begin{proposition2}\label{17-proposition-critete-c-perfection}
Soit $A$ un anneau. Soit $X$ un schéma noethérien.
Soit $K\in \detalecoefs X A$. Alors, les conditions suivantes sont équivalentes
:
\begin{itemize}
\item[(i)] $K$ est $\derivec$-parfait ;
\item[(ii)] $K$ est $\derivec$-pseudo-cohérent et de tor-dimension finie ;
\item[(iii)] $K$ est isomorphe à l'image dans $\detalecoefs X A$ d'un complexe
strictement $\derivec$-parfait.
\end{itemize}
\end{proposition2}

\begin{lemme2}
Soit $A$ un anneau. Soit $X$ un schéma noethérien. Soit $(U_i)_{i\in I}$
une partition finie de $X$ par des sous-schémas réduits. Soit $K\in
\detalecoefs X
A$. On suppose que pour tout $i\in I$, la restriction de $K$ à $U_i$
est $\derivec$-pseudo-cohérente. Alors, $K$ est $\derivec$-pseudo-cohérent.
\end{lemme2}

Par les arguments habituels, on se ramène au cas où la partition de $X$ est
constituée d'un ouvert $U$ et d'un fermé $Z$. On note $i\colon Z\to X$ et
$j\colon U\to X$ les immersions correspondantes. On dispose d'un triangle
distingué dans $\detalecoefs X A$ :
\[j_!j^\star K\to K\to i_\star i^\star K\vers +\]
On sait que $j^\star K$ et $i^\star K$ sont $\derivec$-pseudo-cohérents.
Il est
évident que $j_!$ et $i_\star$ préservent la notion de faisceau plat et
constructible ; au niveau des catégories triangulées, ces foncteurs
préservent donc évidemment la notion de $\derivec$-pseudo-cohérence. Les objets
$j_!j^\star K$ et $i_\star i^\star K$ sont $\derivec$-pseudo-cohérent ;
il en résulte que $K$ aussi est $\derivec$-pseudo-cohérent.

\begin{lemme2}
Soit $A$ un anneau. Soit $X$ un schéma noethérien. Soit $K\in\detalecoefs X A$.
Si $K$ est $\derivec$-parfait, alors $K$ est $\derivec$-pseudo-cohérent.
\end{lemme2}

D'après le lemme précédent, quitte à passer à un recouvrement fini par des
localement fermés convenables, on peut supposer que $K$ est parfait.
Par conséquent, $K$ est pseudo-cohérent (sous-entendu relativement à la
$X_\et$-catégorie $\mathscr C_0$) ; \emph{a fortiori}, $K$ est
$\derivec$-pseudo-cohérent.

\bigskip

Démontrons la proposition~\ref{17-proposition-critete-c-perfection}.
L'implication $\text{(iii)}\Longrightarrow\text{(i)}$ est évidente.
L'implication $\text{(i)}\Longrightarrow\text{(ii)}$ résulte
essentiellement du lemme précédent ; il reste cependant à vérifier que si
$K$ est $\derivec$-parfait, alors il est de tor-dimension finie. Supposons
donc que $K$ est $\derivec$-parfait. Il existe un recouvrement fini
$(U_i)_{i\in I}$ de $X$ par des localement fermés tels que $K_{|U_i}$ soit
parfait pour tout $i\in I$. Pour obtenir (ii) pour $K$, il suffit de
montrer que le complexe parfait $K_{|U_i}$ est de tor-dimension finie, ce
qui résulte aussitôt de \sga{6}{i}{5.8.1}. Il reste à établir l'implication
$\text{(ii)}\Longrightarrow\text{(iii)}$, la plus intéressante pour nous.
Soit $K\in \detalecoefs X A$ un complexe $\derivec$-pseudo-cohérent
et de tor-dimension finie.
Par définition de la $\derivec$-pseudo-cohérence,
on peut remplacer si besoin est $K$ par un
complexe borné supérieurement formé de faisceaux de $A$-modules plats et
constructibles.
Pour tout $n\in\mathbf{Z}$, on peut considérer la troncature canonique
$\tau_{\leq n}K$ de $K$, si on note $Z^n$ le noyau de $K^n\to K^{n+1}$, il
s'agit du sous-complexe suivant de $K$ :
\[\dots \to K^{n-2} \to K^{n-1} \to Z^n \to 0\to \dots\;\text{.}\]
Comme $K$ est de tor-dimension finie, il existe un entier $a\in\mathbf{Z}$
tel que pour tout $A$-module $M$, $\tau_{\leq a}(K\otimes_A M)$ soit
acyclique. En appliquant ceci avec $M=A$, on obtient une résolution plate
de $Z^a$ :
\[\dots\to K^{a-3}\to K^{a-2}\to K^{a-1}\to Z^a \to 0\;\text{.}\]
Ensuite, on obtient aussitôt que pour un $A$-module $M$ quelconque, cette
suite reste exacte après passage au produit tensoriel avec $M$. Par suite,
pour tout $i>0$, $\SheafTor_i^A(Z^a,M)=0$ pour tout $A$-module $M$,
ce qui implique que $Z^a$ est un faisceau de $A$-modules plat. Par
ailleurs, $Z^a$ est le conoyau du morphisme $K^{a-2}\to K^{a-1}$, donc
$Z^a$ est un faisceau constructible. Ainsi, $Z^a$ est plat et constructible.
Le complexe $K$ est quasi-isomorphe au complexe
strictement $\derivec$-parfait
\[\dots \to 0\to Z^a \to K^a \to K^{a+1} \to \dots\;\text{,}\]
ainsi $K$ vérifie la condition (iii).

\bigskip

Nous sommes maintenant en mesure de démontrer le
lemme~\ref{17-lemme-approximation-noetherienne}. Compte tenu de la
proposition~\ref{17-proposition-critete-c-perfection}, il s'agit de montrer
que si $K\in C(X_\et,A)$ est un complexe strictement $\derivec$-parfait,
alors $A$ contient un sous-anneau $A'$ de type fini sur $\mathbf{Z}$ tel
qu'il existe un complexe strictement $\derivec$-parfait $K'\in C(X_et,A')$
et un isomorphisme $K\simeq A\otimes_{A'} K'$. Ceci résulte aussitôt du
lemme suivant :

\begin{lemme2}
Soit $(A_\alpha)_{\alpha\in I}$ un système inductif d'anneaux commutatifs
indexé par un ensemble ordonné filtrant $I$. On en note $A$ la limite
inductive. Soit $X$ un schéma noethérien. Alors, la donnée d'un faisceau de
$A$-modules constructible (resp. plat et constructible) sur $X$ équivaut à
la donnée d'un faisceau de $A_\alpha$-modules constructible (resp. plat et
constructible) sur $X$ pour $\alpha$ assez grand.
\end{lemme2}

Conformément aux grands principes de \ega{iv}{8},
ceci signifie d'une part que si $F$ est un faisceau de $A$-modules
constructible sur $X$, il existe $\alpha\in I$ et $F_\alpha$ un faisceau de
$A_\alpha$-modules constructible sur $X$ tel que $F$ soit isomorphe à
$A\otimes_{A_\alpha} F_\alpha$ et d'autre part que si $\alpha\in I$ et que
$F_\alpha$ et $G_\alpha$ sont deux faisceaux de $A_\alpha$-modules
constructibles sur $X$, si on note
$F_\beta=A_\beta\otimes_{A_\alpha}F_\alpha$ et
$G_\beta=A_\beta\otimes_{A_\alpha}G_\alpha$ pour tout $\beta\geq \alpha$ et
$F=A\otimes_{A_\alpha}F_\alpha$ et $G=A\otimes_{A_\alpha}G_\alpha$, alors
l'application canonique
\[\underset{\beta\geq \alpha}\colim \Hom_{A_\beta}(F_\beta,G_\beta)\to
\Hom_A(F,G)\]
est un isomorphisme. En outre, si $\alpha\in I$ et que $F_\alpha$ est un
faisceau de $A_\alpha$-modules constructible sur $X$, alors
$F=A\otimes_{A_\alpha}F_\alpha$ est plat si et seulement si pour $\beta\geq
\alpha$ assez grand, $F_\beta=A_{\beta}\otimes_{A_\alpha} F_\alpha$ est
plat.

L'énoncé dans le cas non respé résulte de la description des faisceaux de
$B$-modules constructibles (pour tout anneau commutatif $B$) comme conoyau
d'une flèche $B_V\to B_U$ où $U$ et $V$ sont étales et de présentation
finie sur $X$, et du fait que le foncteur $\HH^0(V,-)$ de la catégorie des
faisceaux de groupes abéliens sur $X_\et$ vers celle des groupes abéliens
commute aux limites inductives filtrantes \sga{4}{vii}{3.3}.

Il reste à montrer que si $F_\alpha$ est un faisceau de $A_\alpha$-modules
constructible tel que, avec les notations ci-dessus, $F$ soit $A$-plat,
alors pour $\beta\geq \alpha$ assez grand, $F_\beta$ est $A_\beta$-plat. En
utilisant une décomposition de $X$ en réunion de localement fermés connexes
au-dessus desquels $F_\alpha$ soit localement constant trivialisé par un
revêtement étale, on peut supposer que $X$ est connexe et que
$F_\alpha$ est localement constant et trivialisé par un revêtement étale
$Y\to X$. Pour vérifier la platitude de $F_\beta$, il suffit de l'obtenir
pour \emph{une} fibre $(F_\beta)_{\overline x}$ ; on est ainsi ramené au
lemme suivant :

\begin{lemme2}
Soit $(A_\alpha)_{\alpha\in I}$ un système inductif d'anneaux commutatifs
indexé par un ensemble ordonné filtrant $I$. On en note $A$ la limite
inductive. Soit $X$ un schéma noethérien. Soit $\alpha\in I$, soit
$M_\alpha$ un $A_\alpha$-module de présentation finie. On suppose que
$M=A\otimes_{A_\alpha} M_\alpha$ est $A$-plat. Alors, il existe $\beta\geq
\alpha$ tel que $M_\beta=A_\beta\otimes_{A_\alpha} M_\alpha$ soit
$A_\alpha$-plat.
\end{lemme2}

Les modules considérés étant de présentation finie, le module $M$ (resp.
$M_\beta$) est plat si et seulement s'il est facteur direct d'un module
libre de type fini. On peut conclure en utilisant convenablement
\ega{iv}{8.5.2}.

%>>>2
%>>>1

\ifx\danslelivre\undefined
\appendix
\fi


\section{Produits tensoriels de complexes non bornés}%<<<1
\subsection{\texorpdfstring{$K$-platitude}{K-platitude}}%<<<2

\begin{definition2}
Soit $(\mathscr T,\mathscr A)$ un topos annelé en anneaux commutatifs. On
note $C(\mathscr T,\mathscr A)$ (resp. $K(\mathscr T,\mathscr A)$,
$\dtotale(\mathscr T,\mathscr A)$) la catégorie des complexes de $\mathscr
A$-Modules (resp. la catégorie homotopique correspondance, la catégorie
dérivée associée). On dispose d'un bifoncteur $\otimes_{\mathscr A}$ sur
$C(\mathscr T,\mathscr A)$ induisant un bifoncteur sur $K(\mathscr
T,\mathscr A)$\;\footnote{Plus précisément, si $K$ et $L$ sont deux
complexes de $\mathscr A$-Modules, on peut définir un bicomplexe dont la
composante de bidegré $(p,q)$ est $K^p\otimes_{\mathscr A}L^q$ et alors
$K\otimes_{\mathscr A} L$ est le complexe simple (défini en termes
de sommes) associé à ce bicomplexe. Nous n'imposerons pas ici de
convention de signes au lecteur.}.
Soit $K\in C(\mathscr T,\mathscr A)$. On dit que $K$ est
$K$-plat si le foncteur triangulé $K\otimes_{\mathscr A}-\colon K(\mathscr
T,\mathscr A)\to K(\mathscr T,\mathscr A)$ préserve les
quasi-isomorphismes, autrement dit que pour tout complexe acyclique $L$ de
$C(\mathscr T,\mathscr A)$, le complexe $K\otimes_{\mathscr A} L$ est
acyclique (cf. \cite[definition~5.1]{RUC@Spaltenstein}).
\end{definition2}

La sous-catégorie pleine de $C(\mathscr T,\mathscr A)$ formée des complexes
$K$-plats est stable par limites inductives filtrantes, sommes directes et
facteurs directs. En outre, la sous-catégorie pleine de $K(\mathscr
T,\mathscr A)$ correspondante est une sous-catégorie triangulée de
$K(\mathscr T,\mathscr A)$.

\begin{proposition2}
Soit $(\mathscr T,\mathscr A)$ un topos annelé en anneaux commutatifs. Soit
$K$ un complexe borné supérieurement formé de $\mathscr A$-modules plats
(cf. \sga{4}{v}{1}). Alors, $K$ est $K$-plat.
\end{proposition2}

En utilisant les foncteurs de troncature bête et la stabilité par limites
inductives filtrantes de la $K$-platitude, on se ramène au cas où $K$ est
borné. Comme la $K$-platitude définit une
sous-catégorie triangulée de $K(\mathscr T,\mathscr A)$, on peut procéder à
un dévissage utilisant encore les troncatures bêtes pour se ramener au cas
où $K^q=0$ pour $q\neq 0$. On est alors ramené à montrer que si $F$ est un
$\mathscr A$-Module plat et $L\in C(\mathscr T,\mathscr A)$ un complexe
acyclique, alors $F\otimes_{\mathscr A}L$ est acyclique, ce qui résulte
aussitôt de la définition de la platitude.

%>>>2
\subsection{Résolutions \texorpdfstring{$K$-plates}{K-plates}}%<<<2
\subsubsection{Définition du produit tensoriel dérivé}%<<<3

\begin{theoreme3}\label{17-theoreme-existence-resolutions-k-plates}
Soit $(\mathscr T,\mathscr A)$ un topos annelé en anneaux commutatifs. Il
existe un foncteur $\rho\colon C(\mathscr T,\mathscr A)\to C(\mathscr
T,\mathscr A)$ et une transformation naturelle $\rho K\to K$ pour $K\in
C(\mathscr T,\mathscr A)$ telle que :
\begin{itemize}
\item pour tout $K\in C(\mathscr T,\mathscr A)$, $\rho K$ soit $K$-plat ;
\item pour tout $K\in C(\mathscr T,\mathscr A)$, le morphisme $\rho K\to K$
soit un quasi-isomorphisme ;
\item le foncteur $\rho$ commute aux limites inductives filtrantes.
\end{itemize}
\end{theoreme3}

Ce théorème sera démontré plus bas. Déduisons-en aussitôt la proposition
triviale suivante, qui constitue notre définition du produit tensoriel sur
$\dtotale(\mathscr T,\mathscr A)$ :

\begin{proposition3}\label{17-proposition-existence-produit-tensoriel-derive}
Soit $(\mathscr T,\mathscr A)$ un topos annelé en anneaux commutatifs.
Le foncteur dérivé total à gauche de $\otimes_{\mathscr A}\colon C(\mathscr
T,\mathscr A)\times C(\mathscr T,\mathscr A)\to C(\mathscr T,\mathscr A)$
existe. Plus précisément, pour tous $K$ et $L$ dans $C(\mathscr T,\mathscr
A)$, on note $K\overset{\L}\otimes_{\mathscr A} L=(\rho
K)\otimes_{\mathscr A}(\rho L)\in C(\mathscr T,\mathscr A)$ ; ce bifoncteur
$\overset{\L}\otimes_{\mathscr A}$ commute aux limites inductives
filtrantes en chaque argument et, préservant les quasi-isomorphismes, il
induit un bifoncteur du même nom $\dtotale(\mathscr T,\mathscr A)\times
\dtotale(\mathscr
T,\mathscr A)\to \dtotale(\mathscr T,\mathscr A)$ ; la transformation naturelle
évidente $\overset{\L}\otimes_{\mathscr A}\to\otimes_{\mathscr A}$ 
fait de $\overset{\L}\otimes_{\mathscr A}$ le foncteur dérivé total
à gauche de $\otimes_{\mathscr A}$
(cf.~\cite[remark~7.4, Chapter~~II]{SHT@Goerss-Jardine} pour une définition
des foncteurs dérivés totaux en termes d'extensions de Kan). En outre,
si $K$ et $L$ sont deux objets de $C(\mathscr T,\mathscr A)$ dont l'un
au moins est $K$-plat, alors le morphisme canonique
$K\overset{\L}\otimes_{\mathscr A}L\to K\otimes_{\mathscr A}L$ est un
isomorphisme dans $\dtotale(\mathscr T,\mathscr A)$.
\end{proposition3}

%>>>3
\subsubsection{Modules sur un anneau}%<<<3

On se place ici dans le cas particulier où le topos $\mathscr T$ est
ponctuel. On peut identifier les faisceaux de $\mathscr A$-modules à des
$A$-modules pour un anneau $A$. Le lemme suivant
démontre le théorème~\ref{17-theoreme-existence-resolutions-k-plates} dans ce
cas particulier.

\begin{lemme3}\label{17-lemme-existence-resolutions-k-plates-topos-ponctuel}
Pour tout anneau commutatif $A$, on peut définir un foncteur $\rho_A$ et
une transformation naturelle $\rho_A\to\Id$ de foncteurs de la
catégorie $C(A)$ des
complexes de $A$-modules dans elle-même telle que $\rho_A$ commute aux
limites inductives filtrantes, préserve les monomorphismes,
que pour tout $K\in C(A)$, le morphisme de
complexes $\rho_A(K)\to K$ soit un quasi-isomorphisme, que pour tout
entier relatif $n$, $\rho_A(K)^n$ soit un $A$-module libre, et que
$\rho_A(K)$ soit la limite inductive filtrante de ses sous-complexes
bornés formés de $A$-modules libres (en particulier, $\rho_A(K)$ est
$K$-plat).

On peut définir de telles résolutions $K$-plates $\rho_A$ pour tout anneau
commutatif $A$ de sorte que si $A\to A'$ est un morphisme d'anneaux, on ait
un morphisme fonctoriel $\rho_A(K)\to\rho_{A'}(K)$ dans $C(A)$ pour
$K\in C(A')$, ce morphisme vérifiant une compatibilité évidente
à la composition des morphismes d'anneaux.
\end{lemme3}

On note $G$ le foncteur adjoint à gauche du foncteur d'oubli $\mathrm{oub}$
de la catégorie des $A$-modules vers celle des ensembles pointés. On pose
$F=G\circ \mathrm{oub}$. Si $M$ est un $A$-module, $FM$ est le quotient du
$A$-module libre de base l'ensemble $M$ par le sous-module libre de rang
$1$ engendré par le zéro de $M$. Le morphisme d'adjonction $FM\to M$ est un
épimorphisme.
Pour tout
$A$-module $M$, on note $Z_0$ le noyau de cet épimorphisme $FM\to M$
et on pose $(F'M)_0=FM$.
Ensuite, de façon évidente, pour tout entier naturel $n\geq 1$,
on peut définir par récurrence un objet $(F'M)_n=FZ_{n-1}$,
un morphisme $d_n\colon (F'M)_n\to (F'M)_{n-1}$ et le noyau $Z_n=\ker d_n$.
Il est évident que l'on définit ainsi un complexe $F'M$ concentré en degrés
négatifs ou nuls, muni d'une augmentation $F'M\to M$ qui soit un
quasi-isomorphisme. Comme $F$ préserve les monomorphismes, on voit que $F'$
préserve aussi les monomorphismes.

Soit $K\in C(A)$. Le foncteur $F'$ défini ci-dessus n'est pas additif (à
moins que $A=0$), mais il est tel que $F'(0)=0$. Ainsi, si on applique
terme à terme le foncteur $F'$ aux objets $K_n$ pour tout $n\in\mathbf{Z}$,
on obtient un complexe double dans la catégorie des $A$-modules. On note
$\rho_A(K)$ le complexe simple associé (défini en termes de sommes). On
dispose bien entendu d'un morphisme d'augmentation $\rho_A(K)\to K$. Il
est évident que $\rho_A$ commute aux limites inductives filtrantes,
préserve les monomorphismes et que pour tout entier $n$, $\rho_A(K)^n$
soit une $A$-module libre. Si $K$ est borné supérieurement, le fait que
pour tout $n\in\mathbf{Z}$, le morphisme $\rho_A(K^n)\to K^n$ soit un
quasi-isomorphisme implique, par passage au complexe simple, que
$\rho_A(K)\to K$ est un quasi-isomorphisme. Comme tout complexe de
$A$-module peut s'écrire comme une limite inductive filtrante de
sous-complexes bornés supérieurement, il vient que pour tout $K\in C(A)$,
le morphisme $\rho_A(K)\to K$ est un quasi-isomorphisme. Comme
$\rho_A$ préserve les monomorphismes et commute aux limites inductives
filtrantes, pour montrer que $\rho_A(K)$ est une limite inductive
filtrante de ses sous-complexes bornés formés de $A$-modules libres, on
peut supposer que $K$ est borné supérieurement : le résultat est alors
trivial puisque, dans ce cas, $\rho_A(K)$ est un complexe borné
supérieurement formé de $A$-modules libres.

La dernière assertion concernant le changement d'anneau étant évidente, on
peut considérer que le lemme a été démontré.

%>>>3
\subsubsection{Préfaisceaux de Modules}%<<<3
\label{17-subsubsection-produit-tensoriel-derive-prefaisceaux}

On suppose maintenant que $\mathscr T$ est le topos des préfaisceaux sur
une petite catégorie $\mathscr C$. Le faisceau d'anneaux $\mathscr A$ est
un préfaisceau d'anneaux commutatifs sur $\mathscr C$.

Soit $K\in C(\mathscr T,\mathscr A)$. Pour tout objet $U$ de $\mathscr C$,
$K(U)$ s'identifie à un objet de $C(\mathscr A(U))$. On applique la
construction du
lemme~\ref{17-lemme-existence-resolutions-k-plates-topos-ponctuel} à l'anneau
$\mathscr A(U)$. On pose $(\rho K)(U)=\rho_{\mathscr A(U)}(K(U))\in
C(\mathscr A(U))$. Si $V\to U$ est un morphisme dans $\mathscr C$, on
définit un $\mathscr A(U)$-morphisme $(\rho K)(U)\to (\rho K)(V)$ de la
façon suivante :
\[\rho_{\mathscr A(U)}(K(U))\to \rho_{\mathscr A(U)}(K(V))\to
\rho_{\mathscr A(V)}(K(V))\]
où le morphisme de gauche est induit par la structure de complexes de 
préfaisceaux de $\mathscr A$-modules sur $K$ et la flèche de droite par la
compatibilité de la construction du
lemme~\ref{17-lemme-existence-resolutions-k-plates-topos-ponctuel} au
changement d'anneau. D'après ce lemme, ces morphismes de transition
définissent une structure de préfaisceau sur $\rho K$. Ainsi, on a défini
un objet $\rho K\in C(\mathscr T,\mathscr A)$ et il est muni d'un morphisme
fonctoriel $\rho K \to K$.

Par construction, on peut vérifier les vertus présumées de $\rho$ terme à
terme ; ainsi, ce foncteur $\rho$ permet d'établir le
théorème~\ref{17-theoreme-existence-resolutions-k-plates} dans le cas où le
topos est un topos de préfaisceaux.

%>>>3
\subsubsection{Faisceaux de Modules}%<<<3
\label{17-subsubsection-produit-tensoriel-derive-faisceaux}

\begin{proposition3}\label{17-proposition-stabilite-k-platitude-par-faisceau-associe}
Soit $\mathscr T$ le topos des faisceaux sur un site dont la catégorie
sous-jacente est notée $\mathscr C$. On note $\mathscr T'$ le topos des
préfaisceaux d'ensembles sur $\mathscr C$. Soit $\mathscr A'$ un préfaisceau
d'anneaux commutatifs sur $\mathscr C$. On note $\mathscr A$ le faisceau
d'anneaux $a\mathscr A'$ sur $\mathscr T$ associé à $\mathscr A'$.
Si $K\in C(\mathscr
T',\mathscr A')$ est $K$-plat, alors $aK\in C(\mathscr T,\mathscr A)$
est $K$-plat.
\end{proposition3}

\begin{lemme3}
Avec les notations de la
proposition~\ref{17-proposition-stabilite-k-platitude-par-faisceau-associe},
si $K$ et $L$ sont des objets de $C(\mathscr T',\mathscr A')$ tel que $aK$
soit nul dans $\dtotale(\mathscr T,\mathscr A)$, alors
$a(K\overset{\L}\otimes_{\mathscr A'} L)$
(où le produit tensoriel dérivé au-dessus de
$\mathscr A'$ est celui défini plus haut dans le cas des préfaisceaux) est
nul dans $\dtotale(\mathscr T,\mathscr A)$.
\end{lemme3}

La compatibilité du produit tensoriel dérivé sur $\dtotale(\mathscr T',\mathscr
A')$ avec les limites inductives filtrantes calculées dans $C(\mathscr
T',\mathscr A')$ permet de supposer que $L$ est un complexe borné
supérieurement (utiliser les troncatures canoniques). Quitte à remplacer
$L$ par la résolution $K$-plate $\rho_{\mathscr A'}L$ sus-définie, on peut
ensuite supposer que $L$ est borné supérieurement et constitué de $\mathscr
A'$-Modules plats. En utilisant les troncatures bêtes sur ce complexe $L$,
on peut supposer de plus que $L$ est borné. Finalement, en utilisant les
triangles distingués dans $K(\mathscr T',\mathscr A')$ donnés par les
troncatures bêtes, on peut finalement supposer que $L$ est formé d'un
unique $\mathscr A'$-Module plat placé en degré $0$.

Bref, il faut montrer que si $L$ est un $\mathscr A'$-Module plat et que
$K\in C(\mathscr T',\mathscr A')$ est tel que $aK$ soit un complexe
acyclique dans $C(\mathscr T,\mathscr A)$, alors le complexe de faisceaux 
$a(K\otimes_{\mathscr A'} L)$ est acyclique. On a un isomorphisme
$a(K\otimes_{\mathscr A'} L)\simeq aK\otimes_{\mathscr A}aL$, donc pour
conclure que $a(K\otimes_{\mathscr A'} L)$ est acyclique, il suffit de
montrer que le $\mathscr A$-Module $aL$ est plat, ce qui est vrai d'après
\sga{4}{v}{1.7.1}.

\bigskip

Montrons la
proposition~\ref{17-proposition-stabilite-k-platitude-par-faisceau-associe}.
Notons $\mathscr Z$ la sous-catégorie triangulée de $\dtotale(\mathscr T',\mathscr
A')$ formée des complexes qui sont annulés par le foncteur faisceau associé
$\dtotale(\mathscr T',\mathscr A')\to \dtotale(\mathscr T,\mathscr A)$. Le foncteur
induit $\dtotale(\mathscr T',\mathscr A')/\mathscr Z\to \dtotale(\mathscr T,\mathscr A)$
est évidemment une équivalence de catégories triangulées. D'après le lemme,
le bifoncteur $\overset{\L}\otimes_{\mathscr A'}$ sur $\dtotale(\mathscr
T',\mathscr A')$ passe au quotient par $\mathscr Z$ pour définit un
bifoncteur sur $\dtotale(\mathscr T,\mathscr A)$. La proposition en résulte
aussitôt. En effet, soit $K\in C(\mathscr T',\mathscr A')$ $K$-plat, soit
$L\in C(\mathscr T,\mathscr A)$ tel que $L$ soit nul dans $\dtotale(\mathscr
T,\mathscr A)$. On peut identifier $L$ à un objet $L'$ de $C(\mathscr
T',\mathscr A')$ et cet objet $L'$ appartient à la sous-catégorie
triangulée $\mathscr Z$ de $\dtotale(\mathscr T',\mathscr A')$. Le lemme montre
que $K\overset{\L}\otimes_{\mathscr A'} L'$ appartient à $\mathscr
Z$, autrement dit, $K$ étant $K$-plat, que $K\otimes_{\mathscr A'} L'$
appartient à $\mathscr Z$, c'est-à-dire que le complexe 
de faisceaux $aK\otimes_{\mathscr A} L$ est acyclique, ce qui montre que
$aK\in C(\mathscr T,\mathscr A)$ est $K$-plat.

\begin{corollaire3}
Avec les notations de la
proposition~\ref{17-proposition-stabilite-k-platitude-par-faisceau-associe},
le foncteur $\rho_{\mathscr A}$ qui à un complexe $K\in C(\mathscr
T,\mathscr A)$ associe $a\rho_{\mathscr A'} K'$ où $K'$ est le $\mathscr
A'$-Module défini par $K$ est un foncteur de résolution $K$-plate sur
$C(\mathscr T,\mathscr A$) vérifiant les conditions du
théorème~\ref{17-theoreme-existence-resolutions-k-plates}.
Ainsi, on dispose d'un bifoncteur $\overset{\L}\otimes_{\mathscr  A}$
sur $\dtotale(\mathscr T,\mathscr A)$ et d'un isomorphisme bifonctoriel
\[aK\overset{\L}\otimes_{\mathscr A}aL\simeq
a(K\overset{\L}\otimes_{\mathscr A'}L)\]
dans $\dtotale(\mathscr T,\mathscr A)$ pour tous $K$ et $L$ dans $\dtotale(\mathscr
T',\mathscr A')$.
\end{corollaire3}

Avec l'énoncé de ce corollaire s'achève la démonstration du
théorème~\ref{17-theoreme-existence-resolutions-k-plates}.

%>>>3
%>>>2
\subsection{Compléments}%<<<2
\subsubsection{Homomorphismes internes}%<<<3

\begin{definition3}
Soit $(\mathscr T,\mathscr A)$ un topos annelé en anneaux commutatifs. On
note $\SheafHom_{\mathscr A}$ le foncteur adjoint à droite du foncteur
produit tensoriel $\otimes_{\mathscr A}$ sur $C(\mathscr T,\mathscr A)$,
c'est-à-dire que pour $X$, $Y$ et $Z$ des objets de $C(\mathscr T,\mathscr
A)$, on a un isomorphisme canonique de groupes abéliens :
\[\Hom_{C(\mathscr T,\mathscr A)}(X\otimes_{\mathscr A} Y,Z)\simeq 
\Hom_{C(\mathscr T,\mathscr A)}(X,\SheafHom_{\mathscr A}(Y,Z))\;\text{.}\]
\end{definition3}

On rappelle qu'un objet $L\in C(\mathscr T,\mathscr A)$ est $K$-injectif si
pour tout complexe acyclique $K\in C(\mathscr T,\mathscr A)$, le complexe
de groupes abéliens $\Hom^\bullet_{\mathscr A}(K,L)$ (complexe simple
associé au complexe double $(\Hom_{\mathscr A}(K^{-p},L^q))_{p,q}$) est
acyclique et qu'il existe des foncteurs de résolution
$K$-injectives\;\footnote{Ce résultat est énoncé dans
\cite[§4]{RUC@Spaltenstein} dans le cadre des espaces topologiques annelés,
mais la démonstration peut être étendue au cas des topos annelés. Le
principe de la démonstration est similaire à celui utilisé par Grothendieck
pour montrer l'existence de suffisamment d'injectifs dans les catégories de
faisceaux dans \cite[§I.10]{Tohoku@Grothendieck} ; dans le contexte de l'algèbre
homotopique, cet argument est connu sous le nom de \guil{raisonnement du
petit objet}. On peut trouver une démonstration pour le cas qui nous
intéresse dans \cite{Complexes@Hovey}.}.

La proposition suivante, indiquée pour mémoire, est essentiellement
triviale :

\begin{proposition3}
Soit $(\mathscr T,\mathscr A)$ un topos annelé en anneaux commutatifs.
Le foncteur $\SheafHom_{\mathscr A}$ sur $C(\mathscr T,\mathscr A)$ admet
un foncteur dérivé total à droite $\R\SheafHom_{\mathscr A}\colon
\dtotale(\mathscr T,\mathscr A)^\opp
\times\dtotale(\mathscr T,\mathscr A)\to \dtotale(\mathscr T,\mathscr
A)$, adjoint à droite de $\overset{\L}\otimes_{\mathscr A}$. En
outre, si $K$ et $L$ sont des objets de $C(\mathscr T,\mathscr A)$, avec
$L$ $K$-injectif, alors le morphisme canonique 
\[\SheafHom_{\mathscr A}(K,L)\to\R\SheafHom_{\mathscr A}(K,L)\]
est un isomorphisme dans $\dtotale(\mathscr T,\mathscr A)$ ; si on suppose
de plus que $K$ est $K$-plat, alors $\SheafHom_{\mathscr A}(K,L)$ est
$K$-injectif. Enfin, si $K$, $L$ et $M$ sont trois objets de
$\dtotale(\mathscr T,\mathscr A)$, on a un isomorphisme fonctoriel dans
$\dtotale(\mathscr T,\mathscr A)$ \guil{cher à Cartan} :
\[\R\SheafHom_{\mathscr A}(K\overset{\L}\otimes_{\mathscr A} L,M)
\simeq \R\SheafHom_{\mathscr A}(K,
\R\SheafHom_{\mathscr A}(L,M))\;\text{.}\]
\end{proposition3}

%>>>3
\subsubsection{Compatibilité aux images inverses}%<<<3

\begin{definition3}
Soit $u\colon (\mathscr T,\mathscr A)\to (\mathscr T',\mathscr A')$ un
morphisme de topos annelés en anneaux commutatifs. Le foncteur
$u^\star\colon C(\mathscr T,\mathscr A)\to C(\mathscr T',\mathscr A')$
admet un foncteur dérivé total à gauche que l'on note $\L u^\star$
(construit en appliquant $u^\star$ à une résolution $K$-plate).
\end{definition3}

\begin{proposition3}\label{17-proposition-u-star-preserve-k-platitude}
Soit $u\colon (\mathscr T,\mathscr A)\to (\mathscr T',\mathscr A')$ un
morphisme de topos annelés en anneaux commutatifs. Alors, pour tout $K\in
C(\mathscr T',\mathscr A')$ $K$-plat, $u^\star K\in C(\mathscr T,\mathscr
A)$ est $K$-plat.
\end{proposition3}

Pour démontrer cette proposition, il est commode d'introduire une nouvelle
définition :

\begin{definition3}
Soit $(\mathscr T,\mathscr A)$ un topos annelé en anneaux commutatifs.
Soit $K\in C(\mathscr T,\mathscr A)$. On dit que $K$ est
universellement acyclique si pour tout $\mathscr A$-module $L$, le complexe
$K\otimes_{\mathscr A} L$ est acyclique.
\end{definition3}

\begin{lemme3}
Soit $(\mathscr T,\mathscr A)$ un topos annelé en anneaux commutatifs.
Soit $K\in C(\mathscr T,\mathscr A)$. Les conditions suivantes sont
équivalentes :
\begin{itemize}
\item[(i)] $K$ est universellement acyclique ;
\item[(i')] pour tout $L\in C(\mathscr T,\mathscr A)$, le complexe
$K\otimes_{\mathscr A} L$ est acyclique ;
\item[(ii)] $K$ est $K$-plat et acyclique.
\end{itemize}
\end{lemme3}

L'équivalence entre (i) et (i') est facile : supposant (i), on montre
d'abord que la condition énoncée dans (i') est vérifiée par les complexes
$L$ bornés (dévissage), puis par tous les complexes (passage à la limite
inductive utilisant les troncatures bêtes et canoniques). L'implication
$\text{(i)}\Longrightarrow\text{(ii)}$ est triviale. L'implication
$\text{(ii)}\Longrightarrow\text{(i)}$ résulte aussitôt de la
proposition~\ref{17-proposition-existence-produit-tensoriel-derive}.

\medskip

De ce lemme, on déduit le suivant :

\begin{lemme3}
Soit $(\mathscr T,\mathscr A)$ un topos annelé en anneaux commutatifs.
Soit $K\in C(\mathscr T,\mathscr A)$. Si $K'\to K$ est un
quasi-isomorphisme, avec $K'$ $K$-plat, alors $K$ est $K$-plat si et
seulement si le cône de $K'\to K$ est universellement acyclique.
\end{lemme3}

\begin{lemme3}
Soit $u\colon (\mathscr T,\mathscr A)\to (\mathscr T',\mathscr A')$ un
morphisme de topos annelés en anneaux commutatifs.
Si on note $\rho_{\mathscr A'}$ \guil{le}
foncteur de résolution $K$-plate construit dans le 
paragraphe~\ref{17-subsubsection-produit-tensoriel-derive-faisceaux}, alors
pour tout $K\in C(\mathscr T',\mathscr A')$, l'objet $u^\star\rho_{\mathscr
A'}K\in C(\mathscr T,\mathscr A)$ est $K$-plat.
\end{lemme3}

Compte tenu du procédé de construction de $\rho_{\mathscr A'}$ par
extension aux complexes de faisceaux d'une construction pour les faisceaux
par passage au complexe simple, on peut supposer que $K$ est constitué
d'une unique faisceau placé en degré $0$. Il s'agit de montrer que pour
tout $\mathscr A'$-Modules $F$, le complexe (borné supérieurement)
$u^\star \rho_{\mathscr A'}F$ est constitué de $\mathscr A$-Modules plats.
En revenant à la construction du
lemme~\ref{17-lemme-existence-resolutions-k-plates-topos-ponctuel}, on peut
observer pour tout $n\in\mathbf{Z}$, $(\rho_{\mathscr A'}F)^n$ est un
faisceau de $\mathscr A'$-Modules libres sur un faisceau d'ensembles
pointés ; cette propriété étant évidemment préservée après application de
$u^\star$, on obtient que $(u^\star \rho_{\mathscr A'}F)^n$ est un
$\mathscr A$-Module plat.

\bigskip

Les lemmes précédents réduisent la
proposition~\ref{17-proposition-u-star-preserve-k-platitude} au lemme suivant :

\begin{lemme3}
Soit $u\colon (\mathscr T,\mathscr A)\to (\mathscr T',\mathscr A')$ un
morphisme de topos annelés en anneaux commutatifs. Si $K\in
C(\mathscr T',\mathscr A')$ est universellement acyclique, alors
$u^\star K\in C(\mathscr T,\mathscr A)$ est universellement acyclique.
\end{lemme3}

On peut démontrer ce lemme en suivant la méthode des limites inductives
locales utilisée dans \sga{4}{v}{8.2.9} ; en présence d'une famille
conservative de foncteurs fibres, il est possible de faire plus simple.

\begin{corollaire3}
Si $u$ et $v$ sont des morphismes composables de topos annelés en anneaux
commutatifs, alors on a un isomorphisme de foncteurs
$\L v^\star\circ \L u^\star\iso \L (u\circ v)^\star$.
\end{corollaire3}

\begin{corollaire3}
Soit $u\colon (\mathscr T,\mathscr A)\to (\mathscr T',\mathscr A')$ un
morphisme de topos annelés en anneaux commutatifs.
Pour tous $K$ et $L$ objets de
$\dtotale(\mathscr T',\mathscr A')$, on a un isomorphisme canonique 
\[\L u^\star K\overset{\L}\otimes_{\mathscr A}\L u^\star
L\simeq \L u^\star (K\overset{\L}\otimes_{\mathscr A'}L)\]
dans $\dtotale(\mathscr T,\mathscr A)$.
\end{corollaire3}

%>>>3
%>>>2
%>>>1
\section{Complexes inversibles}%<<<1

\begin{proposition}\label{17-proposition-objets-inversibles-d-a}
Soit $A$ un anneau commutatif. Soit $X\in \dtotale(A)$. Soit $Y\in
\dtotale(A)$. Soit $X\overset{\L}{\otimes}_AY\simeq A$ un isomorphisme
dans $\dtotale(A)$. Alors, il existe une fonction localement constante
$k\colon \SP(A)\to \mathbf{Z}$, un $A$-module inversible $L$ et des
isomorphismes $X\simeq L[k]$ et $Y\simeq L^\vee [-k]$ (le foncteur de
décalage $[k]$ étant défini de façon évidente).
\end{proposition}

Ce résultat apparaît dans \cite[lemma~3.3, Chapter~V]{RD@Hartshorne}
sous des hypothèses supplémentaires disant que $A$ est noethérien et que
$X$ et $Y$ sont dans $\dmc(A)$. Cette version implique
évidemment une autre où on demande à $A$ d'être noethérien et à $X$ et $Y$
d'être des complexes parfaits. Il est alors facile de supprimer l'hypothèse
noethérienne (cf.~\ega{iv}{8}).
Bref, pour achever la démonstration, il suffit de montrer que
dans les conditions de la proposition ci-dessus, le complexe $X$ (et donc
$Y$ par symétrie des rôles) est un complexe parfait. On dit d'un objet $X$
de $\dtotale(A)$ qu'il est de présentation finie si le foncteur
$\Hom_{\dtotale(A)}(X,-)$ de $\dtotale(A)$ vers la catégorie des groupes
abéliens commute aux sommes directes (infinies). On peut montrer que
$X\in\dtotale(A)$ est un complexe parfait si et seulement s'il est de
présentation finie\;\footnote{Il s'agit d'un bon exercice. Toutefois, on
peut aussi obtenir ce critère en utilisant des principes généraux. L'objet
$A$ de $\dtotale(A)$ est un générateur de présentation finie (\emph{i.e.} le
foncteur cohomologique $\Hom_{\dtotale(A)}(A,-)$ commute aux sommes
directes et est conservatif) ; la sous-catégorie triangulée de $\dtotale(A)$
formée des objets de présentation finie est donc l'enveloppe
pseudo-abélienne de la sous-catégorie triangulée engendrée par $A$,
c'est-à-dire la sous-catégorie des complexes parfaits : combiner
\cite[proposition~8.4.1]{TC@Neeman}, \cite[lemma~4.4.5]{TC@Neeman} et
\cite[remark~4.2.6]{TC@Neeman}.}. De l'isomorphisme
$X\overset{\L}{\otimes}_A Y\simeq A$, on tire un isomorphisme de
foncteurs $\R\SheafHom(X,-)\simeq
Y\overset{\L}{\otimes}_A-\colon \dtotale(A)\to\dtotale(A)$. Comme le
foncteur $Y\overset{\L}{\otimes}_A-$ commute évidemment
aux sommes directes,
c'est aussi le cas du foncteur $\Hom_{\dtotale(A)}(X,-)$, ce qui
montre que $X$ est de présentation finie : il s'agit d'un complexe parfait.

\begin{proposition}\label{17-proposition-objets-inversibles-d-t-a}
Soit $(\mathscr T,\mathscr A)$ un topos annelé en anneaux commutatifs. Soit
$X\in\dtotale(\mathscr T,\mathscr A)$.
Soit $Y\in\dtotale(\mathscr T,\mathscr A)$. Soit
$X\overset{\L}{\otimes}_{\mathscr A}Y\simeq \mathscr A$ un
isomorphisme dans $\dtotale(\mathscr T,\mathscr A)$.
Alors, l'objet final de $\mathscr T$ est recouvert par des objets $U$
tels que $X_{|U}$ et $Y_{|U}$ puissent être 
induits par des complexes parfaits $X'$ et $Y'$ de $\dtotale(\mathscr A(U))$
et que l'on ait un isomorphisme $X'\overset{\L}{\otimes}_{\mathscr
A(U)}Y'\simeq \mathscr A(U)$ compatible à l'isomorphisme donné ; la
forme des complexes $X'$ et $Y'$ est donc connue grâce à la
proposition~\ref{17-proposition-objets-inversibles-d-a}.
\end{proposition}

Démontrons la proposition~\ref{17-proposition-objets-inversibles-d-t-a}. On
peut supposer que $\mathscr T$ est le topos des faisceaux sur un site
$\mathscr C$. Notons $\mathscr T'$ le topos des préfaisceaux sur la
catégorie sous-jacente au site $\mathscr
C$, elle aussi notée $\mathscr C$,
et $\mathscr A'$ le préfaisceau d'anneaux sur $\mathscr C$ défini par
$\mathscr A$. On peut identifier $X\overset{\L}\otimes_{\mathscr A}
Y$ au faisceau associé à $\rho_{\mathscr A'}X\otimes_{\mathscr
A'}\rho_{\mathscr A'}Y$ où $\rho_{\mathscr A'}$ est le foncteur de
résolution $K$-plate défini dans le
paragraphe~\ref{17-subsubsection-produit-tensoriel-derive-prefaisceaux}.
Ainsi, l'isomorphisme donné $\mathscr A\iso
X\overset{\L}\otimes_{\mathscr A}Y$ peut être représenté localement
par un $0$-cocycle $s$ de $(\rho_{\mathscr A'}X\otimes_{\mathscr
A'}\rho_{\mathscr A'}Y)(U)=\rho_{\mathscr A(U)}(X(U))\otimes_{\mathscr
A(U)}\rho_{\mathscr A(U)}(Y(U))$.
Soit $U$ un objet de $\mathscr C$ sur lequel une telle description
est possible. Par construction de $\rho_{\mathscr A(U)}$, il vient qu'il
existe des sous-complexes bornés formés de $\mathscr A(U)$-modules libres de
type fini $X'$ de $\rho_{\mathscr A(U)}(X(U))$ et $Y'$ de $\rho_{\mathscr
A(U)}(X(U))$ tel que $X'\otimes_{\mathscr A(U)} Y'$ contienne $s$. Notons
$\mathscr T_{|U}$ le topos des faisceaux sur $\mathscr T$ au-dessus de $U$
(un site sous-jacent est donné par la catégorie $\mathscr C/U$). On dispose
d'un morphisme de topos annelés évident $\pi$ de $(\mathscr T_{|U},\mathscr
A_{|U})$ vers le topos ponctuel muni de l'anneau $\mathscr A(U)$. Posons
$X''=\pi^\star X'$ et $Y''=\pi^\star Y'$ : ce sont des complexes strictement
parfaits sur $\mathscr T_{|U}$. On dispose de morphismes évidents $X''\to
X_{|U}$ et $Y''\to Y_{|U}$ dans $\dtotale(\mathscr T_{|U},\mathscr A_{|U})$
et d'après ce qui précède, on a un morphisme $\mathscr A_{|U}\to
X''\otimes_{\mathscr A_{|U}} Y''$ factorisant l'isomorphisme $\mathscr
A_{|U}\simeq X_{|U}\overset{\L}\otimes_{\mathscr A_{|U}}Y_{|U}$.
Bref, $\mathscr A_{|U}\simeq X_{|U}
\overset{\L}\otimes_{\mathscr A_{|U}}Y_{|U}$ est un
facteur direct de $X''\otimes_{\mathscr A_{|U}} Y''$. Par ailleurs,
l'isomorphisme $\mathscr A_{|U}\iso X_{|U}
\overset{\L}\otimes_{\mathscr A_{|U}}Y_{|U}$ se factorise aussi par
$X''\otimes_{\mathscr A_{|U}} Y_{|U}$, donc $\mathscr A_{|U}$
est un facteur direct de
$X''\otimes_{\mathscr A_{|U}} Y_{|U}$ ; en tensorisant ce fait avec
$X_{|U}$, on obtient que $X_{|U}$ est un facteur direct de $X''$. Soit
$p\colon X''\to X''$ un projecteur dont l'image soit isomorphe à $X_{|U}$.
Comme $X'$ et $Y'$ sont (strictement) parfaits, quitte à remplacer $U$ par
une famille d'objets le recouvrant, on peut supposer que $p$ provient d'un
projecteur $\tilde p$ sur $X'$ dans $\dtotale(\mathscr A(U))$. L'image de
$\tilde p$ est un complexe parfait $\tilde X$ dans $\dtotale(\mathscr
A(U))$ induisant $X_{|U}$. De même, on peut supposer que $Y_{|U}$ est induit
par un complexe parfait de $\mathscr A(U)$-modules $\tilde Y$. Quitte à
raffiner le recouvrement de $U$, on peut donc supposer que l'isomorphisme
$\mathscr A_{|U}\iso X_{|U}\otimes_{\mathscr A_{|U}}Y_{|U}$ est induit
par un morphisme $\tilde s\colon \mathscr A(U)\to
\tilde X\otimes_{\mathscr A(U)}\tilde Y$. Notons
$C$ un cône de $\tilde s$. Le complexe $C$ est parfait et vérifie
$\pi^\star C\simeq 0$. Il en résulte que quitte à raffiner le recouvrement
de $U$, on peut supposer que $C$ est acyclique, c'est-à-dire que $\tilde s$
est un isomorphisme.

%>>>1
\section{Coefficients universels}%<<<1
\subsection{Énoncés pour \texorpdfstring{$\R\SheafHom$}{RHom}}%<<<2

\begin{proposition2}\label{17-proposition-coefficients-universels-1}
Soit $Z$ un schéma noethérien. Soit $A$ un anneau commutatif noethérien.
Soit $X\in \detalemoinsccoefs Z A$. Soit $Y\in\detalebtfcoefs Z A$. On suppose ou
bien que $A$ est un corps ou bien que
$\R\SheafHom(X,M\overset{\L}{\otimes}_{A} Y)$ est borné
indépendamment du $A$-module $M$. Alors, pour tout $M\in\dplus(A)$, le
morphisme canonique
\[M\overset{\L}{\otimes}_A \R\SheafHom(X,Y)\to
\R\SheafHom(X,M\overset{\L}{\otimes}_A Y)\]
est un isomorphisme. En outre, si on est dans le cas où
$\R\SheafHom(X,M\overset{\L}{\otimes}_{A} Y)$ est borné
indépendamment du $A$-module $M$, alors
$\R\SheafHom(X,Y)\in\detalebtfcoefs Z A$.
\end{proposition2}

On note $u_M\colon M\overset{\L}{\otimes}_A \R\SheafHom(X,Y)\to
\R\SheafHom(X,M\overset{\L}{\otimes}_A Y)$ le morphisme
canonique, pour tout $M\in\detalepluscoefs Z A$. L'hypothèse selon laquelle
$X\in\detalemoinsccoefs Z A$ implique que les foncteurs $\SheafExt^q(X,-)$
commutent aux limites inductives filtrantes de faisceaux de $A$-modules sur
$Z$. En particulier, le foncteur $\R\SheafHom(X,-)$ de la catégorie
$\detalepluscoefs Z A$ dans elle-même commute aux sommes directes
représentables. On en déduit que $u_M$ est un isomorphisme si $M$ est un
$A$-module libre. Par suite, si $M$ est un complexe borné de $A$-modules
libres, alors $u_M$ est un isomorphisme ; si $A$ est un corps, on peut
conclure que $u_M$ est un isomorphisme pour tout $M\in \detalepluscoefs Z A$.

On se place dorénavant dans le cas où où
$\R\SheafHom(X,M\overset{\L}{\otimes}_{A} Y)$ est borné
indépendamment du $A$-module $M$.
Soit $M\in\db(A)$. Montrons que $u_M$ est un isomorphisme.
Pour tout entier relatif $q$, il existe un
morphisme $P\to M$ avec $P$ un complexe borné de $A$-modules libres tel que
si on note $C$ un cône de ce morphisme, alors $C\leq q$ (de telles
inégalités sont à comprendre relativement à la $t$-structure canonique).
On considère le carré commutatif :
\[
\xymatrix{
P\overset{\L}{\otimes}_A \R\SheafHom(X,Y)
\ar[d]\ar[r]^{u_P}_{\sim}&
\R\SheafHom(X,P\overset{\L}{\otimes}_A Y) \ar[d] \\
M\overset{\L}{\otimes}_A \R\SheafHom(X,Y)\ar[r]^{u_M} &
\R\SheafHom(X,M\overset{\L}{\otimes}_A Y)
}\]
Notons $N$ le plus petit entier naturel tel que 
$\R\SheafHom(X,V\overset{\L}{\otimes}_A Y)$ soit $\leq N$ pour
tout $A$-module $V$. Par dévissage, on obtient que les cônes des flèches
verticales du diagramme ci-dessus sont $\leq q+N$. Les deux flèches
verticales induisent donc des isomorphismes sur les objets de cohomologie
${\mathscr H}^i$ pour $i\geq q+n+2$. La flèche du haut étant un
isomorphisme, la flèche du bas induit des isomorphismes sur les 
${\mathscr H}^i$ pour $i\geq q+n+2$. Ce fait étant vérifié pour tout
$q\in\mathbf{Z}$, le morphisme du bas est bien un isomorphisme.

On déduit aussitôt de ce qui précède que
$\R\SheafHom(X,Y)\in\detalebtfcoefs Z A$. On sait par ailleurs que
$Y\in\detalebtfcoefs Z A$ et que $X\in \detalemoinscoefs Z A$. Par conséquent, il
existe un entier relatif $N$ tel que si $M\in\dtotale^{\geq c}(A)$ pour un
certain entier $c$, alors la source et le but de $u_M$ sont $\geq c+N$.
En raisonnant comme ci-dessus, on déduit du fait que $u_M$ soit un
isomorphisme pour tout $M\in\db(A)$ que ce résultat vaut en
fait pour tout $M\in\dplus(A)$.

\begin{proposition2}\label{17-proposition-coefficients-universels-2}
Soit $Z$ un schéma. Soit $A$ un anneau commutatif. Soit
$Y\in\detalebtfcoefs Z
A$. Soit $M$ un complexe pseudo-cohérent dans $\dtotale(A)$. Soit
$N\in\dplus(A)$. Alors, le morphisme canonique 
\[ \R\Hom(M,N)\overset{\L}{\otimes}_A Y\to
\R\SheafHom(M,N\overset{\L}{\otimes}_A Y) \]
est un isomorphisme dans $\detalepluscoefs Z A$.
\end{proposition2}

Pour tout $M\in\dtotale(A)$, notons $v_M$ le morphisme canonique \[
\R\Hom(M,N)\overset{\L}{\otimes}_A Y\to
\R\SheafHom(M,N\overset{\L}{\otimes}_A Y)\;\text{.} \] Bien
entendu, si $M\in\dtotale(A)$ est un complexe parfait, $v_M$ est un
isomorphisme. Pour tout $M\in\dtotale(A)$ pseudo-cohérent
et tout entier relatif $q$, il
existe un morphisme $P\to M$ avec $P$ parfait dont le cône soit $\leq q$.
Pour pouvoir déduire que $v_M$ est un isomorphisme pour tout
$M\in\dtotale(A)$ du cas particulier où $M$ est supposé parfait, il suffit
donc de montrer qu'il existe un entier relatif $c$ tel que pour tout entier
$q$ et tout $M\in\dtotale^{\leq q}(A)$, la source et le but de $v_M$
soient $\geq -q+c$. Notons $a$ un entier tel que $N\geq a$ et $b$ un entier
tel que pour tout $A$-module $V$, $V\overset{\L}{\otimes}_A Y\geq b$.
On vérifie aussitôt que $c=a+b$ convient, ce qui achève la démonstration de
la proposition.

\begin{proposition2}\label{17-proposition-coefficients-universels-3}
Soit $Z$ un schéma noethérien. Soit $A$ un anneau commutatif noethérien.
Soit $X\in \detalemoinsccoefs Z A$. Soit $Y\in\detalebtfcoefs Z A$. On suppose ou
bien que $A$ est un corps ou bien que
$\R\SheafHom(X,N\overset{\L}{\otimes}_{A} Y)$ est borné
indépendamment du $A$-module $N$. Soit $B$ une $A$-algèbre.
Alors, pour tout complexe pseudo-cohérent
$M\in\dtotale(B)$ et pour tout $N\in\dplus(B)$, le morphisme canonique
\[\R\Hom_B(M,N)\overset{\L}{\otimes}_A \R\SheafHom_A(X,Y)
\to \R\SheafHom_B(M\overset{\L}{\otimes}_A X,
N\overset{\L}{\otimes}_A Y)\]
est un isomorphisme dans $\detalecoefs Z B$.
\end{proposition2}

Les hypothèses font que $\R\Hom_B(M,N)$ appartient à $\dplus(B)$.
Si $A$ n'est pas un corps, 
on peut appliquer le résultat de la
proposition~\ref{17-proposition-coefficients-universels-2} en remplaçant
respectivement $A$, $X$, $M$ et $N$ par $B$, $M$, $N$ et
$B\overset{\L}{\otimes}_A \R\SheafHom_A(X,Y)$ (on notera que
$\R\SheafHom_A(X,Y)\in\detalebtfcoefs Z A$ d'après la
proposition~\ref{17-proposition-coefficients-universels-1}). On obtient ainsi
un isomorphisme canonique
\[\R\Hom_B(M,N)\overset{\L}{\otimes}_A
\R\SheafHom_A(X,Y)\iso
\R\SheafHom_B(M,N\overset{\L}{\otimes}_A\R\SheafHom_A(X,Y))\;\text{.}\]
Si $A$ est un corps, on établit cet isomorphisme par un argument similaire
en se ramenant au cas où $M$ est un complexe parfait de $B$-modules.
En appliquant maintenant la
proposition~\ref{17-proposition-coefficients-universels-1}, on obtient un
nouvel isomorphisme :
\[\R\SheafHom_B(M,N\overset{\L}{\otimes}_A\R\SheafHom_A(X,Y))\iso
\R\SheafHom_B(M,\R\SheafHom_A(X,N\overset{\L}{\otimes}_A
Y))\;\text{.}\]
Enfin, on utilise l'isomorphisme cher à Cartan :
\[\R\SheafHom_B(M,\R\SheafHom_A(X,N\overset{\L}{\otimes}_A
Y))\simeq
\R\SheafHom_B(M\overset{\L}{\otimes}_A X,
N\overset{\L}{\otimes}_A Y)\;\text{.}\]
L'isomorphisme voulu est l'isomorphisme obtenu en composant les différents
isomorphismes canoniques ci-dessus.

%>>>2
\subsection{Conséquences pour \texorpdfstring{$\R j_\star$}{Rj*} et \texorpdfstring{$i^!$}{i!}}%<<<2

\begin{proposition2}\label{17-proposition-coefficients-universels-r-j-star}
Soit $j\colon U\to X$ une immersion ouverte entre schémas noethériens. Soit
$A$ un anneau commutatif noethérien. Soit $Y\in\detalebtfcoefs U A$. On suppose
ou bien que $A$ est un corps ou bien que $\R j_\star
(M\overset{\L}\otimes_A Y)$ est borné indépendamment du $A$-module
$M$. Alors, pour tout $M\in\dplus(A)$, le morphisme canonique
\[M\overset{\L}{\otimes}_A\R j_\star Y\to
\R j_\star (M\overset{\L}{\otimes}_A Y)\]
est un isomorphisme dans $\detalecoefs X A$. Si $\R j_\star
(M\overset{\L}\otimes_A Y)$ est borné indépendamment du $A$-module
$M$, alors $\R j_\star Y$ appartient à $\detalebtfcoefs X A$.
\end{proposition2}

Ceci résulte de la
proposition~\ref{17-proposition-coefficients-universels-1}, compte tenu de la
formule $\R j_\star Y\simeq \R\SheafHom(j_!A,j_!Y)$ pour tout
$Y\in\detalecoefs U A$.

\begin{proposition2}
Soit $i\colon Z\to X$ une immersion fermée entre schémas noethériens. Soit
$A$ un anneau commutatif noethérien. Soit $Y\in\detalebtfcoefs Z A$. On suppose
ou bien que $A$ est un corps ou bien que $i^!
(M\overset{\L}\otimes_A Y)$ est borné indépendamment du $A$-module
$M$. Alors, pour tout $M\in\dplus(A)$, le morphisme canonique
\[M\overset{\L}{\otimes}_A i^! Y\to i^!
(M\overset{\L}{\otimes}_A Y)\]
est un isomorphisme dans $\detalecoefs X A$. Si $i^!
(M\overset{\L}\otimes_A Y)$ est borné indépendamment du $A$-module
$M$, alors $i^! Y$ appartient à $\detalebtfcoefs Z A$.
\end{proposition2}

Cette fois-ci, on utilise la formule $i^!Y\simeq \R\SheafHom(i_\star
A,i_\star Y)$.

%>>>2
%>>>1
\section{Modules ind-unipotents}%<<<1
\subsection{Définitions}%<<<2
Soit $G$ un groupe topologique. Soit $A$ un anneau. On appellera ici 
$A[G]$-module un $A[G]$-module (à gauche)
au sens où on l'entend usuellement en
considérant $G$ comme groupe discret. Un $A[G]$-module discret est un
$A[G]$-module dans lequel le stabilisateur de tout élément est ouvert.
Un $A[G]$-module est dit trivial si le groupe $G$ agit
trivialement sur lui ; un tel $A[G]$-module est discret.
On note $I_G$ l'idéal d'augmentation de $A[G]$ ; un
$A[G]$-module est trivial si et seulement s'il est annulé par $I_G$.

La catégorie des $A[G]$-modules discrets est une catégorie abélienne de
Grothendieck ; le foncteur d'inclusion de cette
catégorie dans celle des $A[G]$-modules est exact, commute aux limites
inductives, mais en général pas aux limites projectives. La catégorie des
$A[G]$-modules discrets est stable par sous-quotient (mais en général pas
par extensions) dans celle des $A[G]$-modules.

\begin{definition2}
Un $A[G]$-module discret est unipotent s'il admet une filtration finie dont
les quotients successifs sont des $A[G]$-modules triviaux. L'ordre
d'unipotence d'un $A[G]$-module unipotent est la plus petite longueur d'une
telle filtration ; ainsi, un $A[G]$-module trivial non nul est unipotent
d'ordre $1$. Un $A[G]$-module discret est ind-unipotent si tous ses
sous-$A[G]$-modules de type fini sont unipotents.
\end{definition2}

\begin{definition2}\label{17-definition-filtration-unipotence}
Soit $M$ un $A[G]$-module discret. On définit par récurrence
une filtration croissante
$(\Fil_n M)_{n\in\mathbf{Z}}$ de $M$ par des sous $A[G]$-modules, de façon
à ce que $\Fil_0M=0$ et que pour tout entier $n\in\mathbf{N}$, on ait un
isomorphisme canonique $\Fil_{n+1}M/\Fil_nM\simeq \HH^0(G,M/\Fil_nM)\subset
M/\Fil_n M$.
\end{definition2}

La filtration $\Fil_\bullet M$ est évidemment fonctorielle en $M$ au sens
où si $f\colon M\to M'$ est un morphisme de $A[G]$-modules,
$f(\Fil_nM)\subset \Fil_n(M')$ pour tout $n\in\mathbf{Z}$.

%>>>2
\subsection{Propriétés}%<<<2

\begin{proposition2}\label{17-proposition-proprietes-unipotents}
Les notions d'unipotence et d'ind-unipotence des $A[G]$-modules discrets
jouissent des propriétés suivantes :
\begin{itemize}
\item[(i)] Le caractère unipotent des
$A[G]$-modules discrets est stable par sous-quotients
et extensions ; l'indice d'unipotence décroît par passage à un
sous-quotient et est sous-additif vis-à-vis des extensions ;
\item[(ii)] Le caractère ind-unipotent des $A[G]$-modules discrets est
stable par sous-quotients, et, si $G$ est profini, par extensions ;
\item[(iii)] La catégorie des $A[G]$-modules discrets unipotents
(resp. ind-unipotents) est abélienne, le foncteur d'inclusion
dans celle des $A[G]$-modules discrets est exact ;
\item[(iv)] Un $A[G]$-module discret $M$ est
unipotent si et seulement s'il existe un entier naturel $n$ tel que
$\Fil_nM=M$ ; dans ce cas, l'indice d'unipotence de $M$ est le plus petit
de ces entiers $n$ ;
\item[(v)] Le caractère unipotent (resp. ind-unipotent)
d'un $A[G]$-module discret ne dépend pas de l'anneau
des coefficients $A$ et cette notion n'est pas altérée non plus
si on considère $G$ comme groupe discret ;
\item[(vi)] Pour tout entier naturel $n$ et tout $A[G]$-module
discret, $\Fil_n M$ est l'annulateur de $I_G^n$ dans $M$ ;
\item[(vii)] Pour tout entier naturel $n$,
un $A[G]$-module discret est unipotent d'ordre au plus $n$
si et seulement s'il est annulé par l'idéal $I_G^n$. Un $A[G]$-module
discret $M$ est ind-unipotent si et seulement si tout élément de $M$ est
annulé par une puissance de $I_G$, c'est-à-dire que
$M=\cup_{n\in\mathbf{N}}\Fil_n M$ ;
\item[(viii)] Un $A[G]$-module unipotent est ind-unipotent ;
\item[(ix)] Un $A[G]$-module de type fini est unipotent si et seulement
s'il est ind-unipotent ;
\item[(x)] Une somme directe (resp. une limite inductive)
de $A[G]$-modules ind-unipotents est un $A[G]$-module ind-unipotent.
\item[(xi)] Les $A[G]$-modules (discrets) ind-unipotents sont exactement
les $A[G]$-modules limites inductives filtrantes de $A[G]$-modules
(discrets) unipotents.
\end{itemize}
\end{proposition2}

Montrons (i).
Si $(M,F_\bullet M)$ est un $A[G]$-module discret filtré, tout sous-objet
$M'$ (resp. sous-quotient $M''$) de $M$ est naturellement muni d'une
filtration $F_\bullet M'$ (resp. $F_\bullet M''$) telle que pour tout
$n\in\mathbf{Z}$, le morphisme induit au niveau des gradués $\Gr_n M'\to
\Gr_n M$ (resp. $\Gr_n M\to \Gr_n M''$) soit injectif (resp. surjectif) ;
comme tout sous-quotient d'un $A[G]$-module trivial est trivial, il en
résulte aussitôt que le caractère unipotent d'un $A[G]$-module est stable
par sous-quotient et que l'indice d'unipotence décroît dans cette
opération. Il est évident sur la définition que le caractère unipotent des
$A[G]$-modules discrets est stable par extensions, et que l'indice est
sous-additif vis-à-vis d'elles. (viii) est une conséquence triviale de (i).

Montrons la propriété (iv). Soit $M$ un $A[G]$-module discret.
Si $\Fil_n M=M$, la filtration $(\Fil_iM)_{0\leq i\leq n}$ est une
filtration finie de $M$ de longueur $n$ dont les quotients successifs sont
triviaux ; ainsi, $M$ est unipotent d'ordre au plus $n$. Inversement,
soit $M$ un $A[G]$-module muni d'une filtration $(F_i M)_{i\geq 0}$ telle
que $F_0M=0$ et dont les quotients successifs soient triviaux. Une
récurrence évidente sur $i\in\mathbf{N}$ montre que l'on a une inclusion
$F_i M\subset \Fil_i M$, de sorte que si un entier naturel $n$ est tel que
$F_n M=M$, alors $\Fil_n M=M$.

Dans le cas unipotent, la propriété (v) est une conséquence de la propriété
(iv) : la filtration $\Fil_\bullet M$ est la même que l'on considère $M$
comme $A[G]$-module ou comme $\mathbf{Z}[G]$-module, et que l'on considère
$G$ comme un authentique groupe topologique ou comme un groupe discret.
Quand ce ne sera pas pertinent, on ne précisera donc pas systématiquement
dans la suite que les $A[G]$-modules sont discrets.

Concernant la propriété (vi), il est facile de montrer par récurrence sur
$n\in\mathbf{N}$ que $\Fil_n M$ est l'annulateur de $I_G^n$ dans $M$.

Considérons (vii). On déduit trivialement de (iv) et (vi) qu'un
$A[G]$-module discret est unipotent d'ordre au plus $n$ si et seulement
s'il est annulé par $I_G^n$. Montrons l'autre partie de (vii). Soit $M$ un
$A[G]$-module discret ind-unipotent. Pour tout $m\in M$, le
sous-$A[G]$-module de $M$ engendré par $m$ est unipotent, d'après ce qu'on
vient de montrer, ce module est annulé par une puissance de $I_G$, en
particulier, $m$ est annulé par $I_G^n$ pour un certain entier naturel $n$.
Inversement, supposons que tout élément de $M$ soit annulé par un puissance
de $I_G$. Si on considère un sous-$A[G]$-module de type fini $N$ de $M$, en
appliquant l'hypothèse aux éléments d'un ensemble fini de générateurs de
$N$, on obtient que $N$ est annulé par $I_G^n$ pour un certain entier
naturel $n$, et donc que $N$ est unipotent. Ceci achève la démonstration de
(vii). On sait donc qu'un $A[G]$-module discret $M$ est ind-unipotent si et
seulement si $M=\cup_{n\in\mathbf{N}}\Fil_n M$. Comme la filtration $\Fil_n
M$ ne dépend pas de la topologie de $G$ ni de l'anneau des coefficients
$A$, on peut obtenir (v). La propriété (ix) est une conséquence immédiate
de (vii). La stabilité par sommes directes et sous-quotients du caractère
ind-unipotent en résulte aussi, ce qui établit (x).

Supposons $G$ profini
et montrons la stabilité par extensions des $A[G]$-modules discrets
ind-unipotents, ce qui achèvera la démonstration de (ii).
D'après (v), on peut supposer que $A=\mathbf{Z}$. On se
donne une suite exacte courte $0\to M'\to M\to M''\to 0$ de
$\mathbf{Z}[G]$-modules discrets avec $M'$ et $M''$ ind-unipotents.
Soit $N$ un sous-$\mathbf{Z}[G]$-module
de type fini de $M$. Il s'agit de montrer que $N$ est unipotent. On peut
considérer l'image $N''$ de $N$ dans $M''$, et $N'$ le noyau de la
projection $N\to N''$. Le $\mathbf{Z}[G]$-module $N'$ (resp. $N''$) est un
sous-$\mathbf{Z}[G]$-module de $M'$ (resp. $M''$). Par passage à des
sous-objets, $N'$ et $N''$ sont ind-unipotents. Le groupe $G$ étant profini
et $N$ un $\mathbf{Z}[G]$-module discret, $N$ est un groupe abélien de type
fini ; les groupes abéliens $N'$ et $N''$ qui en sont des sous-quotients
sont eux aussi de type fini. D'après (ix), $N'$ et $N''$ sont
unipotents ; d'après (i), $N$ est unipotent, ce qui achève la
démonstration de (ii).

(iii) résulte aussitôt de (i) et (ii). Enfin, (xi) découle de (vii), (viii)
et (x).

\begin{remarque2}
Si $H\to G$ est un morphisme de groupes topologiques et $M$ un
$A[G]$-module unipotent (resp. ind-unipotent), alors, en tant
que $A[H]$-module, $M$ est unipotent (resp. ind-unipotent).
\end{remarque2}

%>>>2
\subsection{Modules ind-unipotents pour un sous-groupe distingué}%<<<2
\label{17-subsection-modules-ind-unipotents-sous-groupe-distingue}

Soit $A$ un anneau. Soit $G$ un groupe topologique. Soit $H$ un sous-groupe
distingué (fermé) de $G$. Si $M$ est un $A[G]$-module discret, on peut se
demander si en tant que $A[H]$-module, $M$ est unipotent (resp.
ind-unipotent).

Dans cette situation, on note $\Fil_\bullet M$ la filtration de la
définition~\ref{17-definition-filtration-unipotence} pour le groupe
$H$. Le fait que $H$ soit distingué dans $G$ permet d'obtenir aussitôt que
cette filtration $\Fil_\bullet M$ est constituée de sous-$A[G]$-modules de
$M$. De cette remarque et de la 
proposition~\ref{17-proposition-proprietes-unipotents}~(iv),
on tire :

\begin{proposition2}
Soit $M$ un $A[G]$-module. Les conditions suivantes sont équivalentes :
\begin{itemize}
\item en tant que $A[H]$-module, $M$ est unipotent ;
\item il existe une filtration finie de $M$ par des sous-$A[G]$-modules
telle que $H$ agisse trivialement sur les quotients successifs.
\end{itemize}
\end{proposition2}

\begin{proposition2}
Soit $M$ un $A[G]$-module. Les conditions suivantes sont équivalentes :
\begin{itemize}
\item en tant que $A[H]$-module, $M$ est ind-unipotent ;
\item tout sous-$A[G]$-module de type fini de $M$ est unipotent pour
$H$.
\end{itemize}
\end{proposition2}

Il s'agit de montrer que si un sous-$A[H]$-module $N$ de $M$ de type fini
est unipotent, alors le sous-$A[G]$-module (de type fini) de $M$ engendré
par $N$ est unipotent. Il suffit en fait d'établir le lemme suivant :

\begin{lemme2}
Soit $M$ un $A[H]$-module unipotent.
En tant que $A[H]$-module, $A[G]\otimes_{A[H]} M$ est unipotent.
\end{lemme2}

On peut conclure facilement en observant simplement que le $A[H]$-module
$A[G]\otimes_{A[H]} M$ s'identifie à une somme directe de copies de $M$.

\bigskip

Maintenant que les définitions d'unipotence et de ind-unipotences pour le
sous-groupe distingué $H$ sont clarifiées, on peut énoncer la proposition
suivante :

\begin{proposition2}\label{17-proposition-proprietes-unipotents-sous-groupe-distingue}
Les énoncés de la proposition~\ref{17-proposition-proprietes-unipotents}
restent vrais si on
remplace \guil{unipotent} par \guil{unipotent pour $H$},
\guil{ind-unipotent} par \guil{ind-unipotent pour $H$} et $I_G$ par $I_H$,
et que, comme ci-dessus, on définit la filtration $\Fil_\bullet$
relativement au groupe $H$.
\end{proposition2}

Compte tenu des clarifications faites ci-dessus, c'est trivial.

%>>>2
%>>>1

\ifx\danslelivre\undefined
\bibliography{../fichiers-configuration-latex/bibliographie}
\bibliographystyle{../fichiers-configuration-latex/smfalpha}
\end{document}
\fi
